\documentclass[twoside,11pt]{article}

\usepackage{blindtext}
\usepackage{graphicx} % Required to insert images
\usepackage{booktabs} % Allows the use of \toprule, \midrule and \bottomrule in tables
\usepackage{mathrsfs}
\usepackage{amssymb}
\usepackage{amsmath}
\usepackage{mathtools}
\usepackage{bbding}
\usepackage{pifont}
\usepackage{bm}
\usepackage{amsbsy}
\usepackage{algorithm}
\usepackage{algorithmicx}
\usepackage{algpseudocode}
\usepackage{float}
\usepackage{geometry}
\usepackage{lipsum}
\usepackage{cases}
\usepackage{bbm}
\usepackage{color}
\usepackage{tabularx}
\usepackage{CJK}
\usepackage{multirow}
\usepackage{setspace}
\usepackage{geometry}
\usepackage{color}
\usepackage{multicol}
\usepackage{extarrows}
%\usepackage{hyperref}
\usepackage{subfigure}
\usepackage{tikz}
\usepackage{multicol}
\usepackage{comment}



\definecolor{blue}{rgb}{0,0,0.9}
\definecolor{red}{rgb}{0.9,0,0}
\definecolor{green}{rgb}{0,0.9,0}

% Any additional packages needed should be included after jmlr2e.
% Note that jmlr2e.sty includes epsfig, amssymb, natbib and graphicx,
% and defines many common macros, such as 'proof' and 'example'.
%
% It also sets the bibliographystyle to plainnat; for more information on
% natbib citation styles, see the natbib documentation, a copy of which
% is archived at http://www.jmlr.org/format/natbib.pdf

% Available options for package jmlr2e are:
%
%   - abbrvbib : use abbrvnat for the bibliography style
%   - nohyperref : do not load the hyperref package
%   - preprint : remove JMLR specific information from the template,
%         useful for example for posting to preprint servers.
%
% Example of using the package with custom options:
%
% \usepackage[abbrvbib, preprint]{jmlr2e}

\usepackage{jmlr2e}

% Definitions of handy macros can go here
%\input{math command}
\newcommand{\blue}[1]{\begin{color}{blue}#1\end{color}}
\newcommand{\magenta}[1]{\begin{color}{magenta}#1\end{color}}
\newcommand{\red}[1]{\begin{color}{red}#1\end{color}}
\newcommand{\green}[1]{\begin{color}{green}#1\end{color}}
\newcommand{\floor}[1]{\lfloor#1\rfloor}
\newcommand{\ceil}[1]{\lceil#1\rceil}
\newcommand{\be}{\begin{equation}}
\newcommand{\ee}{\end{equation}}
\newcommand{\ben}{\begin{enumerate}}
\newcommand{\een}{\end{enumerate}}
\newcommand{\bfr}{\begin{frame}}
\newcommand{\efr}{\end{frame}}
\newcommand{\bit}{\begin{itemize}}
\newcommand{\eit}{\end{itemize}}


\newcommand{\B}{\mathbb B}
\newcommand{\D}{\mathbb D}
\newcommand{\R}{\mathbb R}
\newcommand{\E}{\mathbb{E}}
\newcommand{\bP}{\mathbb{P}}
\newcommand{\Q}{\mathbb{Q}}
\newcommand{\Z}{\mathbb{Z}}
\newcommand{\W}{\mathbb W}
\newcommand{\cA}{\mathcal A}
\newcommand{\cB}{\mathcal B}
\newcommand{\cC}{\mathcal C}
\newcommand{\cD}{\mathcal D}
\newcommand{\cE}{\mathcal E}
\newcommand{\cF}{\mathcal F}
\newcommand{\cG}{\mathcal G}
\newcommand{\cH}{\mathcal H}
\newcommand{\cI}{\mathcal I}
\newcommand{\cJ}{\mathcal J}
\newcommand{\cL}{\mathcal L}
\newcommand{\cM}{\mathcal M}
\newcommand{\cN}{\mathcal N}
\newcommand{\cO}{\mathcal O}
\newcommand{\cP}{\mathcal P}
\newcommand{\cQ}{\mathcal Q}
\newcommand{\cR}{\mathcal R}
\newcommand{\cS}{\mathcal S}
\newcommand{\cT}{\mathcal T}
\newcommand{\cU}{\mathcal U}
\newcommand{\cW}{\mathcal W}
\newcommand{\cX}{\mathcal X}
\newcommand{\cY}{\mathcal Y}
\newcommand{\cZ}{\mathcal Z}
\newcommand{\bfx}{{\bf x}}
\newcommand{\bfu}{{\bf u}}
\newcommand{\wtnabla}{\widetilde\nabla}

\newcommand{\grd}{\nabla}
\newcommand{\bu}{\bullet}
\newcommand{\bbu}{\blue{\bullet}}
\newcommand{\btr}{\blue{\blacktriangleright}}

\newcommand{\ft}{\footnote}
\newcommand{\ftf}{\footnote[frame]}
\newcommand{\frt}{\frametitle}


% Math Operations
\DeclareMathOperator*{\argmin}{argmin}
\DeclareMathOperator*{\argmax}{argmax}
\def\norm#1{\|#1\|}
\def\inprod#1#2{\langle#1,\,#2\rangle}
\def\kron{\mbox{{\scriptsize$\,\otimes\,$}}}

\def\nn{\nonumber}

\def\sig{\sigma}
\def\lam{\lambda} 
\def\alp{\alpha}
\def\eps{\epsilon}
\def\veps{\varepsilon}
\def\bsig{\bar{\sigma}}
\def\beps{\bar{\epsilon}}
\def\heps{\hat{\epsilon}}
\def\Prox{\mathbb{P}}
\def\Mat{{\rm Mat}}
\def\myfbox#1{\fbox{\begin{minipage}{15cm}#1\end{minipage}}}
\renewcommand{\columnseprulecolor}{\color{red}}

% \newtheorem{definition}{Definition}[section]
% \newtheorem{theorem}{Theorem}[section]
% \newtheorem{corollary}{Corollary}[section]
% \newtheorem{lemma}{Lemma}[section]
% \newtheorem{proposition}{Proposition}[section]
% \newtheorem{remark}{Remark}[section]
 \newtheorem{assumption}{Assumption}%[section]
% \newtheorem{thm}{\bf Theorem}
% \newtheorem{example}{Example}[section]

% Vectors
\def\vzero{{\bm{0}}}
\def\vone{{\bm{1}}}
\def\vmu{{\bm{\mu}}}
\def\vtheta{{\bm{\theta}}}
\def\va{{\bm{a}}}
\def\vb{{\bm{b}}}
\def\vc{{\bm{c}}}
\def\vd{{\bm{d}}}
\def\ve{{\bm{e}}}
\def\vf{{\bm{f}}}
\def\vg{{\bm{g}}}
\def\vh{{\bm{h}}}
\def\vi{{\bm{i}}}
\def\vj{{\bm{j}}}
\def\vk{{\bm{k}}}
\def\vl{{\bm{l}}}
\def\vm{{\bm{m}}}
\def\vn{{\bm{n}}}
\def\vo{{\bm{o}}}
\def\vp{{\bm{p}}}
\def\vq{{\bm{q}}}
\def\vr{{\bm{r}}}
\def\vs{{\bm{s}}}
\def\vt{{\bm{t}}}
\def\vu{{\bm{u}}}
\def\vv{{\bm{v}}}
\def\vw{{\bm{w}}}
\def\vx{{\bm{x}}}
\def\vy{{\bm{y}}}
\def\vz{{\bm{z}}}

% Matrix
\def\mA{{\bm{A}}}
\def\mB{{\bm{B}}}
\def\mC{{\bm{C}}}
\def\mD{{\bm{D}}}
\def\mE{{\bm{E}}}
\def\mF{{\bm{F}}}
\def\mG{{\bm{G}}}
\def\mH{{\bm{H}}}
\def\mI{{\bm{I}}}
\def\mJ{{\bm{J}}}
\def\mK{{\bm{K}}}
\def\mL{{\bm{L}}}
\def\mM{{\bm{M}}}
\def\mN{{\bm{N}}}
\def\mO{{\bm{O}}}
\def\mP{{\bm{P}}}
\def\mQ{{\bm{Q}}}
\def\mR{{\bm{R}}}
\def\mS{{\bm{S}}}
\def\mT{{\bm{T}}}
\def\mU{{\bm{U}}}
\def\mV{{\bm{V}}}
\def\mW{{\bm{W}}}
\def\mX{{\bm{X}}}
\def\mY{{\bm{Y}}}
\def\mZ{{\bm{Z}}}

\newcommand{\dataset}{{\cal D}}
\newcommand{\fracpartial}[2]{\frac{\partial #1}{\partial  #2}}

% Heading arguments are {volume}{year}{pages}{date submitted}{date published}{paper id}{author-full-names}

\usepackage{lastpage}
%\jmlrheading{23}{2022}{1-\pageref{LastPage}}{1/21; Revised 5/22}{9/22}{21-0000}{Author One and Author Two}

% Short headings should be running head and authors last names

\ShortHeadings{Nonconvex Stochastic Bregman Proximal Gradient Method in Deep Learning}{Ding, Li, and Toh}
\firstpageno{1}

\begin{document}

\title{Nonconvex Stochastic Bregman Proximal Gradient Method with Application to Deep Learning}

\author{\name Kuangyu Ding \email  kuangyud@u.nus.edu \\
       \addr Department of Mathematics\\
       University of Singapore\\
       10 Lower Kent Ridge Road, Singapore 119076
       \AND
       \name Jingyang Li \email  li\_jingyang@u.nus.edu \\
       \addr Department of Mathematics\\
       University of Singapore\\
       10 Lower Kent Ridge Road, Singapore 119076
       \AND
       \name Kim-Chuan Toh \email  mattohkc@nus.edu.sg \\
       \addr Department of Mathematics Institute of Operations Research and Analytics\\
       University of Singapore\\
       10 Lower Kent Ridge Road, Singapore 119076}


\editor{}


\maketitle

\begin{abstract}%   <- trailing '%' for backward compatibility of .sty file
{The widely used stochastic gradient methods for minimizing nonconvex composite objective functions require the Lipschitz smoothness of the differentiable part. But the requirement} does not hold true for problem classes including quadratic inverse problems and training neural networks. 
To address this issue, we investigate a family of stochastic Bregman proximal gradient (SBPG) methods, which only require smooth adaptivity of the differentiable part. SBPG replaces the upper quadratic approximation used in SGD with the Bregman proximity measure, resulting in a better approximation model that captures the non-Lipschitz gradients of the nonconvex objective. We formulate the vanilla SBPG and establish its convergence properties under nonconvex setting without finite-sum structure. Experimental results on quadratic inverse problems testify the robustness of SBPG. Moreover, we propose a momentum-based version of SBPG (MSBPG) and prove it has improved convergence properties. We apply MSBPG to the training of deep neural networks with a polynomial kernel function, which ensures the smooth adaptivity of the loss function. Experimental results on representative benchmarks demonstrate the effectiveness and robustness of MSBPG in training neural networks. Since the additional computation cost of MSBPG compared with SGD is negligible in large-scale optimization, {
MSBPG can potentially be employed as an universal open-source optimizer in the future.}

\end{abstract}

\begin{keywords}
 Nonconvex stochastic algorithm, Bregman distance, Smooth adaptivity, Deep neural network, Gradient explosion
\end{keywords}

\section{Introduction}\label{sec-introduction}
In this paper, we present and analyze a family of nonconvex stochastic Bregman proximal gradient methods (SBPG) for solving the following generic stochastic minimization problem:
\begin{equation}
\min_{\vx\in \overline{C}}\;\E_{{\bm\xi}}[f(\vx,{\bm\xi})]+R(\vx),
\label{min-prob}
\end{equation}
where $f(\cdot,{\bm\xi})$ is a nonconvex differentiable function on $\overline{C}$, $R$ is a proper lower-semicontinuous convex function, ${\bm\xi}$ is a random variable, and $\overline{C}$ is the closure of $C$, which is a nonempty convex open subset of $\R^d$. We denote $F(\vx)\coloneqq\E_{{\bm\xi}}[f(\vx,{\bm\xi})]$, and $\Phi(\vx)\coloneqq F(\vx)+R(\vx)$. This stochastic minimization problem, where an optimizer has limited access to the distribution of ${\bm\xi}$ and can only draw samples from it, is prevalent in the fields of machine learning and statistics \citep{hastie2009elements,shapiro2021lectures,zhang2004solving}. In many instances, the smooth part of the objective function $F(\vx)$ can be formulated as a 
finite-sum structure $F(\vx)=\frac{1}{n}\sum_{i=1}^nf_i(\vx)$.
%which is highly efficient from a computational standpoint. 
However, when $n$ is extremely large, calculating the true gradient for the smooth part of the objective function becomes extremely expensive. As a result, stochastic first-order methods, which trace back to the work of \citet{robbins1951stochastic}, have become the prevailing approach for solving these large-scale optimization problems. In particular, stochastic (proximal) gradient descent and its numerous variants \citep{duchi2011adaptive,duchi2009efficient,gu2020unified,kingma2014adam,JMLR:v18:16-410,wang2022riemannian} have been widely utilized in large-scale stochastic optimization for machine learning \citep{lecun2015deep,shapiro2021lectures,zhang2004solving}.
% The basic $k$-th iteration takes the form:
% \begin{equation}
% {\vx^{k+1}}=\argmin\left\{\inprod{\widetilde\nabla_k}{\vx-{\vx^k}}+\frac{1}{2\alpha_k}\|\vx-{\vx^k}\|^2+R(\vx)\right\},
% \label{SPG-scheme}
% \end{equation}
% where $\widetilde\nabla_k$ is an estimator of $\nabla F({\vx^k})$.  
From modeling perspective, stochastic gradient descent can be viewed as minimizing a sequence of upper quadratic approximations of the nonconvex objective $\Phi(\vx)$:
\begin{equation}
{\vx^{k+1}}=\argmin_{\vx\in \overline{C}}\left\{\underbrace{F({\vx^k},{\bm\Xi}_{k})+\inprod{\widetilde\nabla_k}{\vx-{\vx^k}}+\frac{1}{2\alpha_k}\|\vx-{\vx^k}\|^2}_{F_{{\vx^k}}(\vx):\;model\;of\;F\;at\;{\vx^k}}+R(\vx)\right\},
\label{convention-model} 
\end{equation}
where $F(\vx^k,{\bm\Xi}_k)\coloneqq\frac{1}{|{\bm\Xi}_k|}\sum_{{\bm\xi}\in{\bm\Xi}_k}f(\vx^k,{\bm\xi})$, ${\bm\Xi_k}$ is the set of samples of ${\bm\xi}$ at the $k$-th iteration, and $\widetilde{\nabla}_k$ is an estimator of the exact gradient $\nabla F(x^k)$. This modeling perspective is well-known in deterministic optimization, and has been used in methods such as Newton method, Gauss-Newton method, bundle methods, and trust-region methods, as discussed in various sources such as \citet{hiriart1993convex,nesterov2003introductory,lin2007trust,paren2022stochastic}.

Despite being widely used, stochastic gradient methods \eqref{convention-model} still have several well-known bottlenecks  both in theory and practice. One of the crucial requirements for analyzing stochastic gradient methods is the assumption of Lipschitz continuity of the gradient of the differentiable part, which is essential for ensuring the convergence of the algorithm. However, this assumption does not always hold true. For example, even the seemingly simple function $F(x)=x^4$ does not admit a globally Lipschitz gradient over $\R$, which highlights the challenges of analyzing stochastic gradient methods. In addition, choosing the appropriate stepsize is another challenge in the practical usage of stochastic gradient methods. The stepsize has a significant impact on the convergence performance of the algorithm, and finding the optimal stepsize can be a time-consuming  process. Engineers may have to conduct multiple experiments to determine the optimal stepsize, further compounding the complexity of the problem at hand.
% \begin{enumerate}
% \item One of the central requirements for analyzing gradient methods, including the stochastic proximal gradient method \cite{duchi2009efficient}, is the assumption of Lipschitz continuity of the gradient of the smooth part \cite{bertsekas1997nonlinear}. 
% That is, for any $x,y \in \mathbb{R}^d$, there exists a Lipschitz constant $L > 0$ such that
% \[
% \|\nabla F(\vx)-\nabla F(\vy)\|\leq L\|\vx-\vy\|.
% \]
% However, it should be noted that even a simple function like $F(x)=x^4$ does not admit a Lipschitz gradient globally,  which can pose challenges in the analysis of gradient methods.
% \item The convergence behavior of the stochastic gradient methods are highly dependent on the choice of stepsize, which can be challenging to select in practice. Engineers may need to conduct multiple experiments to find the optimal stepsize, which can be time-consuming and computationally expensive.
% \end{enumerate} 

To address these issues, classical approaches often resort to either line search or more complicated inner loops, but these methods can negatively impact the efficiency of the algorithm or even becomes intractable in a stochastic setting. For instance, stochastic proximal point algorithm (PPA) model the approximation of $F(\vx)$ in \eqref{convention-model} as $F_{{\vx^k}}(\vx)=F(\vx,{\bm\xi}_k)+\frac{1}{2\alpha_k}\|\vx-{\vx^k}\|^2$ \citep{bertsekas2011incremental,bianchi2016ergodic,patrascu2017nonasymptotic,rockafellar1976monotone}, which makes the selection of stepsize $\alpha_k$ more robust than the original model \eqref{convention-model}.
However, the application of stochastic PPA is limited due to the difficulty of solving the subproblems, particularly when dealing with  complicated objective functions, such as training deep neural networks. In such cases, solving the subproblem is almost as difficult as solving the original problem, rendering the approach impractical. 
Recently, \citet{bauschke2017descent,lu2018relatively} have proposed using Bregman proximity measures 
%to overcome the issue of gradient Lipschitz continuity and 
to relax the assumption of gradient Lipschitz continuity to smooth adaptivity. The Bregman gradient method was first introduced as the mirror descent scheme by \citet{nemirovskij1983problem} for minimizing convex nonsmooth functions.
%, whose definition will be provided in Section \ref{subsec-Bregman-Proximal-mapping}. 
From the modeling perspective, Bregman methods consider the following subproblem at each iteration:
\begin{equation}
{\vx^{k+1}}=\argmin_{\vx\in \overline{C}}\left\{\underbrace{F({\vx^k},{\bm\Xi}_{k})+\inprod{\widetilde\nabla_k}{\vx-{\vx^k}}+\frac{1}{\alpha_k}\cD_\phi(\vx,{\vx^k})}_{F_{{\vx^k}}(\vx):\;model\;of\;F\;at\;{\vx^k}}+R(\vx)\right\},
\label{model-Bregman}
\end{equation}
where $\cD_\phi$ is the Bregman distance induced by the kernel function $\phi$.
%, whose definitions will be provided in Section \ref{section-notation-preliminaries}.
To convey the advantage of the Bregman proximity model, we present a toy example. Consider objective function $F(x)=x^4$ with non-Lipschitz gradient continuity, we compare the effectiveness of the upper quadratic approximation model \eqref{convention-model} and Bregman proximity model \eqref{model-Bregman}.
% , which use the kernel function $\phi(x)=\frac{1}{2}x^2$ and $\phi(x)=\frac{1}{2}x^2+\frac{1}{4}x^4$ respectively.
% We use two different kernel functions in the models. The first kernel function is $\phi(x)=\frac{1}{2}x^2$, which recovers the conventional gradient method, i.e., the model in Equation \eqref{convention-model}. The second kernel function is $\phi(x)=\frac{1}{2}x^2+\frac{1}{4}x^4$. 
As apparent from Figure \eqref{model-base-figure}(a), the Bregman proximity model \eqref{model-Bregman} ($F_2(x)$) with the kernel function $\phi(x)=\frac{1}{2}x^2+\frac{1}{4}x^4$ can provide a more suitable approximation for $F(x)$ than the upper quadratic approximation model \ref{convention-model} ($F_1(x)$), as the yellow curve {stays} closer to the curve of the objective function $F(x)=x^4$. With this improved approximation, ${x^{k+1}}$ generated by the Bregman gradient method can make more significant progress towards approaching the optimal solution ($x^*=0$) than ${x^{k+1}}$ generated by the gradient descent method, as depicted in Figure \ref{model-base-figure}(b).
\begin{figure}
\centering
\subfigure[Two models' approximations of $F(x)$]{
\includegraphics[width=5.5cm,height=4.5cm]{model_base1-eps-converted-to.pdf}
}
\subfigure[Zoomed-in of version plot (a)]{
\includegraphics[width=5.5cm,height=4.5cm]{model_base2-eps-converted-to.pdf}
}
\caption{For function $F(x)=x^4$, which does not admit a globally Lipschitz continuous gradient. We restrict the feasible set to $[-0.5,2]$. Consider the models \eqref{convention-model} and \eqref{model-Bregman} of $F$ at ${x^k}=1$. The Lipschitz constant of $F$ with respect to the kernel $\phi(x)=\frac{1}{2}x^2$ is 48. The Lipschitz constant of $F$ with respect to the kernel $\phi(x)=\frac{1}{2}x^2+\frac{1}{4}x^4$ is 4. 
%So $L_1=24$, $L_2=4$.
The figure in (b) is a zoomed-in version of the plot in (a) for the range $[0.6,1]$. The unique minimum of $F(x)$ is at $x=0$.}
\label{model-base-figure}
\end{figure}

% addressed the first longstanding issue by replacing the conventional upper quadratic approximation of the smooth part of objective function by a more general Bregman proximity measure, in which the Bregman distance can capture the geometry of the smooth part of objective function. 
% Drawing on this inspiration, the assumption of gradient Lipschitz continuity can be relaxed and replaced by the concept of "smooth adaptivity", the definition is given in Section \ref{subsec-Bregman-Proximal-mapping}. 
While several stochastic extensions of Bregman methods that are based on smooth adaptivity assumption have been developed recently, the current literature primarily focuses on stochastic convex problems \citep{dragomir2021fast,hanzely2021fastest,lu2019relative}. The only existing convergence analysis of Bregman methods for nonconvex problems \citep{latafat2022bregman} requires a finite-sum structure and a novel equivalent consensus reformulation. Moreover, it is memory-intensive and requires essentially periodic computation of the full gradient, which is %ineffective
{expensive}
for large-scale problems such as deep neural networks \citep{defazio2019ineffectiveness}.
As we can see, stochastic Bregman methods have not been fully explored in the context of modern large-scale nonconvex problems such as training neural networks, and {rigorous numerical evaluations} of their performance are limited. Furthermore, current literature lacks attention towards the robustness of stochastic Bregman methods, particularly in terms of selecting stepsizes and initial points, which can significantly impact their performance in large-scale problems.  
% Based on the concept of smooth adaptivity, several stochastic extensions have been developed in recent years \cite{lu2019relative,hanzely2021fastest,dragomir2021fast,latafat2022bregman}. However, \cite{lu2019relative,hanzely2021fastest,dragomir2021fast} focused on stochastic convex problems. The convergence analysis of nonconvex problems in \cite{latafat2022bregman} require a finite sum structure and a novel equivalent consensus reformulation for the finite sum case. The method proposed in \cite{latafat2022bregman} requires a significant amount of memory to store certain quantities of one data pass or essentially is a SVRG type \cite{johnson2013accelerating} algorithm to save memory, which is ineffective in deep learning \cite{defazio2019ineffectiveness}.
% In the context of modern large-scale problems such as deep neural networks, stochastic Bregman methods have not been fully explored, and only have primarily been applied to smaller-scale problems. Therefore, numerical illustrations of stochastic Bregman methods may be limited and may not fully capture the range of challenges that can arise in such applications. Moreover, the current body of literature regarding stochastic Bregman methods lacks attention towards their performance robustness, such as the selection of stepsizes and initial points, which can considerably affect the efficacy of these methods.  

In this paper, we consider stochastic Bregman proximal gradient methods (SBPG) for nonconvex problems (without finite-sum structure requirement) with application to the training of deep neural networks. We establish the convergence result of a vanilla SBPG without Lipschitz smoothness assumption for nonconvex problems. Moreover, we propose a momentum-based SBPG (denoted as MSBPG) and prove that it has improved convergence properties compared with vanilla SBPG. We apply MSBPG to the training of deep neural networks with a polynomial kernel function, which ensures the smooth adaptivity of the loss function. According to our analysis, MSBPG can improve the robustness of training neural networks by mitigating the gradient explosion phenomenon and improve generalization performance
by adopting Bregman proximity approximation to the loss function locally. For numerical illustrations, we conduct numerical experiments on quadratic inverse problems (QIP) and testify vanilla SBPG's robustness to stepsize selection and {initial point's scaling}. We also conduct extensive experiments on training CNNs for image classification and LSTMs for language modeling by employing MSBPG, which is well-suited for solving large-scale problems. Experimental results on representative benchmarks show that our
MSBPG has excellent generalization performance, outperforming the most frequently used optimization algorithms, including SGD~\citep{robbins1951stochastic}, Adam~\citep{kingma2014adam}, and AdamW~\citep{loshchilov2017decoupled}. Furthermore, MSBPG is demonstrated to be robustness to large stepsize and initial point's scaling, which are the common reasons behind gradient explosion. 
% Although several methods have been proposed to enhance the robustness of stepsize selection and relieve gradient explosion in training neural networks, such as adaptive gradient methods \cite{duchi2011adaptive,tieleman2012rmsprop,kingma2014adam}, gradient clipping \cite{pascanu2013difficulty}, and normalized gradient \cite{hazan2015beyond}, existing methods  
% However, as point out in \cite{wilson2017marginal}, adaptive gradient methods tend to give a poor generalization in supervised learning.
% In addition to proper initialization, other techniques have been proposed to address the gradient explosion issue, such as normalized gradient \cite{hazan2015beyond} and gradient clipping \cite{pascanu2013difficulty}. However, it should be noted that these techniques are often heuristic in nature, as exemplified in \cite{goodfellow2016deep}.

% A well-known perspective to view subproblem \eqref{SPG-scheme} is the modeling perspective, i.e., we minimize a sequence of random models of $\Phi$. Equivalently, \eqref{SPG-scheme} can be viewed as minimizing a sequence of random models of $\Phi$, that is
% \begin{equation}
% {\vx^{k+1}}=\argmin_{\vx\in\R^d}\left\{\underbrace{F({\vx^k},{\bm\xi}_{k})+\inprod{\widetilde\nabla_k}{\vx-{\vx^k}}+\frac{1}{2\alpha_k}\|\vx-{\vx^k}\|^2}_{F_{{\vx^k}}(\vx):\;model\;of\;F\;at\;{\vx^k}}+R(\vx)\right\},
% \label{convention-model} 
% \end{equation}
% where ${\bm\xi_k}$ is the set of samples of ${\bm\xi}$ at the $k$-th iteration and $F(\vx,{\bm\xi}_k)=\frac{1}{|{\bm\Xi}_k|}\sum_{{\bm\xi}\in{\bm\Xi}_k}f(\vx,{\bm\xi})$. 
% This modeling perspective is well-known in deterministic optimization, and has been used in methods such as Newton's method, Gaussian Newton's method, bundle methods, and trust-region methods, as discussed in various sources such as \cite{hiriart1993convex,nesterov2003introductory,nocedal1999numerical}. 
% More recently, a growing body of literature has investigated this modeling perspective in weakly convex stochastic optimization, including works such as \cite{davis2020nonsmooth,drusvyatskiy2018error,duchi2018stochastic}. 
% When $F_{{\vx^k}}(\vx)=F(\vx,{\bm\xi}_k)+\frac{1}{2\alpha_k}\|\vx-{\vx^k}\|^2$, stochastic proximal point algorithm (PPA) is recovered \cite{bertsekas2011incremental,bianchi2016ergodic,kulis2010implicit,rockafellar1976monotone}, which have shown that the performance of stochastic PPA is more robust than that of the stochastic gradient method with respect to the selection of $\alpha_k$. However, the application of stochastic PPA is limited due to the difficulty of solving the subproblems, especially when the objective function is extremely complicated, such as in deep neural networks. In such cases, solving the subproblem is almost as difficult as solving the original problem, thus rendering the approach impractical. 
% Motivated by the modeling perspective, we ask the following question: given a non-globally Lipschitz smooth function $F$, can we construct an approximation model that can better capture the non-Lipschitz geometry of objective function, while simultaneously ensuring computational tractability of the subproblem? The answer may be partially positive thanks to the use of a properly selected Bregman distance. Contrary to the conventional stochastic gradient method, we now consider the following model-based subproblem:
% \begin{equation}
% {\vx^{k+1}}=\argmin_{\vx\in\R^d}\left\{\underbrace{F({\vx^k},{\bm\xi}_{k})+\inprod{\widetilde\nabla_k}{\vx-{\vx^k}}+\frac{1}{\alpha_k}\cD_\phi(\vx,{\vx^k})}_{F_{{\vx^k}}(\vx):\;model\;of\;F\;at\;{\vx^k}}+R(\vx)\right\},
% \label{model-Bregman}
% \end{equation}
% where $\cD_\phi$ is the Bregman distance induced by kernel function $\phi$, whose definitions will be present in Section \ref{section-notation-preliminaries}.
% We will begin by examining a toy example to illustrate the key concept. 
% Consider function $F(x)=x^4$, which does not admit a globally Lipschitz continuous gradient. We use two different kernel functions in the models. The first kernel function is $\phi(x)=\frac{1}{2}x^2$, which recovers the conventional gradient method, i.e., the model in Equation \eqref{convention-model}. The second kernel function is $\phi(x)=\frac{1}{2}x^2+\frac{1}{4}x^4$. The approximation of the two types of models is demonstrated in Figure \ref{model-base-figure}. From Figure \ref{model-base-figure}(a), we can see that under kernel $\phi(x)=\frac{1}{2}x^2+\frac{1}{4}x^4$, model \eqref{model-Bregman} can better approximate $F$ than model \ref{convention-model}, as the yellow curve is closer to the curve of the objective function $F(x)=x^4$. Thus, ${x^{k+1}}$  generated by model \eqref{model-Bregman} can make larger progress towards approaching 0 compared to model \eqref{convention-model}, as illustrated in Figure \ref{model-base-figure}(b).
% \begin{figure}
% \centering
% %\quad
% \subfigure[Initial point robustness]{
% \includegraphics[width=5.5cm,height=4.5cm]{model_base1.eps}
% }
% \subfigure[Safe step threshold]{
% \includegraphics[width=5.5cm,height=4.5cm]{model_base2.eps}
% }
% \caption{For function $F(x)=x^4$, which does not admit a globally Lipschitz continuous gradient. We restrict the feasible set to $[-0.5,2]$. Consider two models \eqref{convention-model} and \eqref{model-Bregman} of $F$ at ${x^k}=1$. The Lipschitz constant of $F$ with respect to the kernel $\phi(x)=\frac{1}{2}x^2$ is 24. The Lipschitz constant of $F$ with respect to the kernel $\phi(x)=\frac{1}{2}x^2+\frac{1}{4}x^4$ is 4. So $L_1=24$, $L_2=4$. The figure in (b) is a zoomed-in version of the plot in (a) for the range $[0.6,1]$. The unique minimum of $F(x)$ is at $x=0$.}
% \label{model-base-figure}
% \end{figure}

% Deep neural networks (DNNs) have become a widely-used tool for solving complex tasks in various domains. However, optimizing DNNs is a challenging problem due to the non-convexity of the objective function and the high dimensionality of the parameter space. One of the primary challenges is the lack of global Lipschitz gradient of the loss function, which makes stepsize selection sensitive. Moreover, the gradient explosion can occur during the training process with an improper initial point or large initial stepsize. In recent years, several methods have been proposed to enhance the robustness of stepsize selection and relieve gradient explosion in DNNs. The most popular ones that enhancing the stepsize selection are adaptive gradient methods, such as Adagrad \cite{duchi2011adaptive}, RMSPROP \cite{tieleman2012rmsprop}, Adam \cite{kingma2014adam}. However, as point out in \cite{wilson2017marginal}, adaptive gradient methods tend to give a poor generalization in supervised learning. Weight initialization strategies, such as the Xavier initialization \cite{glorot2010understanding} and the He initialization \cite{he2015delving}, can reduce the likelihood of the gradient explosion problem. In addition to proper initialization, other techniques have been proposed to address the gradient explosion issue, such as normalized gradient \cite{hazan2015beyond} and gradient clipping \cite{pascanu2013difficulty}. However, it should be noted that these techniques are often heuristic in nature, as exemplified in \cite{goodfellow2016deep}.

 To summarize, our contributions are as follows:
\begin{itemize}
    \item[1.] We investigate Stochastic Bregman Proximal Gradient (SBPG) method to solve nonconvex problem without finite-sum requirement, which employs 
    Bregman distance to handle the non-Lipschitz gradient continuity. We establish convergence results for the vanilla SBPG in the {sense of expectation}. Further, we propose a momentum-based SBPG {(MSBPG)} that is tailored for modern large-scale applications, and prove that it has improved convergence properties compared to the vanilla SBPG. To our best knowledge, this is the first time that the momentum technique has been integrated into {a stochastic Bregman proximal gradient method.} 
    % Our analysis demonstrates that the momentum-based SBPG achieves superior convergence results compared to its non-momentum counterpart under the expectation sense.
    \item[2.] We apply {MSBPG} to training deep neural networks (DNN), which leverages on a suitable polynomial kernel function to ensure that the DNN's loss function is smooth adaptable with respect to the designed kernel function. MSBPG exhibits good convergence behavior and excellent generalization performance on extensive tasks. Moreover, MSBPG is found to be more robust than the traditional SGD, especially when it comes to stepsize selection and initialization. We highlight that MSBPG is a theoretically derived method that is able to ease the difficulty of selecting stepsize, mitigate gradient explosion, and maintains excellent generalization performance simultaneously. This distinguishes MSBPG from many existing techniques that rely on intuition and empirical observations. %Notably, SBPG is the first theoretically derived method that is easy to implement and robust in terms of both stepsize selection and gradient explosion mitigation. This distinguishes SBPG from many existing techniques that rely on intuition and empirical evidence.  
    \item[3.] We demonstrate the efficiency and robustness of SBPG or MSBPG in a range of applications, including sparse quadratic inverse problems and large-scale deep neural networks. In quadratic inverse problem, SBPG is more robust in terms of both stepsize and initial point selections. In training deep neural networks, MSBPG has been been found to achieve a superior generalization performance compared with some of the most frequently used optimizers {such as Adam and AdamW}, and also exhibits robustness to stepsize selection and initialization. These results highlight the potential of MSBPG as a powerful tool for optimizing complex and large-scale deep neural networks, thus offering a promising direction for future research in this area.  
    %Specifically, we demonstrate superior properties in the robustness of stepsize selection compared to conventional SGD or SPG \cite{duchi2009efficient}. We observe that SBPG is relatively robust to the selection of the initial point, which is not the case for conventional SGD. Specifically, when the norm of the initial point is large, conventional SGD may suffer from gradient explosion phenomena.
     %Specifically, SBPG has been shown to outperform state-of-the-art optimizers in terms of convergence speed and generalization performance in deep neural network applications. Additionally, our method has superior performance in terms of robustness of stepsize and initialization, which helps relieve gradient explosion during training. These results highlight the potential of SBPG as a powerful tool for optimizing complex and large-scale deep neural networks, offering a promising direction for future research in this area.
\end{itemize}
The remainder of this paper is organized as follows. In Section \ref{section-notation-preliminaries}, we present notation, some related preliminaries and our problem setting. In Section \ref{section-vanilla-SBPG}, we first describe SBPG and establish its convergence results in the sense of expectation. Then, we propose a momentum-based SBPG (MSBPG) and prove its improved convergence properties. In Section \ref{DNN implementation}, we adapt MSBPG to the training of deep neural networks and analyze its capacity in 
{mitigating} gradient explosion and improving generalization capacity. In Section \ref{section-experiments}, we present numerical experiments that demonstrate the efficiency and robustness of vanilla SBPG on quadratic inverse problems and {MSBPG on} 
training deep neural networks. Finally, we give some concluding remarks in Section \ref{section-conclusion}, summarizaing our key contributions and outlining promising {topics} for future research.

 
\section{Preliminaries and Problem setting}\label{section-notation-preliminaries}
In this paper, vectors are represented using boldface letters like $\vv$, while scalars are represented using normal font. Given a proper, lower-semicontinuous function $F:\R^d\rightarrow\bar\R := [-\infty,\infty]$, $dom\,F=\{\vx:F(\vx)<\infty\}$. The Fenchel conjugate function of $F$ is defined as $F^*(\vy)=\sup\{\inprod{\vx}{\vy}-F(\vx):\vx\in\R^d\}$. Given a set $\cS\subset\R^d$, $\bar\cS$ denotes its closure, $int\,\cS$ denotes the set of interior points. A function is of class $\cC^k(\cS)$ if it is $k$ times differentiable  and the $k$-th derivative is continuous on $\cS$. We say that $F$ is level bounded if the 
set $\{x:F(\vx)<\alpha\}$ is bounded for any real number $\alpha$. Given a matrix $A$, ${\rm Vec}(A)$ denotes the vectorization of $A$ by column order. ${\rm Mat}(\cdot)$ is the inverse operation of ${\rm Vec}(\cdot)$, which reshapes a vector back into its original matrix form. Define the operator $\text{Diag}(\cdot)$ to {map}  a vector into a diagonal matrix with diagonal elements equal to the corresponding entries of the vector. The Hadamard product is represented by the symbol $\circ$. If we use the notation $\|\cdot\|$ without any additional explanation, we assume that it refers to the Euclidean vector norm for vectors and the Frobenius matrix norm for matrices. 

Let $(\Omega,\cF,\mathbb P)$ be a  probability space. Given a random variable ${\bm\xi}$ and a $\sigma$-algebra $\cF$, we write ${\bm\xi}\vartriangleleft\cF$ if ${\bm\xi}$ is measurable over $\cF$. {Let $\{{\bm\xi}_k\}_{k\geq0}$ be a stochastic process, and $\{\cF_k\}_{k\geq0}$ be a filtration, where $\cF_k$ defined by a $\sigma$-algebra $\cF_k:=\sigma(\bm\xi_0, \ldots, \bm\xi_{k-1}) $ on $\Omega$. The conditional expectation is denoted by $\E[\cdot|\cF_k]$.  For simplicity, we use the notation  $\E[\cdot]$ to denote $\E[\cdot|\cF_{\infty}]$. The sequence $\{\vx^k\}_{k\geq0}$ generated by our proposed method is adapted to the filtration $\{\cF_k\}_{k\geq0}$, i.e. $\vx^k\vartriangleleft\cF_k$, for all $k\geq0$. The notation $\wtnabla_k$ represents an estimator of the exact gradient $\nabla F(\vx^k)$, which satisties $\widetilde\nabla_k\vartriangleleft\cF_{k+1}$. This estimator is applicable to both the vanilla and momentum cases. The stochastic error is denoted by ${\bm\varepsilon}_k=\nabla F(\vx^k)-\wtnabla_k$. The unbiasedness of the stochastic error {$\bm{\varepsilon}_k$} is assumed throughout this paper, i.e., $\E[{\bm\varepsilon}_k|\cF_{k}]=0$.}
The following supermartingale convergence theorem is a fundamental tool in the analysis of stochastic algorithms.
\begin{theorem} \citep{robbins1951stochastic}
Let $\left\{y_k\right\},\left\{u_k\right\},\left\{a_k\right\}$ and $\left\{b_k\right\}$ be non-negative adapted processes with respect to the filtration $\left\{\mathcal{F}_k\right\}$ such that $\sum_k a_k<\infty, \sum_k b_k<\infty$, and for all $k$, $\mathbb{E}\left[y_{k+1} \mid \mathcal{F}_k\right] \leq\left(1+a_k\right) y_k-u_k+b_k$ almost surely. Then, $\left\{y_k\right\}$ converges almost surely to a non-negative finite random variable and $\sum_k u_k<\infty$ almost surely.
\label{martingale-convergence}
\end{theorem}

\subsection{Smooth adaptable functions}
In this subsection, we introduce the concept of smooth adaptivity, which was initially proposed in \cite{bolte2018first}. This concept is an extension of the idea of relative smoothness for convex functions {introduced in} \cite{bauschke2017descent,lu2018relatively}.  We first give the definitions of kernel function and Bregman distance.
\begin{definition}\label{Bregman-distance-def}
(Kernel function and Bregman distance). Let $\cS$ be a nonempty, convex and open subset of $\R^d$. Associated with $\cS$, a function $\phi:\R^d\rightarrow(-\infty,+\infty]$ is called a kernel function if it satisfies the following two conditions:
\begin{enumerate}
    \item $\phi$ is proper, lower-semicontinuous and convex, with $dom\,\phi\subset\bar\cS$, $dom\,\partial\phi=\cS$.
    \item $\phi\in\cC^1(\cS)$ and $int\,dom\,\phi=\cS$.
\end{enumerate}
Denote the class of kernel function {associated with $\cS$}  by $\cM(\cS)$. Given $\phi\in\cM(\cS)$, the Bregman distance \citep{bregman1967relaxation} generated by $\phi$ is defined as $\cD_\phi(\vx,\vy):dom\,\phi\times int\,dom\,\phi\rightarrow[0,+\infty)$, where
\[
\cD_\phi(\vx,\vy)=\phi(\vx)-\phi(\vy)-\inprod{\nabla\phi(\vy)}{\vx-\vy}.
\]
\end{definition}
Bregman distance measures the difference between the value of $\phi$ at $\vx$ and its linear approximation at $\vy$ based on the gradient of $\phi$ at $\vy$. Some basic properties of Bregman distance can be found in \cite{chen1993convergence,teboulle2018simplified}. Some  kernel functions commonly used in optimization are $\frac{1}{2}\|\vx\|^2$, $\frac{1}{2}\|\vx\|^2+\frac{\alpha}{4}\|\vx\|^4$, $-\sum_{i=1}^d\log\vx_i$ and $\sum_{i=1}^d\vx_i\log\vx_i$, where $\frac{1}{2}\|\vx\|^2\in\cM(\R^d)$ recovers the classical half squared Euclidean distance. The kernel function $\frac{1}{2}\|\vx\|^2+\frac{\alpha}{4}\|\vx\|^4\in\cM(\R^d)$ has found applications in various problems, such as quadratic inverse problems, non-negative matrix factorization, and low-rank minimization \citep{bolte2018first,dragomir2021quartic}. The entropy function $\sum_{i=1}^d\vx_i \log \vx_i\in\cM(\R^d_{++})$ is commonly used in applications that involve probability constraints, where the resulting Bregman distance is known as the Kullback–Leibler (KL) divergence. Throughout the paper we will focus on the following pair of functions $(f,\phi)$ satisfying smooth adaptivity condition. We introduce this concept in the following definition: 
%%
\begin{definition}\label{relative-smooth-def}
 (Smooth adaptivity). Given a kernel function $\phi\in\cM(\cS)$, 
 a proper lower-semicontinuous function $f:\R^d\rightarrow(-\infty,+\infty]$ with $dom\,f\supset dom\,\phi $ that  is $\cC^1$ on $\cS$. $f$ is $L$-smooth adaptable with respect to $\phi$ if there exists $L>0$, such that $L\phi+f$ and $L\phi-f$ are convex on $\cS$.
\end{definition}
Alternative definition of smooth adaptivity is the two-side descent lemma \cite[Lemma 2.1]{bolte2018first}. When both $f$ and $\phi$ belong to $\cC^2(\cS)$, we can verify their smooth adaptivity by comparing the Hessians of $f$ and $\phi$.
\begin{lemma}\label{two-side-descent}
$f$ is $L$-smooth adaptable with respect to $\phi\in\cM(\cS)$, if and only if 
\[
|f(\vx)-f(\vy)-\inprod{\nabla f(\vy)}{\vx-\vy}|\leq L\cD_\phi(\vx,\vy),\;\forall\,\vx,\vy\in int\,dom\,\phi.
\]
Moreover, when both $f$ and $\phi$ {belong to} $\cC^2(int\,dom\,\phi)$, then the above is equivalent to
\[
\exists L>0,\;L \nabla^2\phi(\vx)-\nabla^2f(\vx) \succeq 0, \text { for all } \vx \in int\,dom\phi.
\]
\end{lemma}
The following four-point identity is frequently employed in our proofs, and can be easily verified.
\begin{lemma}(Four points identity) 
Given points $\va,\vb,\vc,\vd$ and any convex function $\phi$ which is differentiable at $\va$ and $\vb$, then
\[
\inprod{\nabla\phi(\va)-\nabla\phi(\vb)}{\vc-\vd}=\cD_\phi(\vc,\vb)+\cD_\phi(\vd,\va)-\cD_\phi(\vc,\va)-\cD_\phi(\vd,\vb).
\]
\label{four-point-identity}
\end{lemma}


\subsection{Bregman Proximal Mapping}\label{subsec-Bregman-Proximal-mapping}
Throughout this paper, we make the following basic assumptions.
\begin{assumption}\label{relative-smooth-assumption}
(Basic requirements). In problem \eqref{min-prob}:
\begin{itemize}
 \item[A1.] For {every} fixed $\bm\xi$, $f(\cdot,\bm\xi)$ is a proper lower-semicontinuous function with $dom\,\phi\subset dom\,f(\cdot,\bm\xi)$, {and it} is $\cC^1$ on $int\,C$.
 \item[A2.] The Legendre kernel (Definition \ref{Legendre}) $\phi\in\cM(C)$ is $\mu$-strongly convex {for some} $\mu>0$.  For every fixed ${\bm\xi}$, $f(\cdot,\bm\xi)$ is $L_F$-smooth adaptable with respect to $\phi$, {where $L_F$ is independent of $\bm\xi$.} 
%\[
%\begin{aligned}
%|f(\vx,{\bm\xi})-f(\vy,{\bm\xi})-\nabla f(\vy,{\bm\xi})|\leq L_F\cD_\phi(\vx,\vy).
%\end{aligned}
%\] 
 \item[A3.] $R$ is is a proper, lower semicontinuous and convex function with $dom\,R\cap int\,C\neq\emptyset$.
 \item[A4.] $\inf_{\vx\in\overline{C}}\{\Phi(\vx)\}>-\infty$.  %\item[A5.] For any $\alpha>0$, $\alpha R+\phi$ is supercoercive, i.e.,
% \[
% \lim_{\|\vu\|\rightarrow\infty}\frac{\alpha R(\vu)+\phi(\vu)}{\|\vu\|}=\infty.
% \]
%For a given $\alpha\in(0,\frac{1}{L_F})$, it holds for any $\vg\in C$, $\vy\in int\,C$ that 
% \begin{equation}\label{Breg-prox-assumption}
% \arg\min_{\vx\in C}\;\left\{\inprod{\vg}{\vx-\vy}+\frac{1}{\alpha_k}\cD_\phi(\vx,\vy)+R(\vx)\right\}\subset int\,C.
% \end{equation}
\end{itemize}
\end{assumption}
 % Assumption \ref{relative-smooth-assumption} is a standard assumption for Bregman type methods, see for example \cite{bolte2018first,latafat2022bregman}, which usually holds in practice and makes Bregman type methods well-defined. 
 % %As we will see later, subproblem \eqref{Breg-prox-assumption} is the definition of Bregman proximal mapping. In our proposed method, the main cost are from the computation of the Bregman proximal mapping. 
 % We first recall the definitions of the Legendre function in \cite{latafat2022bregman}, which makes additional supercoercive condition on the concept in \cite{rockafellar1997convex}.
 Assumption \ref{relative-smooth-assumption} is a standard requirement for Bregman-type methods and is usually satisfied in practice. 
 %When $\overline{C}$ is compact, A5 is automatically satisfied. 
 It ensures the well-definedness of Bregman-type methods, as shown in \cite{bolte2018first,latafat2022bregman}. We also recall the definition of the Legendre function in \cite{latafat2022bregman}, which makes additional supercoercive conditions on the concept in \cite{rockafellar1997convex}. %The Bregman proximal mapping is defined in terms of subproblem \eqref{Breg-prox-assumption}, which is the main computational cost in our proposed method.

\begin{definition}\label{Legendre}
(Legendre kernel). Let $\phi:{\overline{C}}\rightarrow(-\infty,\infty]$ be a proper
lower-semicontinuous convex function. 
It is called essentially smooth if $\text{int\;dom}\,\phi$ is nonempty and $\phi$ is differentiable on $\text{int\;dom}\,\phi$, moreover $\lim_{k\rightarrow\infty}\|\nabla\phi({\vx^k})\|=\infty$ whenever $\{{\vx^k}\}_{k\in\mathbb N}$ converges to a boundary point of $\text{dom}\,\phi$. The function $\phi$ is called Legendre function if it is essentially smooth, {strictly} convex on $\text{int\,dom}\,\phi$ and supercoercive, i.e. $\lim_{\|\vx\|\rightarrow\infty}\frac{\phi(\vx)}{\|\vx\|}=\infty$.
\end{definition}

\begin{definition}\label{Bregman-prox-maping}
Given a nonempty convex open set $C$, a proper lower-semicontinuous convex function $R$ and a Legendre kernel function $\phi\in\cM(C)$, $\vx\in int\,\phi$, we denote the Bregman proximal mapping by ${\rm Prox}^\phi_R:=(\nabla\phi+\partial R)^{-1}\nabla\phi$, which is defined as
\begin{equation}
{\rm Prox}^\phi_R(\vx):=\argmin_{\vu\in\overline{C}}\;\{R(\vu)+\cD_\phi(\vu,\vx)\}.
\label{Breg-prox-def}
\end{equation}  
\end{definition}

%\begin{example}
%Some commonly used kernel functions are $\frac{1}{2}\|\vx\|^2$, $\frac{1}{2}\|\vx\|^2+\frac{\alpha}{4}\|\vx\|^4$, $-\sum_{i=1}^d\log(x_i)$ and\\ $\sum_{i=1}^dx_i\log x_i$, where $\frac{1}{2}\|\vx\|^2$ recovers the half squared Euclidean distance. The kernel function $\frac{1}{2}\|\vx\|^2+\frac{\alpha}{4}\|\vx\|^4$ has found applications in various problems, such as quadratic inverse problems, non-negative matrix factorization, and low-rank minimization \cite{bolte2018first,dragomir2021quartic}. The entropy function $\sum_{i=1}^d x_i \log x_i$ is commonly used in applications that involve probability constraints, where the resulting Bregman distance is known as the Kullback–Leibler (KL) divergence.
%\end{example}

Note that the objective function of \eqref{Breg-prox-def} is strictly convex on $\text{dom}\,\phi\cap domR$, therefore \eqref{Breg-prox-def} has at most one solution. To ensure that \eqref{Breg-prox-def} is well-defined, the following result claims that ${\rm Prox}^\phi_{\alpha R}(\vx)$ is well-defined for any $\alpha>0$, and moreover ${\rm Prox}^\phi_{\alpha R}(\vx)\in\text{int\;dom}\,\phi$ under standard assumptions. The proof can be found in Appendix \ref{appendix-preliminary}.
\begin{lemma}
Suppose Assumption \ref{relative-smooth-assumption} holds. Then \eqref{Breg-prox-def} has a unique solution. Moreover, the solution ${\rm Prox}^\phi_{\alpha R}(\vx)\in C$. 
\label{well-def-Breg-prox}
\end{lemma}
The following proposition for Bregman proximal mapping generalizes the nonexpansive property of the classical proximal mapping (in the case $\phi(\vx)=\frac{1}{2}\|\vx\|^2$). This property is commonly used in convergence proofs. The proof of the following proposition can be found in Appendix \ref{appendix-preliminary}.
\begin{proposition}
 Suppose Assumption \ref{relative-smooth-assumption} holds. Let $\vx_i^+:={\rm Prox}^\phi_R(\nabla\phi^*(\vx_i))$, $i=1,2$. Then $\|\vx^+_1-\vx^+_2\|\leq\frac{1}{\mu}\|\vx_1-\vx_2\|$.
\label{nonexpansive}
\end{proposition}


In this paper, we make the assumption that $R$ and $\phi$ are simple enough so that \eqref{Breg-prox-def}  either has a closed-form solution or {admits} an efficient subroutine to solve it. Using the definition of the Bregman proximal mapping, we can then define the Bregman gradient mapping associated with \eqref{min-prob}. This mapping measures the 
{solution accuracy of}
the methods we propose. Note that $\phi$ is a Legendre kernel, which implies that $\phi^*\in\cC^{1}(\R^d)$ is strictly convex and $(\nabla\phi)^{-1}=\nabla\phi^*$ \citep[Corollary 13.3.1, Theorem 26.5]{rockafellar1997convex}. Therefore, the following concept is well-defined. 

\begin{definition}[Bregman Gradient Mapping]
Given $\alpha>0$, a nonempty convex open set $C$ and a Legendre kernel function $\phi\in\cM(C)$, the Bregman gradient mapping associated with \eqref{min-prob} is defined as follows
\[
\mathcal G_\alpha(\vx)=\frac{\vx-{\rm Prox}^\phi_{\alpha R}\left(\nabla\phi^*(\nabla \phi(\vx)-\alpha\nabla F(\vx))\right)}{\alpha}.
\]
To simplify notation, we use $\cG(\vx)$ to denote  $\cG_1(\vx)$ when $\alpha=1$.
\label{Breg-prox-gradient-mapping}
\end{definition}
When the kernel function $\phi(\vx) = \frac{1}{2}\|\vx\|^2$, the resulting Bregman Gradient Mapping becomes equivalent to the classical Gradient Mapping \citep{nesterov2003introductory,nesterov2005smooth}, which measures the solution's accuracy for proximal gradient methods.
\begin{definition}
(Limiting subdifferential \citep[Definition 8.3]{RockWets98}) Consider a function $f:\R^d\rightarrow\bar\R$ and a point $x$, the regular subdifferential is defines as
\[
\hat\partial f(\vx)=\{\vv:f(\vy)\geq f(\vx)+\inprod{\vv}{\vy-\vx}+o(\|\vy-\vx\|)\}.
\]
The limiting subdifferential is defined as 
\[
\partial f(\vx)=\{\vv:\vx_n\rightarrow \vx,f(\vx_n)\rightarrow f(\vx),\vv_n\in\hat\partial f(\vx_n),\;and\;\vv_n\rightarrow \vv\}.
\]
\end{definition}
By Fermat's rule \cite[Theorem 10.1]{RockWets98}, the set of critical point of $\Phi$ is given by
\[
{\rm crit}\,\Phi=\left\{x\in\R^d:\;0\in\partial\Phi(x)\equiv\nabla F(x)+\partial R(x)\right\}.
\]
The Bregman Gradient Mapping can also be used to evaluate the 
{solution accuracy} for Bregman methods. Let $\vx^+ = \text{Prox}^\phi_{\alpha R}(\nabla\phi^*(\nabla \phi(\vx) - \alpha\nabla F(\vx)))$. From Definition \ref{Breg-prox-gradient-mapping} and equation \eqref{Breg-prox-def}, it can be easily verified by definition that $0 \in \partial\Phi(\vx) \Leftrightarrow 0 = \mathcal{G}_{\alpha}(\vx).$ 
{Hence,} $0 \in \partial\Phi(\vx^+)$ for any $\alpha > 0$. 
The proof of this result is omitted for brevity. Furthermore, if $\nabla\phi$ is $L_\phi$-Lipschitz continuous, then the following proposition holds, implying that $\|\cG_\alpha(\vx)\|$ can be used as a reasonable criterion {to measure the accuracy of $\vx$.}
\begin{proposition}
\label{prop-Bregman-mapping-subdifferential}
Suppose Assumption \ref{relative-smooth-assumption} holds and that $\nabla\phi$ is $L_\phi$ Lipschitz continuous. Then, we have the following inequality:
\[
{\rm dist}\left(0,\partial\Phi({\vx^+})\right)\leq(1+\alpha L_F)L_\phi\|\cG_\alpha(\vx)\|.
\]
\end{proposition}
 We also define the stochastic counterpart of Definition \ref{Breg-prox-gradient-mapping}, which is commonly utilized to evaluate the accuracy of solutions for nonconvex stochastic proximal gradient methods, as discussed in \cite{ghadimi2016mini}.
\begin{definition}\label{stoc-Bregman-grad-map}
(Stochastic Bregman Gradient Mapping).
Given $\alpha>0$ a nonempty convex open set $C$ and a Legendre kernel function $\phi\in\cM(C)$, the stochastic Bregman gradient mapping associated with \eqref{min-prob} is defined as follows
 \[
 \widetilde{\mathcal G}_{\alpha}({\vx}):=\frac{{\vx}-{\rm Prox}^\phi_{\alpha R}\left(\nabla\phi^*\left(\nabla\phi({\vx})-\alpha\widetilde\nabla\right)\right)}{\alpha},\;\text{where }\widetilde{\nabla}\;\text{is an estimator of }\nabla F(x).
 \]   
\end{definition}

\section{Stochastic Bregman Proximal Gradient Method}\label{section-vanilla-SBPG}
In this section, we will study the Stochastic Bregman Proximal Gradient method (SBPG) with the following update scheme:
\begin{equation}
{\vx^{k+1}}=\argmin_{\vx\in\overline{C}}\;R(\vx)+\inprod{\widetilde\nabla_k}{\vx-{\vx^k}}+\frac{1}{\alpha_k}\cD_\phi(\vx,{\vx^k}).
\label{vanilla-SBPG}
\end{equation}
We call {the above method as} "vanilla" SBPG in this section, meaning that the method we study is a basic version without any additional techniques such as variance reduction, momentum, etc., except for the use of mini-batches. In this case, we suppose the following assumptions.
\begin{assumption}\label{variance-assumption}
(Noise requirement).
The estimator satisfies the following two conditions:
$$\E[\widetilde{\nabla}_k|\cF_k]=\nabla F(x_k) \quad \text{and}\quad 
\mathbb E[\|\widetilde\nabla_k-\nabla F({\vx^k})\|^2|\cF_k]\leq\frac{\sigma^2}{m_k},$$ where $m_k$ is the size of the mini-batch in the $k$-th iteration.    
\end{assumption}
 {Note} that we do not assume a finite-sum structure for $F(\vx)$ in this section. The solution of \eqref{vanilla-SBPG} can be written in the form of the Bregman proximal mapping. This is stated in the following proposition. 
 \begin{proposition} 
  \label{subprob-closed-form-sol} 
 Suppose Assumption \ref{relative-smooth-assumption} holds. Then the solution of \eqref{vanilla-SBPG} can be written as the following Bregman proximal mapping:
 \[
 {\vx^{k+1}}={\rm Prox}^\phi_{\alpha_k R}\left(\nabla\phi^*\left(\nabla\phi({\vx^k})-\alpha_k\widetilde\nabla_k\right)\right).
 \]
 \end{proposition}
 \begin{proof}
From the optimality condition of the main subproblem \eqref{vanilla-SBPG}, we have
 \[
 0\in\partial R({\vx^{k+1}})+\wtnabla_k+\frac{1}{\alpha_k}\left(\nabla\phi({\vx^{k+1}})-\nabla\phi({\vx^k})\right).
 \]
Let $\vu^{k+1}:={\rm Prox}^\phi_{\alpha_k R}\left(\nabla\phi^*\left(\nabla\phi({\vx^k})-\alpha_k\widetilde\nabla_k\right)\right)$. From the definition of Bregman proximal mapping, we have
\[
\vu^{k+1}=\arg\min_\vu\;\left\{\alpha_kR(\vu)+{\cD_\phi}\left(\vu,\nabla\phi^*\left(\nabla\phi({\vx^k})-\alpha_k\widetilde\nabla_k\right)\right)\right\},
\]
which is equivalent to
\[
0\in\alpha_k\partial R(\vu^{k+1})+\nabla\phi(\vu^{k+1})-\nabla\phi\big(\nabla\phi^*(\nabla\phi({\vx^k})-\alpha_k\wtnabla_k)\big).
\]
Note that the function $\phi^*$ is the Fenchel conjugate of the Legendre kernel $\phi$, which implies that $\nabla\phi(\nabla\phi^*(\vw))=\vw$ for all $\vw\in\R^d$, as stated in \cite[Corollary 13.3.1, Theorem 26.5]{rockafellar1997convex}. Furthermore, since the objective function in \eqref{vanilla-SBPG} is strictly convex, there exists a unique solution to the inclusion above. By comparing the two inclusions, we can conclude that $\vu^{k+1}={\vx^{k+1}}$.
 \end{proof}
Based on Proposition \ref{subprob-closed-form-sol} and definition of the definition of $\widetilde{\cG}_\alpha(\vx)$, we can easily observe that $\vx^{k+1}=\vx^k-\alpha_k\widetilde{\cG}_{\alpha_k}(\vx^k)$. We can derive the following proposition, which bounds the difference between $\mathcal{G}_{\alpha}(\vx)$ and $\widetilde{\mathcal{G}}_{\alpha}(\vx)$ directly from Proposition \ref{nonexpansive}. The proof is omitted for brevity.
\begin{proposition}
\label{gradient-map-gap}
Suppose Assumption \ref{relative-smooth-assumption} holds. At the $k$-th step, we have the estimation:
    \[
    \|\mathcal G_{\alpha_k}({\vx^k})-\widetilde{\mathcal G}_{\alpha_k}({\vx^k})\|\leq\frac{1}{\mu}\|\nabla F({\vx^k})-\widetilde{\nabla}_k\|=\frac{\|{\bm\varepsilon}_k\|}{\mu},
    \]
    where 
    {$\bm{\varepsilon}_k=\nabla F({\vx^k})-\widetilde{\nabla}_k.$}
\end{proposition}
Before presenting the main convergence result, we state the following one-step descent lemma below.
\begin{lemma}
Suppose Assumption \ref{relative-smooth-assumption} holds. The sequence generated by SBPG satisfies the following condition: 
\[
\Phi({\vx^{k+1}})\leq\Phi({\vx^k})-\frac{1}{\alpha_k}\cD_\phi({\vx^k},{\vx^{k+1}})-\left(\frac{1}{\alpha_k}-L_F\right)\cD_\phi({\vx^{k+1}},{\vx^k})+\inprod{{\bm\varepsilon}_k}{{\vx^{k+1}}-\vx^{k}}.
\]
\label{basic-descent-lemma-2}
\end{lemma}
\begin{proof}
By the optimality condition of \eqref{vanilla-SBPG}, we obtain that 
\[
0\in\partial R({\vx^{k+1}})+\wtnabla_k+\frac{1}{\alpha_k}\left(\nabla\phi({\vx^{k+1}})-\nabla\phi({\vx^k})\right).
\]
Appealing to the convexity of $R$, we have
\[
R(\vx)-R({\vx^{k+1}})\geq\big\langle{-\wtnabla_k-\frac{1}{\alpha_k}\left(\nabla\phi({\vx^{k+1}})-\nabla\phi({\vx^k})\right)},\,{\vx-{\vx^{k+1}}}\big\rangle.
\]
By the four points identity and the definition of ${\bm\varepsilon}_k$, we get 
\[
 R(\vx)-R({\vx^{k+1}})\geq\frac{1}{\alpha_k}\left[{\cD_\phi}({\vx^{k+1}},{\vx^k})+{\cD_\phi}(\vx,{\vx^{k+1}})-{\cD_\phi}(\vx,{\vx^k})\right]-\inprod{\nabla F({\vx^k})}{\vx-{\vx^{k+1}}}+\inprod{{\bm\varepsilon}_k}{\vx-{\vx^{k+1}}}.
\]
Set $\vx={\vx^k}$ in the above inequality, we have the following inequality:
\[
R({\vx^k})-R({\vx^{k+1}})\geq\frac{1}{\alpha_k}\left[{\cD_\phi}({\vx^{k+1}},{\vx^k})+{\cD_\phi}({\vx^k},{\vx^{k+1}})\right]-\inprod{\nabla F({\vx^k})}{{\vx^k}-{\vx^{k+1}}}+\inprod{{\bm\varepsilon}_k}{{\vx^k}-{\vx^{k+1}}}.
\]
By the smooth adaptivity of $F$, we have 
\[
F({\vx^{k+1}})\leq F({\vx^k})+\inprod{\nabla F({\vx^k})}{{\vx^{k+1}}-{\vx^k}}+L_F{\cD_\phi}({\vx^{k+1}},{\vx^k}).
\]
Combining the above two inequalities above, we complete the proof.
\end{proof}

\subsection{Convergence analysis of SBPG}

 In this subsection, we establish the convergence results for SBPG, which is an extension of the convergence result in \cite{ghadimi2016mini}, in which the lassical Lipschitz gradient assumption is required. In many literature, the bounded sequence assumption is often required in the convergence analysis of stochastic algorithms. However, in this section, we relax this assumption and prove that under a certain condition, the sequence generated by \eqref{vanilla-SBPG} is bounded almost surely. We need the following result to bound the stochastic error term $\inprod{{\bm\varepsilon}_k}{{\vx^{k+1}}-{\vx^k}}$ in Lemma \ref{basic-descent-lemma-2}.
\begin{lemma}
Suppose Assumption \ref{relative-smooth-assumption}, \ref{variance-assumption} hold. We have the following estimation of the error term:
\[
\mathbb E\left[\inprod{{\bm\varepsilon}_k}{{\vx^{k+1}}-{\vx^k}}\right]\leq\frac{\alpha_k}{\mu}\E[\|{\bm\varepsilon}_k\|^2]\leq\frac{\alpha_k\sigma^2}{\mu m_k}.
\]
\label{inprod-bound}
\end{lemma}
\begin{proof}
Define ${\bar\vx}^{k+1}:= {{\rm Prox}_{\alpha_kR}^\phi} (\nabla\phi^*(\nabla\phi({\vx^k})-\alpha_k\nabla F({\vx^k})))$. By Proposition \ref{subprob-closed-form-sol} and 
the optimality condition for {${\bar\vx}^{k+1}$}, we have
\[
0\in\partial R({\bar\vx}^{k+1})+\nabla F({\vx^k})+\frac{1}{\alpha_k}(\nabla\phi({\bar\vx}^{k+1})-\nabla\phi({\vx^k})).
\]
Similarly,
\[
0\in\partial R({\vx^{k+1}})+\wtnabla_k+\frac{1}{\alpha_k}(\nabla\phi({\vx^{k+1}})-\nabla\phi({\vx^k})).
\]
By the monotonicity of $\partial R$ and Lemma \ref{four-point-identity}, we have
\[
\big\langle{{\bar\vx}^{k+1}-{\vx^{k+1}}},\,
{-{\bm\varepsilon}_k-\frac{1}{\alpha_k}(\nabla\phi({\bar\vx}^{k+1})-\nabla\phi({\vx^{k+1}}))}
\big\rangle
\geq0.
\]
Therefore, 
\[
\inprod{{\vx^{k+1}}-{\bar\vx}^{k+1}}{{\bm\varepsilon}_k}\geq
\inprod{{\bar\vx}^{k+1}-{\vx^{k+1}}}{\frac{1}{\alpha_k}(\nabla\phi({\bar\vx}^{k+1})-\nabla\phi({\vx^{k+1}}))}\geq\frac{\mu}{\alpha_k}\|{\bar\vx}^{k+1}-{\vx^{k+1}}\|^2.
\]
By Cauchy-Schwarz inequality, we get $\|{\bar\vx}^{k+1}-{\vx^{k+1}}\|\leq\frac{\alpha_k}{\mu}\|{\bm\varepsilon}_k\|$.  

Now, we are ready to prove Lemma \ref{inprod-bound}. From the definition, we know that ${\bar\vx}^{k+1}\vartriangleleft\cF_k$. Therefore, $\E[\inprod{{\bm\varepsilon}_k}{{\vx^k}-{\bar\vx}^{k+1}}]=\E[\E[\inprod{{\bm\varepsilon}_k}{{\vx^k}-{\bar\vx}^{k+1}}|\cF_k]]=\E[\inprod{\E[{\bm\varepsilon}_k|\cF_k]}{{\vx^k}-{\bar\vx}^{k+1}}]=0$, {where} the first equality is from the tower rule of conditional expectation, the second comes from {the fact} that ${\vx^k}-{\bar\vx}^{k+1}\vartriangleleft\cF_k$. Hence,
{\small
\[
\mathbb E\left[\inprod{{\bm\varepsilon}_k}{{\vx^{k+1}}-\vx^{k}}\right]=\E\left[\inprod{{\bm\varepsilon}_k}{{\vx^{k+1}}-{\bar\vx}^{k+1}}\right]-\E\left[\inprod{{\bm\varepsilon}_k}{{\vx^k}-{\bar\vx}^{k+1}}\right]\leq\frac{\alpha_k}{\mu}\E[\|{\bm\varepsilon}_k\|^2]\leq\frac{\alpha_k\sigma^2}{\mu m_k},
\]}
which completes the proof.
\end{proof}
 
\begin{lemma}[Bounded sequence]
Suppose Assumption \ref{relative-smooth-assumption}, \ref{variance-assumption} hold. If $\sum_k\frac{\alpha_k}{m_k}<\infty$, $\sup_k\alpha_k\leq\bar\alpha<\frac{1}{L_F}$, then,
 \begin{itemize}
     \item[1.] $\sum_{k=0}^\infty \E[\cD_\phi({\vx^{k+1}},{\vx^k})]<\infty$.
     \item[2.] If $\Phi$ is level bounded, then $\{{\vx^k}\}_{k\geq0}$ is bounded almost surely.
 %In particular, $\{{\vx^k}\}_k\leq0$ is bounded.
   \end{itemize}
   \label{bounded sequence-vanilla}
\end{lemma}
\begin{proof}
By Cauchy-Young inequality, we have 
\[
|\inprod{{\bm\varepsilon}_k}{{\vx^k}-{\vx^{k+1}}}|\leq\frac{\mu}{2\alpha_k}\|{\vx^k}-{\vx^{k+1}}\|^2+\frac{\alpha_k}{2\mu}\|{\bm\varepsilon}_k\|^2\leq \frac{1}{\alpha_k}{\cD_\phi}({\vx^k},{\vx^{k+1}})+\frac{\alpha_k}{2\mu}\|{\bm\varepsilon}_k\|^2.
 \]
 By Lemma \ref{basic-descent-lemma-2}, we have
 \begin{equation}
 \left(\frac{1}{\alpha_k}-L_F\right){\cD_\phi}({\vx^{k+1}},{\vx^k})\leq\Phi({\vx^k})-\Phi({\vx^{k+1}})+\frac{\alpha_k}{2\mu}\|{\bm\varepsilon}_k\|^2.
 \label{bound-sequence-recursion}
 \end{equation}
Taking conditional expectation for both sides of \eqref{bound-sequence-recursion}, we get
\[
\E\left[\left(\frac{1}{\alpha_k}-L_F\right){\cD_\phi}({\vx^{k+1}},{\vx^k})|\cF_k\right]\leq\Phi({\vx^k})-\E[\Phi({\vx^{k+1}})|\cF_k]+\frac{\alpha_k}{2\mu}\E[\|{\bm\varepsilon}_k\|^2|\cF_k].
\]
Since $\sum_{k\geq0}\frac{\alpha_k}{2\mu}\E[\|{\bm\varepsilon}_k\|^2|\cF_k]\leq\sum_{k\geq0}\frac{\alpha_k\sigma^2}{2\mu m_k}<\infty$, applying Theorem \ref{martingale-convergence}, we have that $\Phi({\vx^k})$ converges and $\sum_{k\geq0}\E\left[\left(\frac{1}{\alpha_k}-L_F\right){\cD_\phi}({\vx^{k+1}},{\vx^k})|\cF_k\right]<\infty$  almost surely.  By the tower rule of conditional expectation, we have $\sum_{k=0}^\infty \E[{\cD_\phi}({\vx^{k+1}},{\vx^k})]<\infty$. Since $\Phi({\vx^k})$ converges {almost surely}, thus $\{\Phi({\vx^k})\}_{k\geq0}$ is bounded almost surely. By the level boundness of $\Phi$, we deduce that $\{{\vx^k}\}_{k\geq0}$ is bounded almost surely.
\end{proof}

Now, we present our main convergence result for the vanilla SBPG, which is in the sense of expectation.
\begin{theorem}[Convergence result in expectation]
  Suppose Assumption \ref{relative-smooth-assumption}, \ref{variance-assumption} hold, $\alpha_k<\frac{1}{L_F}\min\{1,\frac{1}{\mu}\}$. Define a random variable $r$ with the distribution $\mathbb{P}\{r=k\}=\frac{\alpha_k}{\sum_{k=0}^{N-1}\alpha_k}$ for $k=0,...,N-1$. Then,
\be\label{vSBPG-convergence-thm}
 \E[\|\widetilde{\mathcal G}_{\alpha_r}(\vx^{r})\|^2]\leq\frac{2\Phi(x^0)-2\Phi^*+2\sum_{k=0}^{N-1}\frac{\alpha_k\sigma^2}{\mu m_k}}{\mu\sum_{k=0}^{N-1}\alpha_k}.
\ee
  Moreover, if $\Phi$ is level bounded, $\sum_{k}\frac{\alpha_k}{m_k}<+\infty$ and $\sum_{k}\alpha_k=+\infty$, then the sequence $\{{\vx^k}\}_{k\geq0}$ is bounded almost surely, and the right hand side of \eqref{vSBPG-convergence-thm} converges to zero.
 \label{converge-thm-vanilla-expt}
\end{theorem}
\begin{proof}
    Note that ${\vx^{k+1}}={\vx^k}-\alpha_k\widetilde{\cG}_{\alpha_k}({\vx^k})$ and by the strongly convexity of $\phi$, Lemma \ref{basic-descent-lemma-2} yields 
\[
\begin{aligned}
    \mu(\alpha_k-\frac{L_F\alpha_k^2}{2})\|\widetilde{\cG}_{\alpha_k}({\vx^k})\|^2&\leq\frac{1}{\alpha_k}\cD_\phi({\vx^k},{\vx^{k+1}})+\left(\frac{1}{\alpha_k}-L_F\right)\cD_\phi({\vx^{k+1}},{\vx^k})\\
    &\leq\Phi({\vx^k})-\Phi({\vx^{k+1}})+\inprod{{\bm\varepsilon}_k}{{\vx^{k+1}}-{\vx^k}}.
\end{aligned}
\]
Taking expectations, telescoping from $k=0...N-1$, and using Lemma \ref{inprod-bound}, we obtain
\be
 \sum_{k=0}^{N-1}\mu(\alpha_k-\frac{\mu L_F\alpha_k^2}{2})\E[\|\widetilde{\mathcal G}_{\alpha_k}({\vx^k})\|^2]\leq\Phi(\vx^0)-\Phi(\vx^N)+\sum_{k=0}^{N-1}\frac{\alpha_k\sigma^2}{\mu m_k}.
 \label{proof-vsbpg-thm}
\ee
By utilizing the inequality $\alpha_k - \frac{\mu L_F \alpha_k^2}{2} \geq \frac{\alpha_k}{2}$, the condition $\Phi(\vx^N) \geq \Phi^*$, and considering the definition of the random variable $r$, we can derive \eqref{vSBPG-convergence-thm} from \eqref{proof-vsbpg-thm}. 
\end{proof}


\begin{remark}
We give some remarks for Theorem \ref{converge-thm-vanilla-expt}.
\begin{itemize} 
\item[1.] The mini-batch setting is a crucial component for ensuring the convergence, as it allows us to control the stochastic error term in Lemma \ref{basic-descent-lemma-2} and provide a bound for $\E[\|\widetilde{\cG}_{\alpha_k}({\vx^k})\|^2]$ that converges to zero as $k$ tends to infinity. If $m_k=1$ for all $k$, then the upper bound for $\E[\|\widetilde{\cG}_{\alpha_k}({\vx^k})\|^2]$ will not converge to zero, no matter how $\{\alpha_k\}$ is selected. 
\item[2.] In \cite{ghadimi2016mini}, a similar convergence result is established for nonconvex stochastic proximal gradient methods, but our analysis {differs in a crucial aspect in that} we do not assume the Lipschitz continuity of $F(\vx)$. Instead, we assume that $F(\vx)$ is smooth adaptable, which is a more relaxed assumption. Moreover, we provide a specific choice of stepsizes $\{\alpha_k\}$ and mini-batch sizes $\{m_k\}$ that guarantee 
{the convergence of $\E[\|\widetilde{\cG}_{\alpha_k}({\vx^k})\|^2]$ to $0$},
as well as the almost sure boundedness of the sequence.
% \item[3.] \blue{In Theorem \ref{converge-thm-vanilla-expt}, the random variable $r$ is originally defined assuming prior knowledge of $L_F$. The similar definition of an auxiliary random variable for the output index, requiring known Lipschitz constant, can be found in various references such as \cite[Theorem 2]{ghadimi2016mini}. It should be noted that this requirement of known $L_F$ is not necessary. In fact, when $L_F$ is unknown, we can consider selecting a step size satisfying $\alpha_k<\frac{1}{L_F}\min\{1,\frac{1}{\mu}\}$. In this case, we define a random variable $r$ with the distribution $\mathbb{P}\{r=k\}=\frac{\alpha_k}{\sum_{k=0}^{N-1}\alpha_k}$ for $k=0,...,N-1$. By utilizing \eqref{proof-vsbpg-thm}, we can deduce the following inequality:
% \[
%  \E[\|\widetilde{\mathcal G}_{\alpha_r}(\vx^{r})\|^2]\leq\frac{2\Phi(x^0)-2\Phi^*+2\sum_{k=0}^{N-1}\frac{\alpha_k\sigma^2}{\mu m_k}}{\mu\sum_{k=0}^{N-1}\alpha_k}.
% \]
% It's worth noting that the definition of is independent of $\mu$ and $L_F$.}
\end{itemize}
\label{Vanilla-convergence-expt-remark}
\end{remark}

Based on Proposition \ref{gradient-map-gap} and Theorem \ref{converge-thm-vanilla-expt}, we can derive the following convergence result using the measure ${\cG}_{\alpha_r}(\vx^{r})$. This can be obtained by observing that $\|\cG_{\alpha_r}(\vx^r)\|^2\leq2\|\widetilde\cG_{\alpha_r}(\vx^r)\|^2+2\|\cG_{\alpha_r}(\vx^r)-\widetilde\cG_{\alpha_r}(\vx^r)\|^2.$
\begin{corollary}
Under the conditions in Theorem \ref{converge-thm-vanilla-expt}, we have
\[
\E[\|{\mathcal G}_{\alpha_r}(\vx^{r})\|^2]
\;\leq\;
\frac{4\Phi(\vx^0)-4\Phi^*+4\sum_{k=0}^{N-1}\frac{\alpha_k\sigma^2}{\mu m_k}}{\sum_{k=0}^{N-1}\mu\alpha_k}+\frac{2\sigma^2}{\mu^2m_k}.
 \]
\end{corollary}

\begin{comment}
\subsection{Almost sure convergence results}
We can establish almost sure convergence under certain conditions. In particular, Lemma \ref{bounded sequence-vanilla} provides the conditions under which $\{{\vx^k}\}_{k\geq0}$ is bounded almost surely. We define $\Omega$ to be the limit point set of $\{{\vx^k}\}_{k\geq0}$. 
\begin{theorem}
 Given a $\mu$-strongly convex kernel function $\phi$, the Bregman distance satisfies the following condition: for any $\vx\in dom\,\phi$, $\vy\in int\,dom\,\phi$, $\cD_\phi(\vx,vy)\rightarrow0$ whenever $x\rightarrow y$. Suppose $\Phi$ is level bounded, $\sup_k\alpha_k\equiv\alpha<\frac{1}{L_F}$ and $\sum_{k\geq0}\frac{1}{m_k}<\infty$, hen, the following results hold almost surely:
 \begin{itemize}
     \item[1.] $\sum_{k\geq0}\|\mathcal G_{\alpha}({\vx^k})\|^2<\infty$, $\|\mathcal G_{\alpha}({\vx^k})\|\rightarrow0$.
     \item[2.] Any limit point $\bar x\in\Omega$ is a critical point.
     \item[3.] $\Phi({\vx^k})$ converges to some value $\Phi_*$. For any $\bar x\in\Omega$, $\Phi(\bar x)=\Phi_*$.
 \end{itemize}
 \label{vanilla-as-convergence}
\end{theorem}
\begin{proof}
    From \eqref{bound-sequence-recursion}, we can derive the inequality
    \[
    \mu(\alpha-\alpha^2L_F)\|\widetilde{\mathcal{G}}_\alpha({\vx^k})\|^2\leq(\frac{1}{\alpha}- L_F)\cD_\phi({\vx^{k+1}},{\vx^k})\leq\Phi({\vx^k})-\Phi({\vx^{k+1}})+\frac{\alpha}{2\mu}\|{\bm\varepsilon}_k\|^2.
    \]
   Then, by applying Proposition \ref{gradient-map-gap}, we obtain
    \[
    \mu(\alpha-\alpha^2L_F)\|{\mathcal{G}}_\alpha({\vx^k})\|^2\leq2\left(\Phi({\vx^k})-\Phi({\vx^{k+1}})+\frac{\alpha}{2\mu}\|{\bm\varepsilon}_k\|^2\right)+\frac{2(\alpha-\alpha^2L_F)}{\mu}\|{\bm\varepsilon}_k\|^2.
    \]
    Since $\mathcal{G}_\alpha({\vx^k})\vartriangleleft\cF_k$, we can take the conditional expectation on both sides of the above inequality to obtain
    \[
    \E[2\Phi({\vx^{k+1}})|\cF_k]\leq2\Phi({\vx^k})-\mu(\alpha-\alpha^2L_F)\|{\mathcal{G}}_\alpha({\vx^k})\|^2+\frac{(4\alpha-2\alpha^2L_F)\sigma^2}{\mu m_k}.
    \]
    By applying Theorem \ref{martingale-convergence}, we conclude that $\sum_{k\geq0}\|{\mathcal{G}}_\alpha({\vx^k})\|^2<\infty$ almost surely. Consequently, $\cG_{\alpha}({\vx^k})\rightarrow0$ almost surely, which verifies Claim 1.

    Since $\sum_k\|\mathcal G_\alpha({\vx^k})\|^2<\infty$ a.s., then $\mathcal G_\alpha({\vx^k})\rightarrow0$ a.s. For any $\bar x\in\Omega$, there exists $\{x^{k_j}\}$ such that $ x^{k_j}\rightarrow\bar x$, 
\[
0=\lim_{k_j\rightarrow\infty}\|\mathcal G_\alpha(x^{k_j})\|\geq\|\mathcal G_\alpha(\bar x)\|,\;\text{{(continuity of $\mathcal{G}_\alpha(x)$)}. }
\]
Therefore, $\bar x$ is a stationary point, which proves Claim 2.

Let $x=\bar x$ in Lemma \ref{basic-descent-lemma-1},
$$
 R(\bar x)-R({\vx^{k+1}})\geq\frac{1}{\alpha_k}\left[\cD_\phi({\vx^{k+1}},{\vx^k})+\cD_\phi(\bar x,{\vx^{k+1}})-\cD_\phi(\bar x,{\vx^k})\right]-\inprod{\nabla F(\vx^k)}{\bar \vx-{\vx^{k+1}}}+\inprod{{\bm\varepsilon}_k}{{\vx^{k+1}}-\bar x},
 $$
 Using $L_F$-smoothness of $F$, we have:
 \[
 \begin{aligned}
 F({\vx^{k+1}})&\leq F({\vx^k})+\inprod{\nabla F({\vx^k})}{{\vx^{k+1}}-{\vx^k}}+L_F\cD_\phi({\vx^{k+1}},{\vx^k})\\
 F(\bar x)&\geq F({\vx^k})+\inprod{\nabla F({\vx^k})}{\bar x-{\vx^k}}-L_F\cD_\phi(\bar x,{\vx^k}),
 \end{aligned}
 \]
 which implies:
 \[
 F(\bar x)-F({\vx^{k+1}})\geq\inprod{\nabla F({\vx^k})}{\bar \vx-{\vx^{k+1}}}-L\cD_\phi({\vx^{k+1}},{\vx^k})-L_F\cD_\phi(\bar x,{\vx^k}).
 \]
  Therefore,
 \[
 \Phi(\bar x)-\Phi({\vx^{k+1}})\geq \frac{1}{\alpha_k}\left[\cD_\phi({\vx^{k+1}},{\vx^k})+\cD_\phi(\bar x,{\vx^{k+1}})-\cD_\phi(\bar x,{\vx^k})\right]-L_F\cD_\phi({\vx^{k+1}},{\vx^k})-L_F\cD_\phi(\bar x,{\vx^k})+\inprod{{\bm\varepsilon}_k}{{\vx^{k+1}}-\bar x}.
 \]
 Note the $\alpha_k\leq\frac{1}{L_F}$, then we obtain
 \[
 \E[\Phi({\vx^{k+1}})-\Phi(\bar x)|\cF_k]\leq\frac{1}{\alpha_k}\left(\cD_\phi(\bar x,{\vx^k})-\E[\cD_\phi(\bar x,{\vx^{k+1}})|\cF_k]\right)+L_F\cD_\phi(\bar x,{\vx^k})+\frac{\alpha\sigma^2}{\mu m_k}.
 \]
 Removing the non-positive term in the right-hand-side, we have
 \begin{equation}
 \E[\Phi({\vx^{k+1}})-\Phi(\bar x)|\cF_k]\leq\frac{1}{\alpha}\cD_\phi(\bar x,{\vx^k})+L_F\cD_\phi(\bar x,{\vx^k})+\frac{\alpha\sigma^2}{\mu m_k}.
 \label{vanilla-as-proof-1}
 \end{equation}
 The conditions in Lemma \ref{bounded sequence-vanilla} are clearly satisfied, and therefore $\Phi({\vx^k})$ converges. Assume $x^{k_j}\rightarrow\bar x$, then applying \eqref{vanilla-as-proof-1} yields:
 \[
 \begin{aligned}
 \E[\Phi_*-\Phi(\bar x)]&=\E[\lim_{k_j\rightarrow\infty}\Phi(x^{k_j+1})-\Phi(\bar x)]\overset{\text{Fatou's lemma}}{\leq}\lim_{k_j\rightarrow\infty}\E[\Phi(x^{k_j+1})-\Phi(\bar x)]\\
  &\leq(\frac{1}{\alpha}+L_F)\lim_{k_j\rightarrow\infty}\E[\cD_\phi(\bar x,x^{k_j})]+\lim_{k_j\rightarrow\infty}\frac{\alpha\sigma^2}{\mu m_k}\\
 &=(\frac{1}{\alpha}+L_F)\E[\lim_{k_j\rightarrow\infty}\cD_\phi(\bar x,x^{k_j})]+\lim_{k_j\rightarrow\infty}\frac{\alpha\sigma^2}{\mu m_k}=0.
 \end{aligned}
 \]
The second equality holds by appealing to the dominated convergence theorem. Since $x^{k_j}\rightarrow\bar x$, we have that $\cD_\phi(\bar x,x^{k_j})\rightarrow0$, and therefore the sequence $\{\cD_\phi(\bar x,x^{k_j})\}_{k_j}$ is bounded \red{[additional assumption is required]}. This allows us to apply the dominated convergence theorem. So, $\Phi_*-\Phi(\bar x)=\E[\Phi_*-\Phi(\bar x)]\leq0$.
 On the other hand,
 \[
\Phi_*=\lim_{k_j\rightarrow\infty}\Phi(x^{k_j})\overset{l.s.c.}{\geq}\Phi(\lim_{k_j\rightarrow\infty} x^{k_j})=\Phi(\bar x),
 \]
which implies $\Phi(\bar x)=\Phi_*$, thereby proving Claim 3.
\end{proof}

\subsection{Global convergence under KL property}
\label{KL-vanilla-section}
In this section, we establish global convergence results for vanilla SBPG applied to nonconvex problems under the KL property assumption. To establish global convergence of a first order method using KL property, a key requirement is the sufficient descent lemma, which requires the identification of a decreasing Lyapunov function. However, this is difficult to achieve in vanilla stochastic first order methods due to the noise in the gradients, which causes oscillations in all quantities unless additional techniques are employed.

To the best of our knowledge, there is currently limited literature on investigating the KL property in the context of stochastic gradient descent methods without sufficient descent property, except for a recent preprint \cite{chouzenoux2023kurdyka}. However, our analysis differs from \cite[4.3]{chouzenoux2023kurdyka} as we rely on a different sufficient descent property and bounded subgradient property. Moreover, we provide a concrete setting for vanilla SBPG that satisfies the conditions, while \cite{chouzenoux2023kurdyka} provides a variance reduction example \cite{schmidt2013fast}, where a strong growth condition is required for each individual function. To exploit the KL property, we require three conditions to hold for the sequence generated by the algorithm: a sufficiently descent condition, a bounded subgradient condition, and a continuity condition. However, in stochastic algorithms, we can only expect a "nearly" sufficient descent condition, which presents a significant challenge, even in the inexact deterministic method. The continuity condition has already been proven in Theorem \ref{vanilla-as-convergence}.


\begin{lemma}
    Suppose the conditions in Theorem \ref{vanilla-as-convergence} hold, except that $\sum_{k\geq0}\frac{1}{{m_k}}<\infty$ is replaced by $\sum_{k\geq0}\frac{1}{\sqrt{m_k}}<\infty$. Additionally, assume that $\phi$ is locally Lipschitz smooth and $dom \phi=\R^d$. Then, there exist positive constants $\alpha$ and $\beta$, as well as positive ${\epsilon_k}$ such that $\sum_{k\geq0}\epsilon_k<\infty$, and the following hold almost surely:
    \begin{itemize}
        \item[1] (Nearly sufficiently descent)
        \[
        \Phi({\vx^{k+1}})\leq\Phi({\vx^k})-\alpha\|{\vx^{k+1}}-{\vx^k}\|^2+\epsilon_k.
        \]
        \item[2] (Bounded subgradient)
        \[
        \|w^{k+1}\|\leq\beta\|{\vx^{k+1}}-{\vx^k}\|+\epsilon_k,\;\forall\;w^{k+1}\in\partial\Phi({\vx^{k+1}}).
        \]
    \end{itemize}
    \label{three-conditions-KL}
\end{lemma}


\begin{remark}
In the inexact deterministic method, one can control the error by increasing the accuracy of the subproblem solved. For example, the popular relative criterion $\epsilon_k\leq\eta_k\|{\vx^{k+1}}-{\vx^k}\|^2$ can be employed. In this case, Claim 1 in Lemma \ref{three-conditions-KL} becomes the sufficient descent property when $\eta_k\leq\alpha$, meaning that the objective function is non-increasing. However, the situation in stochastic methods is different because we can not ensure that for any $k$, $\epsilon_k\leq\eta_k\|{\vx^{k+1}}-{\vx^k}\|^2$ due to the unknown stochasticity. This is why we consider an absolute inexact criterion here.
\end{remark}
Because this paper does not focus on KL property in stochastic algorithms, as a preliminary study, we only consider the smooth case, as in \cite{chouzenoux2023kurdyka}. Therefore, for the rest of this subsection, we assume $R(\vx)=0$, i.e., $\Phi$ is smooth. We define the following event:
\[
{\bm\xi}:=\left\{\omega:\;\Phi({\vx^k})<\Phi_*\;for\;all\;k\geq0 \right\},
\]
where $\Phi_*$ is defined in Theorem \ref{vanilla-as-convergence} as the limit of the function value. We can assume that the event ${\bm\xi}$ has a positive probability, which is reasonable because as the iterates get closer to a critical point, the function value should become more stable and less likely to oscillate. Increasing the mini-batch size can help achieve this stability. While we cannot guarantee that ${\bm\xi}$ has positive probability in practice, analyzing the algorithm under this assumption can still provide insights into its behavior. We have the following global convergence result under KL property.
\begin{theorem}
Suppose the conditions in Lemma \ref{three-conditions-KL} hold, and $\Phi$ is a KL function. Under the event ${\bm\xi}$ with positive probability, the following results hold almost surely:
\begin{itemize}
    \item[1.] $\sum_{k\geq0}\|{\vx^{k+1}}-{\vx^k}\|<\infty$.
    \item[2.] The sequence $\{{\vx^k}\}$ converges to a critical point.
\end{itemize}
\label{vanilla-global-convergence}
\end{theorem}
\begin{proof}
Define $\eta_k=\sum_{i\geq k}\epsilon_i$, $\hat\Phi_k=\Phi({\vx^k})+\eta_k$ and $\bar\Phi_k=\Phi_*+\eta_k$, then $\lim_{k\rightarrow\infty}\eta_k=0$. From Claim 1 in Lemma \ref{three-conditions-KL}, we have that $\hat\Phi_{k+1}\leq\hat\Phi_{k}-\alpha\|{\vx^{k+1}}-{\vx^k}\|^2$. Note that $\varphi$ is concave function and $\hat\Phi_k-\bar\Phi_k=\Phi({\vx^k})-\Phi_*$, so
\[
\begin{aligned}
 \varphi(\hat\Phi_k-\bar\Phi_k)-\varphi(\hat\Phi_{k+1}-\bar\Phi_{k+1})&\geq\varphi^\prime(\hat\Phi_k-\bar\Phi_k)(\hat\Phi_k-\hat\Phi_{k+1})\\
 &\geq\alpha\varphi^\prime(\Phi({\vx^k})-\Phi_*)\|{\vx^{k+1}}-{\vx^k}\|^2\\
 &\geq\frac{\alpha\|{\vx^{k+1}}-{\vx^k}\|^2}{\|\nabla\Phi({\vx^k})\|}\geq\frac{\alpha\|{\vx^{k+1}}-{\vx^k}\|^2}{\beta\|x^{k}-\v\vx^{k-1}\|+\epsilon_k},
 \end{aligned}
\]
where $\varphi^\prime(\Phi({\vx^k})-\Phi_*)$ is well-defined on event ${\bm\xi}$. The rest of the proof is almost standard, see for example, \cite{attouch2009convergence,attouch2013convergence}, except that there is an extra inexact term $\epsilon_k$ in denominator. Denote $\Delta_{p,q}:=\varphi(\hat\Phi_p-\bar\Phi_p)-\varphi(\hat\Phi_{q}-\bar\Phi_{q})$, so $\|{\vx^{k+1}}-{\vx^k}\|^2\leq\frac{\beta}{\alpha}\|{\vx^k}-\v\vx^{k-1}\|\Delta_{k,k+1}+\frac{\epsilon_k}{\alpha}\Delta_{k,k+1}$. Then, by $\sqrt{a+b}\leq\sqrt{a}+\sqrt{b}$ and Cauchy-Young inequality, we have the inequality:
\[
\begin{aligned}
 \|{\vx^{k+1}}-x^{k}\|\leq\frac{1}{2}\|{\vx^k}-\vx^{k-1}\|+\frac{\beta}{\alpha}\Delta_{k,k+1}+\frac{1}{2\beta}\epsilon_k+\frac{\beta}{\alpha}\Delta_{k,k+1}.   
\end{aligned}
\]
For any $\ell$ and $N$ with $\ell<N$, 
\[
\sum_{k=\ell}^{N-1}\|{\vx^{k+1}}-x^{k}\|\leq\frac{1}{2}\sum_{k=\ell}^{N-1}\|{\vx^k}-\vx^{k-1}\|+\frac{2\beta}{\alpha}\Delta_{\ell,N}+\frac{1}{2\beta}\sum_{k=\ell}^{N-1}\epsilon_k, 
\]
then, 
\[
\sum_{k=\ell}^{N-1}\|{\vx^{k+1}}-x^{k}\|\leq\|x^\ell-x^{\ell-1}\|+\frac{4\beta}{\alpha}\Delta_{\ell,N}+\frac{1}{\beta}\sum_{k=\ell}^{N-1}\epsilon_k, 
\]
Note that $\Delta_{\ell,N}\leq\varphi(\hat\Phi_\ell-\bar\Phi_\ell)$ and let $N$ tend to infinity. Then we have $\sum_{k=0}^\infty\|{\vx^{k+1}}-{\vx^k}\|<\infty$. Therefore, ${{\vx^k}}_{k}$ is a Cauchy sequence and must converge. By Claim 2 in Theorem \ref{vanilla-as-convergence}, we can conclude that ${\vx^k}$ converges to some critical point. This completes the proof.
\end{proof}

\end{comment}
\subsection{Momentum based Stochastic Bregman Gradient Descent Method}
Remark \ref{Vanilla-convergence-expt-remark} suggests that increasing the mini-batch size $m_k$ is almost necessary for the error bound in Theorem \ref{converge-thm-vanilla-expt} to converge to zero. However, using a large mini-batch size in each iteration can be computationally expensive in the modern large-scale problems, e.g. training deep neural network. In this part, we resort to the momentum technique to address this issue. Specifically, we consider using a stochastic moving average estimator (SMAE) for the true gradient given by:
\be
\vv^k=(1-\beta_k)\vv^{k-1}+\beta_k\widetilde\nabla_k,\quad \mbox{where}\quad \E[\widetilde\nabla_k|\cF_k]=\nabla F({\vx^k}),
\label{SMAE}
\ee
where $\vv^{k-1}$ can be viewed as the momentum which contain the information of all historical stochastic gradients, and $\E[\|\widetilde\nabla_k\|^2|\cF_k]\leq\frac{\sigma^2}{m_k}$. We expect that incorporating the SMAE technique can achieve a certain level of variance reduction without increasing the mini-batch size. In our approach, we utilize the gradient estimator $\vv^k$ within SBPG, and we refer to the resulting method as MSBPG. Specifically, we consider the following update scheme:
\begin{equation}
{\vx^{k+1}}=\argmin_{\vx\in\overline{C}}\;R(\vx)+\inprod{\vv^k}{\vx-{\vx^k}}+\frac{1}{\alpha_k}\cD_\phi(\vx,{\vx^k}).
\label{momentum-SBPG}
\end{equation}
We need the following assumption that the difference of gradients of $F$ can be bounded by the Bregman distance. 
\begin{assumption}\label{momentum-assumption}
 There exists $\kappa>0$, such that $\|\nabla F(\vx)-\nabla F(\vy)\|^2\leq\kappa \cD_\phi(\vx,\vy)$ for all $\vx\in dom\,\phi$, $\vy\in int\,dom\,\phi$.   
\end{assumption}
\begin{remark}
  This assumption generalizes the {case of Lipschitz kernel function.} Assume that $F$ is $L_F$ smooth adaptable to $\phi$, it can be easily shown that if $\phi$ has $L_\phi$-Lipschitz gradient, this assumption holds for $\kappa\geq\frac{2L_F^2L_\phi^2}{\mu}$. In this paper, we are more interested in polynomial kernel functions. For functions with polynomially bounded growth rates, this assumption is not restrictive. For example, consider the one-dimensional objective function $F(x)=\frac{1}{4}x^4$ and the kernel function $\phi(x)=\frac{1}{2}x^2+\frac{1}{8}x^8$. Then, by \cite[Proposition 2.1]{lu2018relatively}, we know that $F$ is smooth adaptable with respect to $\phi$. Simple algebra shows that $\cD_\phi(x,y)=\frac{1}{8}(x-y)^2(x^6 + 2 x^5 y + 3 x^4 y^2 + 4 x^3 y^3 + 5 x^2 y^4 + 6 x y^5 + 7 y^6 + 4)$ and $(F'(x)-F'(y))^2=(x-y)^2(x^2+xy+y^2)^2$. Numerical computation shows that $(x^6 + 2 x^5 y + 3 x^4 y^2 + 4 x^3 y^3 + 5 x^2 y^4 + 6 x y^5 + 7 y^6 + 4)-(x^2+xy+y^2)^2\geq3.71$. Therefore, $(F'(x)-F'(y))^2\leq8\cD_\phi(x,y)$, which holds globally
  {for any $x$ and $y$}. 
\end{remark}

Next we present a recursion lemma that allows us to estimate the accuracy of the SMAE. While similar lemmas have been proposed in the literature, such as in \cite{wang2017stochastic}, their bounds are not directly applicable in the Bregman setting. As a result, we have developed a version of the recursion lemma that is tailored to our specific context. %The proof of this lemma is included in Appendix \ref{appendix-SBPG-proof}.
\begin{lemma}
\label{SMAE-recursion-lemma}
The following recursion holds
{\small
\[
\E[\|\vv^k-\nabla F({\vx^k})\|^2|\cF_k]\leq (1-\beta_k)\|\vv^{k-1}-\nabla F(\vx^{k-1})\|^2+\beta_k^2\E[\|\widetilde\nabla_k-\nabla F({\vx^k})\|^2|\cF_k]+\frac{\|\nabla F(\vx^{k-1})-\nabla F({\vx^k})\|^2}{\beta_k}.
\]
}
\end{lemma}
\begin{proof}
Note that $\vv^k-\nabla F(\vx^k)=(1-\beta_k)(\vv^{k-1}-\nabla F(\vx^{k-1}))+(1-\beta_k)(\nabla F(\vx^{k-1})-\nabla F(\vx^k))+\beta_k(\widetilde\nabla_k-\nabla F(\vx^k))$, and $\E[\widetilde\nabla_k-\nabla F(\vx^k)|\cF_k]=0$. Then we have
\[
\begin{aligned}
&\E[\|\vv^k-\nabla F(\vx^k)\|^2|\cF_k]\\
={}&(1-\beta_k)^2\|\vv^{k-1}-\nabla F(\vx^{k-1})\|^2+(1-\beta_k)^2\|\nabla F(\vx^{k-1})-\nabla F(\vx^k)\|^2+\\&\beta_k^2\E[\|\widetilde\nabla_k-\nabla F(\vx^k)\|^2|\cF_k]+2(1-\beta_k)^2\inprod{\vv^{k-1}-\nabla F(\vx^{k-1})}{\nabla F(\vx^{k-1})-\nabla F(\vx^k)}\\
\leq{}&(1-\beta_k)^2\|\vv^{k-1}-\nabla F(\vx^{k-1})\|^2+(1-\beta_k)^2\|\nabla F(\vx^{k-1})-\nabla F(\vx^k)\|^2+\\
&\beta_k^2\E[\|\widetilde\nabla_k-\nabla F(\vx^k)\|^2|\cF_k]+\beta_k(1-\beta_k)\|\vv^{k-1}-\nabla F(\vx^{k-1})\|^2+\frac{(1-\beta_k)^3}{\beta_k}\|\nabla F(\vx^{k-1})-\nabla F(\vx^k)\|^2\\
={}&(1-\beta_k)\|\vv^{k-1}-\nabla F(\vx^{k-1})\|^2+\beta_k^2\E[\|\widetilde\nabla_k-\nabla F(\vx^k)\|^2|\cF_k]+\frac{(1-\beta_k)^2\|\nabla F(\vx^{k-1})-\nabla F(\vx^k)\|^2}{\beta_k}\\
\leq{}&(1-\beta_k)\|\vv^{k-1}-\nabla F(\vx^{k-1})\|^2+\beta_k^2\E[\|\widetilde\nabla_k-\nabla F(\vx^k)\|^2|\cF_k]+\frac{\|\nabla F(\vx^{k-1})-\nabla F(\vx^k)\|^2}{\beta_k}.
\end{aligned}
\]
This completes the proof.
\end{proof}

Now we are ready to provide the convergence result for our momentum based SBPG. 
\begin{theorem}
Suppose Assumption \ref{relative-smooth-assumption}, \ref{variance-assumption} and \ref{momentum-assumption} hold. Let $\alpha_k=c\mu\beta_{k+1}$ for any $c\in(0,\frac{1}{2\sqrt{\mu\kappa}}]$. Then
\[
\E[\|\widetilde{\mathcal G}_{\alpha_r}(\vx^{r})\|^2]\leq\frac{\Phi^0-\Phi^*+c\|\vv_0-\nabla F(\vx^0)\|^2+c\sum_{k=1}^N\frac{\beta_k^2\sigma^2}{m_k}}{\sum_{k=0}^{N-1}\frac{\mu\alpha_k}{8}}=\mathcal{O}\left(\frac{1}{\sum_{k=0}^{N-1}\alpha_k}+\frac{\sum_{k=0}^{N-1}\frac{\alpha_k^2}{m_{k+1}}}{\sum_{k=0}^{N-1}\alpha_k}\right),
\]
where $r$ is a random variable with distribution $\mathbb{P}\{r=k\}=\frac{\alpha_k}{\sum_{k=0}^{N-1}\alpha_k}$, for $k=0,...,N-1$.
\label{momentum-convergence-thm}
\end{theorem}
\begin{proof}
 From Lemma \ref{basic-descent-lemma-2} and Cauchy-Young inequality, we have
\[
\begin{aligned}
\Phi({\vx^{k+1}})&\leq\Phi({\vx^k})-\frac{1}{\alpha_k}{\cD_\phi}({\vx^k},{\vx^{k+1}})+\frac{\mu}{4\alpha_k}\|{\vx^k}-{\vx^{k+1}}\|^2+\frac{\alpha_k}{\mu}\|{\bm\varepsilon}_k\|^2\\
&\leq\Phi({\vx^k})-\frac{1}{2\alpha_k}{\cD_\phi}({\vx^k},{\vx^{k+1}})+\frac{\alpha_k}{\mu}\|{\bm\varepsilon}_k\|^2.
\end{aligned}
\]
where we have defined ${\bm\varepsilon}_k:=\nabla F({\vx^k})-\vv^k$. Summing the above inequality over $k=0,\ldots,N-1$ and rearranging the terms, we get
\[
\sum_{k=0}^{N-1}\frac{1}{2\alpha_k}{\cD_\phi}({\vx^k},{\vx^{k+1}})\leq\Phi(\vx^0)-\Phi^*+\sum_{k=0}^{N-1}\frac{\alpha_k}{\mu}\|{\bm\varepsilon}_k\|^2.
\]
By applying Lemma \ref{SMAE-recursion-lemma}, we can obtain the following inequality:
{
\[
\beta_k\E[\|{\bm\varepsilon}_{k-1}\|^2]
\;\leq\; \E[\|{\bm\varepsilon}_{k-1}\|^2] - \E[\|{\bm\varepsilon}_{k}\|^2]
+\beta_k^2\E[\|\widetilde\nabla_k-\nabla F({\vx^k})\|^2] 
+\E\left[\frac{\|\nabla F(\vx^{k-1})-\nabla F({\vx^k})\|^2}{\beta_k}\right].
\]}
Hence
\small{
\[
\sum_{k=0}^{N-1}\beta_{k+1}
\E[\|{\bm\varepsilon}_{k}\|^2]=\sum_{k=1}^N\beta_k\E[\|{\bm\varepsilon}_{k-1}\|^2]\leq\|{\bm\varepsilon}_0\|^2+\sum_{k=1}^N\beta_k^2\E[\|\widetilde\nabla_k-\nabla F({\vx^k})\|^2]+\sum_{k=1}^N\E\left[\frac{\|\nabla F(\vx^{k-1})-\nabla F({\vx^k})\|^2}{\beta_k}\right].
\]}
Since $\frac{\alpha_k}{\mu}=c\beta_{k+1}$ for some constant $c$, we get the following inequality:
\small{
\[
\sum_{k=0}^{N-1}\frac{1}{2\alpha_k}\E[{\cD_\phi}({\vx^k},{\vx^{k+1}})]\leq\Phi(x^0)-\Phi^*+c\left(\|{\bm\varepsilon}_0\|^2+\sum_{k=1}^N\beta_k^2\E[\|\widetilde\nabla_k-\nabla F({\vx^k})\|^2]+\sum_{k=1}^N\E\left[\frac{\|\nabla F(\vx^{k-1})-\nabla F({\vx^k})\|^2}{\beta_k}\right]\right).
\]
}
By using Assumption \ref{momentum-assumption}, we obtain that
\[
\frac{\|\nabla F({\vx^k})-\nabla F({\vx^{k+1}})\|^2}{\beta_{k+1}}\leq\frac{\kappa}{\beta_{k+1}}{\cD_\phi}({\vx^k},{\vx^{k+1}}).
\]
Combining above two inequalities, we get
\[
\sum_{k=0}^{N-1}\frac{1}{2\alpha_k}\E[{\cD_\phi}({\vx^k},{\vx^{k+1}})]\leq\Phi(\vx^0)-\Phi^*+c\left(\|{\bm\varepsilon}_0\|^2+\sum_{k=1}^N\beta_k^2\E[\|\widetilde\nabla_k-\nabla F({\vx^k})\|^2]+\sum_{k=0}^{N-1}\frac{\kappa}{\beta_{k+1}}\E[{\cD_\phi}({\vx^k},{\vx^{k+1}})]\right).
\]
{Since $c\leq\frac{1}{2\sqrt{\mu\kappa}}$ and $\frac{\alpha_k}{\mu}=c\beta_{k+1}$, we can deduce that 
 $\frac{c\kappa}{\beta_{k+1}}\leq \frac{1}{4\alpha_k}$}. Using this condition, we obtain the inequality:
\[
\sum_{k=0}^{N-1}\frac{1}{4\alpha_k}\E[{\cD_\phi}({\vx^k},{\vx^{k+1}})]\leq\Phi(\vx^0)-\Phi^*+c\left(\|{\bm\varepsilon}_0\|^2+\sum_{k=1}^N\beta_k^2\E[\|\widetilde\nabla_k-\nabla F({\vx^k})\|^2]\right). 
\]
Note that ${\cD_\phi}({\vx^k},{\vx^{k+1}})\geq\frac{\mu}{2}\|{\vx^k}-{\vx^{k+1}}\|^2=\frac{\mu\alpha_k^2}{2}\|\widetilde{\mathcal{G}}_{\alpha_k}({\vx^k})\|^2$ and by the definition of the random variable $a$, we get
\[
\E[\|\widetilde{\mathcal G}_{\alpha_a}(\vx^{a})\|^2]\leq\frac{\Phi^0-\Phi^*+ {c\|\bm{\varepsilon}_0\|^2}
+c\sum_{k=1}^N\frac{\beta_k^2\sigma^2}{m_k}}{\sum_{k=0}^{N-1}\frac{\mu\alpha_k}{8}}=\mathcal{O}\left(\frac{1}{\sum_{k=0}^{N-1}\alpha_k}+\frac{\sum_{k=0}^{N-1}\frac{\alpha_k^2}{m_{k+1}}}{\sum_{k=0}^{N-1}\alpha_k}\right),
\]     
which completes the proof.
\end{proof}

\begin{remark}
Now we give some remarks for Theorem \ref{momentum-convergence-thm}.
\begin{itemize}
\item[1.] When the sequence $\{\vx_k\}$ is bounded, an alternative assumption to Assumption \ref{momentum-assumption} is that $C=\R^d$ and $\phi$ has a locally Lipschitz gradient, as made in \cite[Theorem 4.1]{bolte2018first} and \cite[Theorem 4.7]{latafat2022bregman}. Under these assumptions, we can conclude that there exists a compact set $\cU$ containing $\{\vx_k\}$. Therefore, there exists $L_{\phi,\cU}>0$ such that $\nabla\phi$ is Lipschitz continuous over $\cU$, and we can derive that $\|\nabla F(\vx)-\nabla F(\vy)\|^2\leq L_F^2L_{\phi,\cU}^2\|\vx-\vy\|^2\leq\frac{2L_F^2L_{\phi,\cU}^2}{\mu}\cD_\phi(\vx,\vy)$ holds. %However, we Assumption \ref{momentum-assumption} does not require boundness assumption.
\item[2.] Compared to the results presented in Theorem \ref{converge-thm-vanilla-expt}, it is worth noting that even when keeping $m_k$ constant, we can still achieve convergence to zero error bound by carefully selecting the stepsize sequence $\{\alpha_k\}$. A typical stepsize condition is that $\sum_{k}\alpha_k=\infty,\;\sum_{k}\alpha_k^2<\infty$, which coincides with the classical stepsize condition that guarantees a sufficient but not too fast decrease of the stepsize, as discussed in \cite{bertsekas2000gradient}.  Therefore, by incorporating the momentum technique, we can achieve an improved convergence with negligible additional computation costs. This desirable convergence property theoretically supports the use of SBPG with momentum for large-scale problems, such as deep neural networks, without the need to increase the mini-batch size. 
\end{itemize}
\end{remark}

\section{Application of MSBPG in deep neural networks}\label{DNN implementation}
In this section, we present a detailed description of MSBPG applied to training deep neural networks. Throughout this section, we assume that the optimization domain $\overline{C}$ is the entire space $\R^d$, so that $\phi\in\cM(\R^d)$ and $F\in\cC^1(\R^d)$. For simplicity, we omit the explicit mention of the feasible set $\R^d$ in this section. The optimization problem we consider here is given by:
\be 
\min_{\mW}\;\underbrace{\frac{1}{N}\sum_{i=1}^N\cL(\mathcal{DNN}(\mW,\vx_i),y_i)}_{F(\mW)}+\lambda\|\mW\|_1,
\label{DNN-loss-model}
\ee 
where $\mathcal{DNN}({\mW},\vx)$ is the neural network function with training parameters $\mW$ and input data $\vx$, $\mathcal{L}$ is the loss function that measures the difference between the output of the neural network $\mathcal{DNN}({\mW},\vx_i)$ and the label $y_i$, $F(\mW)$ is the training loss evaluated on the training dataset $\{(\vx_i,y_i)\}_{i=1}^N$, and $\lambda \|\mW\|_1$ is the $L_1$ regularization term that is often used to avoid overfitting in training deep neural networks \citep{ng2004feature}. To illustrate the neural network function $\mathcal{DNN}({\mW},\vx)$, in the $L$-layer fully connected neural network, we have $\mW=[\mW_1,\mW_2,\cdots,\mW_L]$ and
\begin{equation}
\mathcal{DNN}(\mW,\vx)=\sigma_L(\mW_L(\sigma_{L-1}(\mW_{L-1}(...(\sigma_1(\mW_1\vx))...)))),
\label{DNN-def}
\end{equation}
where $\sigma_i$ is the nonlinear activation function. In this paper, we focus on smooth activation functions.

At the $k$-th iteration, MSBPG has the following update scheme:
\begin{eqnarray}
	\vv_{k} &=& (1-\beta_k)\vv^{k-1}+\beta_k\widetilde\nabla_k
	\\[5pt]
    {\mW}^{k+1} &=& \argmin_{\mW}\;\inprod{\vv^k}{{\mW}-{\mW}^k}+\frac{1}{\alpha_k}\cD_\phi({\mW},{\mW}^k)+\lambda\|\mW\|_1,
    \label{DNN-sub-ori}
\end{eqnarray}
where $\widetilde\nabla_k$ is mini-batch gradient computed by automatic
differentiation \citep{griewank2008evaluating}. Omitting all the constants, the subproblem takes the form of:
\be 
{\mW}^{k+1}=\argmin_{\mW}\;\phi({\mW})+\inprod{\vp^k}{{\mW}}+\alpha_k\lambda\|{\mW}\|_1,
\label{DNN-sub}
\ee
where $\vp^k=\alpha_k\vv^k-\nabla\phi({\mW}^k)$. Here we adopt the kernel function $\phi({\mW})=\frac{1}{2}\|{\mW}\|^2+\frac{\delta}{r}\|{\mW}\|^{r}$ ($r \ge 2$) for training neural networks, and
then we have an explicit solution for \eqref{DNN-sub} in Proposition \ref{DNN-subpro-solution-prop}. %The derivation is similar to that of \cite[Proposition 5.1]{bolte2018first}, so we omit it here. 
\begin{proposition}
Given $\vp^k \in \R^d$, positive constant $\alpha_k$, $\lambda$, and the kernel function $\phi({\mW})=\frac{1}{2}\|{\mW}\|^2+\frac{\delta}{r}\|{\mW}\|^{r}$ $(r \geq 2,\ \delta>0)$. The solution of the subproblem \eqref{DNN-sub} is given by 
\[
{\mW}^{k+1}=-t^*\vp^+,
\]
where $t^*$ is the unique positive real root of the equation
\begin{eqnarray}
(\delta\|\vp^+\|^{r-2})t^{r-1}+t-1=0,
\label{eq-t*}
\end{eqnarray}
and $\vp^+$ is given by 
\[
\vp^+ = \argmin_\vp\Big\{\frac{1}{2}\|\vp-\vp^k\|^2+\alpha_k\lambda\|\vp\|_1\Big\}
\]
which has an explicit expression given by $\vp_j^+ = \operatorname{sign}(\vp_j^k) \max(|\vp_j^k|-\alpha_k\lambda,0)$ for 
the $j$-th coordinate.
\begin{proof}
The optimality condition of \eqref{DNN-sub} is given by 
\[
0=\mW^{k+1}(1+\delta\|\mW^{k+1}\|^{r-2})+\vp^k+\alpha_k\lambda{\bm\Gamma}^k,\;\text{where }{\bm\Gamma}^k\in\partial\|\cdot\|_1(\mW^{k+1}).
\]
Let $\vp^+=\vp^k+\alpha_k\lambda{\bm\Gamma}^k$. By the optimality condition, we have $\mW^{k+1}=-t\vp^+$ for some positive scalar $t$, and 
\[
(-t-\delta\|\vp^+\|^{r-2}t^{r-1}+1)\vp^+=0.
\]
If $\vp^+\neq0$, then $\delta\|\vp^+\|^{r-2}t^{r-1}+t-1=0$. If $\vp^+=0$, then $\mW^{k+1}=-t\vp^+=0$. Since $t>0$, then we have $\partial\|\cdot\|_1(\mW^{k+1})=\partial\|\cdot\|_1(-t\vp^+)=-\partial\|\cdot\|_1(\vp^+)$. Recall the definition of $\vp^+$, we have 
\[
\vp^+=\vp^k+\alpha_k\lambda{\bm\Gamma}^k\in\vp^k-\alpha_k\lambda\partial\|\cdot\|_1(\vp^+),
\]
which is sufficient and necessary 
optimality condition of the convex optimization problem:
\[
\vp^+ = \argmin_\vp\left\{\frac{1}{2}\|\vp-\vp^k\|^2+\alpha_k\lambda\|\vp\|_1\right\}.
\]
This completes the proof by noting the the above minimization problem is the well-known soft threshold operator, see for example \cite{friedman2010regularization}. 
\end{proof}

\label{DNN-subpro-solution-prop}
\end{proposition}

\begin{example}
In the absence of regularization, that is, when $\lambda=0$, 
then $\bm{p}^+ = \bm{p}^k$ and 
the update formula for MSBPG  at the $k$-th iteration simplifies to ${\mW}^{k+1}=-t^* \vp^k$, where $t^*$ is the positive root of the equation
\eqref{eq-t*}. In this case,
$\mW^{k+1} = t^*(\nabla\phi(\mW^k)-\alpha_k\vv^k)$.
\end{example}
\begin{example}
If we set the $L_1$ regularization parameter $\lambda$ to zero and choose the kernel function simply as the Euclidean distance, i.e. $\delta=0$, then SBPG  reduces to SGD with momentum. Specifically, we have $t^*=1$ and the update
\[
{\mW}^{k+1}={\mW}^k-\alpha_k\vv^k.
\]
\end{example} 

\paragraph{Determining degree of kernel function}
We now turn our attention to selecting the appropriate parameter $r$ for the kernel function. Intuitively, in order to bound the Hessian of the loss function in \eqref{DNN-loss-model}, particularly when the number of layers $L$ in \eqref{DNN-def} is large, $r$ should also be chosen to be larger, so that $\nabla^2F\preceq\frac{1}{\alpha}\nabla^2\phi$ holds globally for some $\alpha>0$. However, in this case, a significant numerical issue may arise when computing $\|{\mW}\|^{r-2}$. This problem can be avoided if the deep neural network exhibits some special structure such that a moderate $r$ can make $F(\mW)$ smooth adaptable with respect to $\phi(\mW)$. For simplicity of analysis, we assume all the given label $y_i$ as zero and consider a sum of squares error loss function. Then, we have a two-layer model defined as follows:
\be 
\min_{\mW=(\vu,\vv)}\;F(\mW)=\frac{1}{2}\left\|\sigma\left(\Mat(\vu)(g(\vv))\right)\right\|^2,
\label{general-two-layer-model}
\ee 
where $\vv\in\R^n$, $\vu\in\R^{km}$, $g:\R^n\rightarrow\R^m$, $\sigma:\R\rightarrow\R$, $\Mat(\vu)\in\R^{k\times m}$ and $\sigma(\cdot)$ is a coordinate-wise operator. Notably, any deep neural network can be recast as the two-layer model given by \eqref{general-two-layer-model}. For instance, if we define $\vv=(\mW_1,...,\mW_{L-1})$, $\vu={\rm Vec}(\mW_L)$, $g(\mW_1,...,\mW_{L-1})=\sigma_{L-1}(\mW_{L-1}(...(\sigma_1( \mW_1\vx))...))$, then model \eqref{DNN-def} can be reformulated as \eqref{general-two-layer-model}. We make the following assumptions in this section, which guarantees that we can find a polynomial kernel function $\phi$ with a moderate degree, such that $F$ in \eqref{general-two-layer-model} is smooth adaptable to $\phi$.
\begin{assumption}\label{act-func}
 %The first and second derivatives of the activation function $\sigma^\prime(\cdot)$ and $\sigma^{\prime\prime}(\cdot)$ are globally bounded.

 $\sigma$ is twice differentiable and $\sigma'$ and $\sigma\cdot\sigma''$ are globally bounded.
\label{bound_assumption_activation}
\end{assumption}
\begin{assumption}
 $g$ is twice differentiable. All partial derivatives of order zero, one, and two of $g$ are globally bounded.
\label{bound_assumption_bn}
\end{assumption}

\begin{remark}
Now we give some remarks on the above assumptions.
\begin{itemize}
\item[1.] Assumption \ref{bound_assumption_activation} is typically valid for various commonly used smooth activation functions.  For example, the sigmoid function $\sigma(x)=\frac{1}{1+e^{-x}}$ satisfies global boundedness for both $\sigma$ and $\sigma''$. Certain activation function may not have bounded function value, such as {GELU} \citep{hendrycks2016gaussian}, which takes the formulation of $\sigma(x)=x\Phi(x)$ where $\Phi$ is the standard Gaussian cumulative distribution function. Nonetheless, the product $\sigma\cdot\sigma''$ is globally bounded. Another type of activation function satisfying Assumption \ref{bound_assumption_activation} is the smoothed ReLU function, for example, the following smoothed ReLU function, which we will consider in our numerical experiments:
\[
\sigma_\epsilon(x) = \left\{
\begin{array}{cc}
     0 & x \le 0 \\
     x^3\left(\frac{1}{\epsilon^2}-\frac{x}{2\epsilon^3}\right) & 0 < x \le \epsilon \\
     x - \frac{\epsilon}{2} & x > \epsilon.
\end{array}\right.
\]
We observe that as $\epsilon$ tends to zero, $\sigma_\epsilon$ converges to the ReLU function. It is straightforward to verify that $\sigma_\epsilon\cdot\sigma_\epsilon''$ is globally bounded.  
Specifically, $\frac{3}{4}$ is a uniform bound on $\sigma_\epsilon\cdot\sigma_\epsilon''$ for $\epsilon\in(0,\frac{1}{2})$. 

\item[2.] In many popular neural network frameworks, batch normalization (BN) layers \citep{ioffe2015batch} are often used before the fully connected layers. For example, in the VGG \citep{simonyan2014very} and ResNet \citep{he2016deep}, BN layers are usually used before the last linear layer. In this case, we can treat all layers except the last one as one layer, which can be modeled as \eqref{general-two-layer-model}. It is expected that the BN layer can make the function $g$ sufficiently smooth, thereby satisfying  Assumption \ref{bound_assumption_bn}.
%For example, for the BN-type function $b(x)=\frac{x-\mu}{\sqrt{(x-\mu)^2+1}}$, we have $|b^\prime(x)|=\mathcal O(1/|x|)$ and $|b^{\prime\prime}(x)|=\mathcal O(1/|x|^2)$. %Another example satisfying Assumption \ref{bound_assumption_bn} is the penultimate activate function is sigmoid function. 
\end{itemize}
\end{remark}

By applying the chain rule, we can compute the Hessian of $F$ and determine a suitable degree parameter $r$ in the kernel function, which will ensure that $\nabla^2F$ is bounded by $\nabla^2\phi$ globally. Consequently, $F$ is smooth adaptable with respect to $\phi$.
%%with not so large degree $r$. %The proof can be found in Appendix \ref{appendix-DNN}. 
In order to compute the Hessian of $F$, two formulas are required, which can be verified directly.\vspace{-0.1cm}
\begin{lemma}
Let $\vu\in\R^{km}$, $\vg\in\R^m$, $\mA \in\R^{n\times m}$ and $\vb\in\R^k$. Consider two linear maps: $\vu\mapsto\Mat(\vu)\vg$ and $\vu\mapsto \mA(\Mat(\vu))^T\vb$, then, the Jacobian of the two maps are given by
\[
J_{\vu}\left[\Mat(\vu)\vg\right]=\vg^T\otimes\mathbb{I}_k,
\]
\[
J_{\vu}[\mA(\Mat(\vu))^T\vb]=\mA \otimes \vb^T.
\]
\label{Jac-Hess-aux-lemma}
\end{lemma}\vspace{-0.5cm}

\begin{proposition}
Suppose Assumptions \ref{bound_assumption_activation} and \ref{bound_assumption_bn} hold. Then, for any given $\delta>0$ and any $r\geq4$, the function $F$ defined in \eqref{general-two-layer-model} is smooth adaptable with respect to $\phi(\mW)=\frac{1}{2}\|\mW\|^2+\frac{\delta}{r}\|\mW\|^r$.
\label{DNN-kernel-degree-prop}
\end{proposition}
\begin{proof}
We denote $\Mat(\vu)$ by $\mM$. The Jacobian of $g$ is denoted by $Jg$, while its transpose is denoted by $J^Tg$. $\mathbb{I}_k$ is $k\times k$ identity matrix. Using Lemma \ref{Jac-Hess-aux-lemma}, we can compute the Jacobian and Hessian of $F$ as follows:

{\noindent\bf Jacobian of $F$ :}\\
\be 
\begin{aligned}
\frac{\partial F}{\partial\vu}&=\left(g(\vv)\otimes\mathbb I_k\right)\left[\sigma^\prime(\mM g(\vv))\circ\sigma(\mM g(\vv))\right],\\
\frac{\partial F}{\partial\vv}&=J^Tg(\vv)\mM^T\left[\sigma^\prime(\mM g(\vv))\circ\sigma(\mM g(\vv))\right].
\end{aligned}
\label{Jacobian_F}
\ee 

{\noindent\bf Hessian of $F$ :}\\
\be 
\begin{aligned}
&\frac{\partial^2F}{\partial\vu^2}=(1)+(2),\\
\text{where }(1)&=\left(g(\vv)\otimes\mathbb I_k\right)\text{Diag}\big(\sigma(\mM g(\vv))\circ\sigma^{\prime\prime}(\mM g(\vv))\big)(\vg^T(v)\otimes\mathbb I_k)\\
(2)&=\left(g(\vv)\otimes\mathbb I_k\right)\text{Diag}\big(\sigma^\prime(\mM g(\vv))\circ\sigma^{\prime}(\mM g(\vv))\big)(\vg^T(v)\otimes\mathbb I_k).
\end{aligned}
\label{Hessian_Fuu}
\ee 

\be 
\begin{aligned}
&\frac{\partial^2F}{\partial\vu\partial\vv}=(1)+(2)+(3),\\
\text{where }(1)&=\left(J^Tg(\vv)\right)\otimes\left[\sigma^\prime(\mM g(\vv))\circ\sigma^{\prime}(\mM g(\vv))\right]^T\\
(2)&=J^Tg(\vv)\mM^T\text{Diag}[\sigma(\mM g(\vv))\circ\sigma^{\prime\prime}(\mM g(\vv))]\left(g^T(\vv)\otimes\mathbb I_k\right)\\
(3)&=J^Tg(\vv)\mM^T\text{Diag}[\sigma^\prime(\mM g(\vv))\circ\sigma^{\prime}(\mM g(\vv))]\left(g^T(\vv)\otimes\mathbb I_k\right).
\end{aligned}
\label{Hessian_F_uv}
\ee 
\be 
\begin{aligned}
&\frac{\partial^2F}{\partial\vv^2}=(1)+(2)+(3),\\
\text{where }(1)&=D^2g(\vv)\left[\mM^T[\sigma^\prime(\mM g(\vv))\circ\sigma(\mM g(\vv))]\right]=\sum d_i\nabla^2g_i(\vv)\\
(2)&=J^Tg(\vv)\mM^T\text{Diag}[\sigma(\mM g(\vv))\circ\sigma^{\prime\prime}(\mM g(\vv))]\mM Jg(\vv)\\
(3)&=J^Tg(\vv)\mM^T\text{Diag}[\sigma^\prime(\mM g(\vv))\circ\sigma^{\prime}(\mM g(\vv))]\mM Jg(\vv),
\end{aligned}
\label{Hessian_F_vv}
\ee 
where $\vd=\mM^T[\sigma^\prime(\mM g(\vv))\circ\sigma(\mM g(\vv))]$. Now, we are ready to prove this proposition. For any $\vw\in\R^{km+n}$ and $\vh=[\vh^u;\vh^v]\in\R^{km+n}$, it suffices to prove that $\inprod{\nabla^2F(\vw)\vh}{\vh}=\mathcal O(\inprod{\nabla^2\phi(\vw)\vh}{\vh})$. From \eqref{Hessian_Fuu}\eqref{Hessian_F_uv}\eqref{Hessian_F_vv} and Assumption \ref{bound_assumption_activation}, \ref{bound_assumption_bn}, we can easily get $\inprod{\nabla^2F(\vw)\vh}{\vh}=\mathcal O((1+\|\vw\|^2)\|\vh\|^2)$. On the other hand, $\nabla^2\phi(\vw)=I(1+\|\vw\|^{r-2})+(r-2)\|\vw\|^{r-4}\vw\vw^T$. Hence $\inprod{\nabla^2\phi(\vw)\vh}{\vh}\geq(1+\|\vw\|^{r-2})\|\vh\|^2$. So, we only require $r-2\geq2$. This completes the proof.   
\end{proof}


\paragraph{Layerwise kernel function}
In Proposition \ref{DNN-subpro-solution-prop}, the kernel function $\phi({\mW})=\frac{1}{2}\|{\mW}\|^2+\frac{\delta}{r}\|{\mW}\|^{r}$ is used, which means we adopt the same Bregman distance for all layers of deep neural networks. However, different layers have different optimization property for deep neural networks \citep{you2019large}, and computing $\|\mW\|^r$ with $r>2$ may result in numerical issues for neural networks with millions of parameters, such as in VGG \citep{simonyan2014very}.
% would need to compute the norm of the whole weight $\mW=(\mW_1,\mW_2,\cdots,\mW_L)$ at the beginning of each iteration, which is computational expensive. 
To take advantage of the layerwise structure of neural networks, we design a layerwise kernel function for a $L$-layer neural network as follows:
\be
\phi({\mW})=\sum_{i=1}^L\phi_i({\mW}_i),\quad \phi_i({\mW}_i)=\frac{1}{2}\|{\mW}_i\|^2+\frac{\delta}{r}\|{\mW}_i\|^{r}.
\label{kernel-nn}
\ee
Note that $\delta$ and $r$ can vary from layer to layer, here we take the same $\delta$ and $r$ for different layers for simplicity.
Then, we have the Bregman distance taking the form $\cD_\phi=\sum_{i=1}^L\cD_{\phi_i}$. By employing this Bregman distance in subproblem \eqref{DNN-sub-ori}, our MSBPG  algorithm can be implemented in a layerwise manner. See the details in Algorithm \ref{dnn-alg}.

\begin{algorithm}[ht] 
    \caption{Momentum based Stochastic Bregman Proximal Gradient (MSBPG) for training neural networks}
    \begin{algorithmic}[1]
        \State \textbf{Input: }Total number of iterations $K$, stepsize $\alpha_k$, momentum parameter $\beta_k$, $\delta$ and $r$ to determine the kernel function $\phi$.
        \State \textbf{Initialize: }Set $\mW=\mW^0$, $\vv_0=0$.
        \For {$k=0, \cdots, K-1$}
        \State Compute mini-batch gradient $\widetilde\nabla_k$;
        \State Compute SMAE: $\vv^k=(1-\beta_k)\vv^{k-1}+\beta_k\widetilde\nabla_k$;
        \For {$i=1,\ldots,L$}
        \State $\vp_i^k = \alpha_k \vv_i^k - \nabla\phi(\mW^{k}_i)$;
        \State $\vp_i^+ = \argmin_{\vp_i}\{\frac{1}{2}\|\vp_i-\vp_i^k\|^2+\alpha_k\lambda\|\vp_i\|_1\}$;
        \State Solve $(\delta \|\vp^+_i\|^{r-2})t_i^{r-1} + t_i -1 =0$ to get $t_i^k$;
        \State $\mW_i^{k+1} = - t_i^k \vp_i^+$;
        \EndFor
        \EndFor
        \State {\bf Output: }$\mW^1, \cdots, \mW^K$.
    \end{algorithmic}\label{dnn-alg}
\end{algorithm}

% At the $k$-th step, for each layer $i=1,\ldots,L$:

% \noindent{\bf Step 1: } Compute stochastic gradient estimator $\vg_i^k$, $\nabla\phi_i({\mW}_i^k)={\mW}_i^k(1+\delta_i\|{\mW}_i^k\|^{r_i-2})$, and $\vp_i^k=\alpha_k \vg_i^k-\nabla\phi_i({\mW}_i^k)$.

% \noindent{\bf Step 2: }Compute the unique positive root $t^k_i$ of the univariate equation: $$(\delta_i\|\vp_i^k\|^{r-2})t^{r-1}+t-1=0.$$

% \noindent{\bf Step 3: } ${\mW}_i^{k+1}=-t^k_i\vp_i^k$.

\paragraph{Mitigating gradient explosion}
In the training of deep neural networks, gradient explosion is a common undesired phenomenon, where the gradients of the loss function grow exponentially from layer to layer, leading to numerical instability or even collapse of the training process \citep{hochreiter1991untersuchungen,manchev2020target}. The reasons for gradient explosion include selecting a large stepsize and choosing an improper initialization for the model's parameters \citep{pascanu2013difficulty}. In the following, we will show that MSBPG  provides a novel approach to mitigate gradient explosion. Considering MSBPG  without regularization, the update rule is given by:
\be
{\mW}_i^{k+1}=-t_i^k\vp_i^k=t_i^k\left((1+\delta\|{\mW_i}^k\|^{r-2}){\mW_i}^k-\alpha_k\vv_i^k\right),
\label{Bregman-GD}
\ee
where $t_i^k\in(0,1)$ is the unique positive root of 
\be 
\left(\delta\left\|(1+\delta\|{\mW}_i^k\|^{r-2}){\mW}_i^k-\alpha_k\vv_i^k\right\|^{r-2}\right)t^{r-1}+t-1=0.
\label{Bregman_sub_equation}
\ee
Combining \eqref{Bregman-GD} and \eqref{Bregman_sub_equation}, we have the following equivalent implicit update scheme for the $i$-th layer:
\be
\mW_i^{k+1}=\frac{1+\delta\|\mW_i^k\|^{r-2}}{1+\delta\|\mW_i^{k+1}\|^{r-2}}\mW_i^k-\frac{\alpha_k}{1+\delta\|\mW_i^{k+1}\|^{r-2}}\vv_i^k.
\label{implicit-sbpg}
\ee
%Denote the coefficient $\delta\left\|(1+\delta\|{\mW}^k\|^{r-2}){\mW}^k-\alpha_k\vg^k\right\|^{r-2}$ by $M$.  We make the following observations: if $M\rightarrow\infty$, then $t^k\rightarrow0+$. On the other hand, if $t^k(1+\delta\|{\mW}^k\|^{r-2})<1$, then SBPG can also be regarded as a weight decay framework. This condition tends to hold at the early stages of the optimization process. Without loss of generality, we assume $\delta=1$ and consider the following two cases:
%\begin{itemize}
%    \item[1.] Suppose that $\|\vg^k\|\gg\max\{\|{\mW}^k\|,1\}$. Then, we have $t^k\approx0$. Consequently, $t^k(1+\delta\|{\mW}^k\|^{r-2})<1$.  This implies that the weight decay effect dominates in the early stage of SBPG.
%    \item[2.] When $\vg^k\approx0$, it is likely that ${\mW}^k$ is approaching a critical point. Let $\beta=\|{\mW}^k\|$, and let $t^k$ be approximately the unique root of $(1+\beta^{r-2})^{r-2}\beta^{r-2}t^{r-1}+t-1=0$ in the interval $[0,1]$, denoted by $t^*$. It can be easily verified that $t^*=\frac{1}{1+\beta^{r-2}}$. Therefore, $t^k(1+\delta\|{\mW}^k\|^{r-2})\approx1$.
%\end{itemize}

% If the norm of the stochastic gradient $\|\vg^k\|$ is extremely large, the update parameter $t^k$ in Equation \eqref{Bregman_sub_equation} can be approximated as 
% \[
% t^k=\left(\frac{1-t^k}{\delta\|\vp^k\|^{r-2}}\right)^{\frac{1}{r-1}}\approx\frac{1}{\delta^{\frac{1}{r-1}}\|\vp^k\|^{\frac{r-2}{r-1}}},
% \]
% Using this approximation, we can then approximate the next weight ${\mW}^{k+1}$ as
% \[
% {\mW}^{k+1}=-t^k\vp^k\approx\frac{\vp^k}{\delta^{\frac{1}{r-1}}\|\vp^k\|^{\frac{r-2}{r-1}}}=\frac{\vp^k}{\|\vp^k\|}\left(\frac{\|\vp^k\|}{\delta}\right)^{\frac{1}{r-1}}.
% \]
% In situations where the gradient $\vg^k$ has a large magnitude, denoted by $\|\vg^k\|$, the magnitude of $\vp^k$, denoted by $\|\vp^k\|$, is expected to be approximately proportional to $\|\vg^k\|$, i.e. $\|\vp^k\|={\Theta}(\|\vg^k\|)$. As a result, $\|{\mW}^{k+1}\|$ would be approximately proportional to $\|\vg^k\|^{\frac{1}{r-1}}$, i.e. $\|{\mW}^{k+1}\|={\Theta}(\|\vg^k\|^{\frac{1}{r-1}})$.
% \red{
% For the coefficient of the weight term $\frac{1+\delta\|\mW_i^k\|^{r-2}}{1+\delta\|\mW_i^{k+1}\|^{r-2}}$, since the training of neural network takes a small stepsize, the change of the norm of the weight $\|\mW_i\|$ is marginal at each iteration, making this coefficient a value approximately 1. For the coefficient of the gradient $\frac{\alpha_k}{1+\delta\|\mW_i^{k+1}\|^{r-2}}$, we can observe that it act as a scaling for the gradient. 
% With large stepsize or large initial point, gradient tends to explode while
% the norm of the weight $\|\mW_i\|$ tends to be large at the same time. By scaling the gradient with $\frac{\alpha_k}{1+\delta\|\mW_i^{k+1}\|^{r-2}}$, the weight is relieved from moving excessively in the direction of the gradient, which can prevent the collapse of the training process.}
% \fbox{Replace the above by the following?}
{It is observed in practice that with large stepsize or large initial point, the gradient $\vv_i^k$ tends to explode if no scaling is done, while
the norm of the weight $\|\mW_i^{k+1}\|$ also tends to be large. In \eqref{implicit-sbpg}, we see that by scaling the gradient with $\frac{1}{1+\delta\|\mW_i^{k+1}\|^{r-2}}$, the weight $\mW_i^{k+1}$ is relieved from moving excessively in the direction of the gradient to avoid 
rapid growth of its norm. At the same time, we can also see that
if the norm  $\|\mW_i^{k+1}\|$ does not change drastically, 
the coefficient of $\mW_i^k$ in \eqref{implicit-sbpg} will 
be maintained to be approximately 1. 
Thus in \eqref{implicit-sbpg}, 
we see an automatic scaling of the gradient to avoid rapid growth
of the weight and hence also mitigating subsequent gradient explosion.}

Experimental results in Section \ref{DNN-experiment} indeed verify MSBPG 's ability to mitigate gradient explosion for training deep neural networks. An intuitive illustration of MSBPG 's ``pull-back" ability is given in Figure \ref{additional-figure} in Appendix \ref{app-fig}, and this ``pull-back" ability originates from the Bregman proximity model and the polynomial kernel function we adopt.
% This is achieved by taking a scaled gradient, which can be interpreted as a "pull-back" operation in the context of Bregman gradient descent mapping. In other words, the mapping constrains the movement of the point with a large gradient, relieving it from moving too far in that direction. Increasing the value of $r$ generally improves the effectiveness of this explosion relieveion mechanism. This intuition is illustrated in , which can be found in the Appendix.

\paragraph{Improving generalization capacity}
From \eqref{implicit-sbpg} we can see that MSBPG  employs a scaling for the gradient during the update of the parameters as a result of adopting a Bregman proximity model and a polynomial kernel function. 
This scaling  not only helps to mitigate the gradient explosion phenomenon, but also can improve the generalization capacity of MSBPG . To be specific, at the beginning of the training process, the initial weight of each layer $\mW_i$ tends to have a larger norm $\|\mW_i\|$ and therefore MSBPG  takes a cautious update at the beginning phase of training. As the training goes on, due to the effect of regularization, either $L_1$ or $L_2$ regularization, the norm of each layer's weight $\|\mW_i\|$ becomes smaller, and MSBPG  can then take bolder update of the parameters. This implicit training strategy of MSBPG  is in agreement with the idea of a heuristic deep learning training technique called ``learning rate warm-up" \citep{gotmare2018closer}, which benefits the training stability and generalization performance \citep{liu2019variance}. Experimental results in Section \ref{DNN-experiment} also testify the excellent generalization capacity of MSBPG .


\section{Numerical experiments}\label{section-experiments}
In this section, we conduct numerical experiments to showcase the effectiveness and robustness of  MSBPG in comparison to modern solvers commonly used in deep learning. We assess the impact of stepsize and initial point selection on the performance of our method. Our experiments consist of two parts. In the first part, we use a quadratic inverse problem as a toy example to illustrate the capabilities of vanilla SBPG. The second part is the main focus of this section, where we evaluate the performance of MSBPG in training deep neural networks. The experiments for the quadratic inverse problem are conducted using MATLAB R2021b on a Windows workstation equipped with a 12-core Intel Xeon E5-2680 @ 2.50GHz processor and 128GB of RAM. For the deep learning experiments, we conducted the experiments using PyTorch running on a single RTX3090 GPU.

\subsection{Quadratic inverse problem}
The quadratic inverse problem, as formulated in \cite{bolte2018first}, is given by:
\[
\min\left\{\Phi(\vx):=\underbrace{\frac{1}{4}\sum_{i=1}^n(\inprod{A_i\vx}{\vx}-b_i)^2}_{F(\vx)}+\lambda R(\vx): \vx\in\R^d\right\},
\]
which has {practical applications \citep{beck2013sparsity} that} includes the phase retrieval problem as a special case \citep{luke2017phase}. In this experiment, we consider the $L_1$ regularization $R(\vx)=\|\vx\|_1$ with $\lambda=1\times10^{-3}$, and solve the quadratic inverse problem using SBPG and stochastic {(Euclidean)} proximal gradient (SPG) method \citep{bertsekas2011incremental}. Notably, SPG is a special case of SBPG, in which $\phi(\vx)=\frac{1}{2}\|\vx\|^2$. Since the smooth term in the objective function $F(\vx)$ does not admit a globally Lipschitz continuous gradient, we employ the kernel function $\phi(\vx)=\frac{1}{2}\|\vx\|^2+\frac{1}{r}\|\vx\|^r$ with $r=4$. It has been shown in \cite{lu2018relatively} that any $r\geq4$ guarantees that $F$ is $\phi$-smooth adaptable globally. Moreover, according to \cite{bolte2018first}, the smooth adaptable constant $L_F$ can be chosen such that $L_F\geq\sum_{i=1}^n(3\|A_i\|^2+\|A_i\||b_i|)$ for $r=4$. In this experiment, we randomly generate the data by the following ${\tt MATLAB}$ commands: 
\[
\begin{aligned}
&{\tt ai = randn(d,1);Ai = ai*ai';}\\
&{\tt x\_true = sprandn(d,1,density\_x); b\_i = x\_true'*(Ai*x\_true);}
\end{aligned}
\]

The {true} solution for the quadratic inverse problem is chosen as a sparse vector $\vx^*$ that satisfies $\inprod{A_i\vx^*}{\vx^*}=b_i$ for $i=1,\dots,n$. We set the mini-batch size for all algorithms to be $m=1$. To evaluate the effectiveness of each algorithm, we use the following criterion that takes into account the possibility of critical points being local minimum or saddle points:
\[
\epsilon_k=\max\left\{\|\cG({\vx^k})\|,\;\epsilon_{\tt obj}:=\frac{{\tt obj}_k-{\tt obj}_*}{1+{\tt obj}_*}\right\},
\]
 where ${\tt obj}_k=\Phi({\vx^k})$ and ${\tt obj}_*=\Phi(\vx^*)$. 
 {The term $\|\cG(\vx^k)\|$ measures the stationarity of the solution, while a small $\eps_{\tt obj}$ indicates that the solution is a "nearly" global minimum.}

We conduct experiments on a problem with data size $d=100$ and ${\tt density\_x=0.05}$. All methods are run until they reach an accuracy of $\epsilon_k\leq0.01$ within a time limit of 30 seconds. To ensure statistical significance, we run each algorithm 10 times and report the median value. The results are presented in Figure \ref{robust-figure}. For Figures \ref{robust-figure}(a), we randomly select initial points within a ball centered at the origin with radius $1\times10^{-2}$. We use the stepsize schedule of $\alpha_k=\max\left\{10^{-4},\frac{\alpha_0}{\sqrt{1+k}}\right\}$, where $\alpha_0$ is the initial stepsize. For Figure \ref{robust-figure}(b), we set constant stepsize schedule $1\times10^{-3}$. 
{For Figures \ref{robust-figure}(c), we randomly select initial points within a ball centered at the origin with radius $1\times10^{-2}$. We use constant stepsize schedule. To prevent excessively small stepsizes that can slow down all methods, we set a lower bound for the stepsize.

Figure \ref{robust-figure}(a) demonstrates that SBPG has a larger range of convergent stepsizes than SPG, indicating that SBPG is more robust in terms of stepsize selection. The impact of the initial stepsize on the performance of the algorithms is reported in this figure. Figure \ref{robust-figure}(b) shows that SBPG is much more robust than SPG in terms of initial point selection. Specifically, SBPG exhibits high resilience to initial point selection and avoids causing the training to collapse. Figure \ref{robust-figure}(c) reveals that a larger degree $r$ in the kernel function increases the safe stepsize threshold. These observations are partly explained in Section \ref{DNN implementation}. Since a large stepsize and a large radius of the initial point tend to lead to gradient explosion, Bregman proximal mapping helps to pull back the iterate and guide it towards a better solution.

%We compare five methods for solving sparse quadratic inverse problems of varying sizes, as shown in Table \ref{QIP-table}. The initial points for all methods are randomly selected within a ball centered at the origin with radius $1\times10^{-2}$. For each example, we run each algorithm 10 times and report the median value. To generate $A_i$, we use the ${\tt MATLAB}$ commands ${\tt ai=sprandn(d,n,0.1);Ai=ai*ai';}$ for problems with $d\geq1000$, and ${\tt ai=randn(d,n);Ai=ai*ai';}$ for those with $d\leq500$. The algorithm stops when $\epsilon_k<10^{-4}$ or the CPU time exceeds 360 seconds. We observe that SVRGBPG and SVRGBPG-bb perform the best over all data sets. The stepsizes for BPG, BPGe, SBPG, and SVRGBPG are selected from the grid: $\{10^{-5},5\times10^{-5},10^{-4},5\times10^{-4},10^{-3},5\times10^{-3},10^{-2},5\times10^{-2},10^{-1},5\times10^{-1},10^{0}\}$ and the best one is chosen. The initial stepsize for SVRGBPG-bb in all tests is $1\times10^{-3}$. Compared to SVRGBPG, we can see that the BB strategy can automatically select a "best" stepsize, thereby saving time in tuning the stepsize. However, for some problems, such as $d=2000,n=5000$, all methods fail to obtain a desirable accurate solution within 360 seconds, as they stagnate at a spurious local minimum.


\begin{figure}
\centering
\subfigure[Stepsize robustness]{
\includegraphics[width=4.5cm]{robust_step-eps-converted-to.pdf}
%\caption{fig1}
}
%\quad
\subfigure[Initial point robustness]{
\includegraphics[width=4.5cm]{robust_initial-eps-converted-to.pdf}
}
\subfigure[Safe stepsize threshold]{
\includegraphics[width=4.5cm]{step_threshold-eps-converted-to.pdf}
}
\caption{Comparison of SBPG and SPG in terms of their robustness with respect to stepsize and initial point selction. A method is considered non-convergent if it fails to reach an accuracy of $\epsilon_k<10^{-2}$ within 30 seconds or if it collapses. Generally, choosing large stepsize and large radius for the initial point can cause an algorithm to collapse. The safe stepsize threshold is the maximum stepsize (constant schedule) that a method does not collapse. We run 10 tests for each algorithm and report the median of the results.}
\label{robust-figure}
\end{figure}

%\input{table}


\subsection{Deep neural network} \label{DNN-experiment}

For the evaluation of MSBPG's performance on training deep neural networks, we consider the model with $L_2$ regularization here for its better generalization capacity:
\be 
\min_{\mW}\;\underbrace{\frac{1}{N}\sum_{i=1}^N\cL(\mathcal{DNN}(\mW,\vx_i),y_i)}_{F(\mW)}+\lambda\|\mW\|_2^2.
\label{exp-dnn-model}
\ee 
We employ MSBPG to solve this large-scale problem. Following AdamW \citep{loshchilov2017decoupled}, we also conduct decoupled weight decay at the end of each iteration for $L_2$ regularization and do not consider the $L_2$ regularization term when solving the subproblems. The detailed algorithm is summarized in Algorithm \ref{dl-alg-sbpg}. At iteration $k$, MSBPG  first uses automatic differentiation to compute the mini-batch gradient $\widetilde{\nabla}_k$. Then, it maintains a bias-corrected gradient estimator $\bar{\vv}^k$ \citep{kingma2014adam} and use it to calculate the layerwise $\vp_i^k$. With $\vp_i^k$, MSBPG  solves a univariate equation to get $t_i^k$ and update the weight of the $i$-th layer to $\mW_i^k$. In the end, MSBPG  conducts decoupled weight decay as 
$L_2$ regularization.

\begin{algorithm}[ht] 
    \caption{MSBPG with $L_2$ regularization}
    \begin{algorithmic}[1]
        \State \textbf{Input: }Total number of training epochs $K$, momentum coefficient $\beta$, stepsize $\alpha_k$, weight decay coefficient $\gamma$, $\delta$ and $r$ to determine the kernel function $\phi$.
        \State \textbf{Initialize: } Set $\mW=\mW^0$, $\vv^0=\bf{0}$.
        \For {$k=1, \cdots, K$}
        \State Compute mini-batch gradient $\widetilde{\nabla}^k$;
        \State $\vv^k = \beta\vv^{k-1} + (1-\beta) \widetilde{\nabla}^k$, $\bar{\vv}^k = \vv^k / (1 - \beta^k)$;
       \For {$i=1, \cdots, L$} 
        \State $\vp_i^k = \alpha_k \bar{\vv}_i^k - \nabla\phi(\mW^{k-1}_i)$;
        \State Solve $(\delta \|\vp^k_i\|^{r-2})t_i^{r-1} + t_i -1 =0$ to get $t_i^k$;
        \State $\widetilde{\mW}_i^{k} = - t_i^k \vp_i^k$;
        \EndFor
        \State $\mW^k = \widetilde{\mW}^k - \alpha_k \gamma \mW^{k-1}$;
        \EndFor
        \State {\bf Output: $\mW^1, \cdots, \mW^K$}
    \end{algorithmic}\label{dl-alg-sbpg}
\end{algorithm}


We conducted extensive experiments on several representative benchmarks, including VGG16~\citep{simonyan2014very}, ResNet34~\citep{he2016deep} on CIFAR10 dataset \citep{krizhevsky2009learning}, ResNet34~\citep{he2016deep}, DenseNet121~\citep{huang2017densely} on CIFAR100 dataset \citep{krizhevsky2009learning},
and LSTMs~\citep{schmidhuber1997long} on the Penn Treebank dataset~\citep{marcinkiewicz1994building}. We compare MSBPG  with the most popular optimization algorithms used for training neural networks, including SGD \citep{sutskever2013importance}, Adam \citep{kingma2014adam}, and AdamW \citep{loshchilov2017decoupled}. Experimental results show that MSBPG  has excellent convergence performance and best generalization capacity for both the task that SGD dominates (image classification with CNNs) and the task that Adam dominates (language modeling with LSTMs). We also conducted experiments to compare MSBPG  with SGD on different initial stepsizes and different scales of initial point. Our experimental results demonstrate the robustness of MSBPG  in training neural networks.

Before getting into the details of our experiments, we first make a clarification about the activation function. The frequently used activation function ReLU in VGG, ResNet, and DenseNet takes the form of $\operatorname{ReLU}(x) = \max(0, x)$, which is not continuously differentiable. Here we design a smoothing approximation of ReLU with coefficient $\epsilon$, which is twice continuously differentiable and satisfies our Assumption \ref{act-func}, namely,
\[
\sigma_\epsilon(x) = \left\{
\begin{array}{cc}
     0 & x \le 0 \\
     x^3\left(\frac{1}{\epsilon^2}-\frac{x}{2\epsilon^3}\right) & 0 < x \le \epsilon \\
     x - \frac{\epsilon}{2} & x > \epsilon.
\end{array}\right.
\]
This activation function has gradient taking the form:
\[
\sigma_\epsilon'(x) = \left\{
\begin{array}{cc}
     0 & x \le 0 \\
     x^2(\frac{3}{\epsilon^2}-\frac{2x}{\epsilon^3}) & 0 < x \le \epsilon \\
     1 & x > \epsilon.
\end{array}\right.
\]
Note that as $\epsilon$ tends to 0, this twice continuously differentiable activation function tends to the ReLU function. We conducted experiments with VGG16 on the CIFAR10 dataset, where we replace all the activation functions in VGG16 by $\sigma_\epsilon$ defined above. As shown in Figure \ref{act-func}, our algorithm MSBPG 's performance does not degrade as $\epsilon$ tends to 0. Therefore, in the subsequent experiments, we use the original neural network architectures with the ReLU activation function to evaluate our method MSBPG .

\begin{figure*}[th]
		%	\setcounter{figure}{3}
		%\setlength{\belowcaptionskip}{-2pt}
		\begin{center}
			\setlength{\tabcolsep}{0.0pt}  
			\scalebox{1}{\begin{tabular}{cccc}
					\includegraphics[width=0.25\linewidth]{Train_qelu_result_step.pdf}&
					\includegraphics[width=0.25\linewidth]{Test_qelu_result_step.pdf}&
					\includegraphics[width=0.25\linewidth]{Train_qelu_result_cos.pdf} &
					\includegraphics[width=0.25\linewidth]{Test_qelu_result_cos.pdf}\\
					\multicolumn{2}{c}{\footnotesize{(a) Step decay.}} &  \multicolumn{2}{c}{\footnotesize{(b) Cosine annealing.}}\\
			\end{tabular}}
		\end{center}
	\vspace{-0.8em}
		\caption{Training and test accuracy (\%) of VGG16 on CIFAR10 dataset under two frequently used training settings. Here the activation function of VGG16 adopts smoothed ReLU activation function $\sigma_\epsilon$ with different choices of $\epsilon$ ($\epsilon=0$ denotes adopting the original ReLU activation function).} \label{fig:act-func}
\end{figure*}

\paragraph{CNNs on image classification} We experimented with VGG16, ResNet34 on the CIFAR10 dataset, and ResNet34, DenseNet121 on CIFAR100 dataset. SGD usually has better generalization performance than adaptive gradient algorithms such as Adam and AdamW when training CNNs on image classification tasks and therefore is the default optimizer in these scenarios \citep{he2016deep,zhou2020towards}. We used two dominant experimental settings for training neural networks, including reducing the stepsize to 0.1 times its original value near the end of training \citep{zhuang2020adabelief,chen2021closing,luo2019adaptive} and adopting a cosine annealing schedule for the stepsizes \citep{loshchilov2016sgdr,loshchilov2017decoupled}. The purpose of these two training strategies is to accelerate the convergence of the optimization algorithms so as to give a fair comparison of their generalization capacity. We use the default training hyperparameters of SGD, Adam, and AdamW in these settings \citep{he2016deep,zhuang2020adabelief,chen2021closing}, and set MSBPG 's learning rate (initial stepsize) as 0.1, momentum coefficient $\beta$ as 0.9, weight decay coefficient $\gamma$ as $1\times 10^{-3}$. For the layerwise kernel function $\phi_i({\mW}_i)=\frac{1}{2}\|{\mW}_i\|^2+\frac{\delta}{r}\|{\mW}_i\|^{r}$, we set $r=4,\ \delta=1\times 10^{-2}$ for VGG16 and $r=6,\ \delta=1 \times 10^{-3}$ for ResNet34 on CIFAR10 dataset, and $r=4,\ \delta=1\times 10^{-2}$ for ResNet34 and $r=4,\ \delta=1 \times 10^{-3}$ for DenseNet121 on CIFAR100 dataset. From the experimental results in Figure \ref{fig:cnn1}, \ref{fig:cnn2}, \ref{fig:cnn3}, \ref{fig:cnn4}, we can see that MSBPG  attains 100\% training accuracy in all the training settings, unlike Adam which fails to fully converge with the training strategy of reducing the learning rate to 0.1 times of original value at the 150th epoch. Furthermore, MSBPG  consistently achieves the best generalization performance for all experimental settings, and attains at least 0.5\% test accuracy improvement compared with the second best optimization algorithm. This generalization advantage of MSBPG  
can be attributed to the Bregman proximity model we adopt.

\begin{figure*}[th]
		%	\setcounter{figure}{3}
		%\setlength{\belowcaptionskip}{-2pt}
		\begin{center}
			\setlength{\tabcolsep}{0.0pt}  
			\scalebox{1}{\begin{tabular}{cccc}
					\includegraphics[width=0.25\linewidth]{VGG-BN_Train_result_step.pdf}&
					\includegraphics[width=0.25\linewidth]{VGG-BN_Test_result_step.pdf}&
					\includegraphics[width=0.25\linewidth]{ResNet_Train_result_step.pdf} &
					\includegraphics[width=0.25\linewidth]{ResNet_Test_result_step.pdf}\\
						\multicolumn{2}{c}{\footnotesize{(a) VGG16 on CIFAR10}} &  \multicolumn{2}{c}{\footnotesize{(b) ResNet34 on CIFAR10}}\\
			\end{tabular}}
		\end{center}
	\vspace{-0.8em}
		\caption{Training and test accuracy (\%) of CNNs on CIFAR10 dataset with learning rate reduced to 0.1 times of the original value at the 150th epoch.} \label{fig:cnn1}
	\end{figure*}

 
 \begin{figure*}[th]
		%	\setcounter{figure}{3}
		%\setlength{\belowcaptionskip}{-2pt}
		\begin{center}
			\setlength{\tabcolsep}{0.0pt}  
			\scalebox{1}{\begin{tabular}{cccc}
					\includegraphics[width=0.25\linewidth]{new100_ResNet_Train_step.pdf}&
					\includegraphics[width=0.25\linewidth]{new100_ResNet_Test_step.pdf}&
					\includegraphics[width=0.25\linewidth]{new100_DenseNet_Train_step.pdf} &
					\includegraphics[width=0.25\linewidth]{new100_DenseNet_Test_step.pdf}\\
						\multicolumn{2}{c}{\footnotesize{(a) ResNet34 on CIFAR100}} &  \multicolumn{2}{c}{\footnotesize{(b) DenseNet121 on CIFAR100}}\\
			\end{tabular}}
		\end{center}
	\vspace{-0.8em}
		\caption{Training and test accuracy (\%) of CNNs on CIFAR100 dataset with learning rate reduced to 0.1 times of the original value at the 150th epoch.} \label{fig:cnn3}
	\end{figure*}



 \begin{figure*}[th]
		%	\setcounter{figure}{3}
		%\setlength{\belowcaptionskip}{-2pt}
		\begin{center}
			\setlength{\tabcolsep}{0.0pt}  
			\scalebox{1}{\begin{tabular}{cccc}
					\includegraphics[width=0.25\linewidth]{VGG-BN_Train_result_cos.pdf}&
					\includegraphics[width=0.25\linewidth]{VGG-BN_Test_result_cos.pdf}&
					\includegraphics[width=0.25\linewidth]{ResNet_Train_result_cos.pdf} &
					\includegraphics[width=0.25\linewidth]{ResNet_Test_result_cos.pdf}\\
						\multicolumn{2}{c}{\footnotesize{(a) VGG16 on CIFAR10}} &  \multicolumn{2}{c}{\footnotesize{(b) ResNet34 on CIFAR10}}\\
					%\small (a) VGG16 on CIFAR10. & \small (b) ResNet34 on CIFAR10. \\
					%\small (a) VGG16 on CIFAR10. & \small (b) ResNet34 on CIFAR10. \\
			\end{tabular}}
		\end{center}
	\vspace{-0.8em}
		\caption{Training and test accuracy (\%) of CNNs on CIFAR10 dataset with learning rate using the cosine annealing schedule.} \label{fig:cnn2}
	\end{figure*}

 \begin{figure*}[th]
		%	\setcounter{figure}{3}
		%\setlength{\belowcaptionskip}{-2pt}
		\begin{center}
			\setlength{\tabcolsep}{0.0pt}  
			\scalebox{1}{\begin{tabular}{cccc}
					\includegraphics[width=0.25\linewidth]{new100_ResNet_Train_cos.pdf}&
					\includegraphics[width=0.25\linewidth]{new100_ResNet_Test_cos.pdf}&
					\includegraphics[width=0.25\linewidth]{new100_DenseNet_Train_cos.pdf} &
					\includegraphics[width=0.25\linewidth]{new100_DenseNet_Test_cos.pdf}\\
						\multicolumn{2}{c}{\footnotesize{(a) ResNet34 on CIFAR100}} &  \multicolumn{2}{c}{\footnotesize{(b) DenseNet121 on CIFAR100}}\\
			\end{tabular}}
		\end{center}
	\vspace{-0.8em}
		\caption{Training and test accuracy (\%) of CNNs on CIFAR100 dataset with learning rate using the cosine annealing schedule.} \label{fig:cnn4}
	\end{figure*}

\paragraph{LSTMs on language modeling}
To further evaluate the performance of MSBPG , we conducted experiments on LSTM with the Penn Treebank dataset, and report the training and test perplexity (lower is better). Adam generally have better generalization capacity than SGD on language modeling \citep{fu2016using,siami2019performance}, and therefore is the default optimization algorithm for training LSTMs.
Here we follow the commonly used experimental setting for training LSTMs \citep{zhuang2020adabelief,chen2021closing}, which reduces the stepsize to 0.1 times its original value two times (at 75th epoch and 150th epoch) during the training process. We also conducted experiments with the cosine annealing learning rate (stepsize) schedule \citep{loshchilov2016sgdr}, which is the most frequently used learning rate schedule in practice. For training hyperparameters, we use the default settings for SGD, Adam, and AdamW in training 1-, 2-, 3-layer LSTMs \citep{zhuang2020adabelief,chen2021closing}. For MSBPG , we set its learning rate as 25, 80, 80 for 1-, 2-, 3-layer LSTMs with momentum coefficient $\beta=0.9$, weight decay coefficient $\gamma=2\times 10^{-6}$. For the layerwise kernel function $\phi_i({\mW}_i)=\frac{1}{2}\|{\mW}_i\|^2+\frac{\delta}{r}\|{\mW}_i\|^{r}$, we set $r=4$ and $\delta=1\times 10^{-6}$. From Figure \ref{fig:lstm1} and Figure \ref{fig:lstm2} we can see that MSBPG  converges well on training dataset for 1-, 2-, 3-layer LSTMs with both two training strategies. On the other hand, SGD with the cosine annealing learning rate schedule fails to get fully converged on the training dataset as shown in Figure \ref{fig:lstm2}. Moreover, MSBPG  consistently achieves the best generalization performance for all the experiments with at least 1 unit of test perplexity lower. This excellent generalization capacity again can be attributed to the Bregman proxmity model we employ.

 \begin{figure*}[th]
		%	\setcounter{figure}{3}
		%\setlength{\belowcaptionskip}{-2pt}
		\begin{center}
			\setlength{\tabcolsep}{0.0pt}  
			\scalebox{1}{\begin{tabular}{ccc}
            \includegraphics[width=0.33\linewidth]{Train_lstm_1layer_step.pdf}&
            \includegraphics[width=0.33\linewidth]{Train_lstm_2layer_step.pdf}&
            \includegraphics[width=0.33\linewidth]{Train_lstm_3layer_step.pdf}\\

            \includegraphics[width=0.33\linewidth]{Test_lstm_1layer_step.pdf}&
            \includegraphics[width=0.33\linewidth]{Test_lstm_2layer_step.pdf}&
            \includegraphics[width=0.33\linewidth]{Test_lstm_3layer_step.pdf}\\
            \footnotesize{(a) 1-layer LSTM} &  \footnotesize{(b) 2-layer LSTM}
            &  \footnotesize{(c) 3-layer LSTM}
            \\
			\end{tabular}}
		\end{center}
	\vspace{-0.8em}
		\caption{Training and test perplexity (lower is better) of LSTMs on Penn Treebank dataset with learning rate reduced to 0.1 times of the original value at the 75th epoch and 150th epoch.} \label{fig:lstm1}
\end{figure*}

 \begin{figure*}[th]
		%	\setcounter{figure}{3}
		%\setlength{\belowcaptionskip}{-2pt}
		\begin{center}
			\setlength{\tabcolsep}{0.0pt}  
			\scalebox{1}{\begin{tabular}{ccc}
            \includegraphics[width=0.33\linewidth]{Train_lstm_1layer_cos.pdf}&
            \includegraphics[width=0.33\linewidth]{Train_lstm_2layer_cos.pdf}&
            \includegraphics[width=0.33\linewidth]{Train_lstm_3layer_cos.pdf}\\

            \includegraphics[width=0.33\linewidth]{Test_lstm_1layer_cos.pdf}&
            \includegraphics[width=0.33\linewidth]{Test_lstm_2layer_cos.pdf}&
            \includegraphics[width=0.33\linewidth]{Test_lstm_3layer_cos.pdf}\\
            \footnotesize{(a) 1-layer LSTM} &  \footnotesize{(b) 2-layer LSTM}
            &  \footnotesize{(c) 3-layer LSTM}
            \\
			\end{tabular}}
		\end{center}
	\vspace{-0.8em}
		\caption{Training and test perplexity (lower is better) of LSTMs on Penn Treebank dataset with learning rate using the cosine annealing schedule.} \label{fig:lstm2}
\end{figure*}

\paragraph{Robustness to initial point scale and stepsize}
As demonstrated in Section \ref{DNN implementation}, MSBPG  can mitigate
the problem of gradient explosion. Generally, choosing large stepsize and large initial point scale will lead to gradient explosion. Here we conduct experiments with VGG16 on CIFAR10 to verify MSBPG 's robustness in training neural networks. To be specific, we compare the performance of MSBPG  and SGD with different scales of initial point and different stepsizes, since MSBPG  and SGD have the same default learning rate ($1\times 10^{-1}$). Adaptive gradient algorithms, on the other hand, have different scale of default learning rate ($1\times 10^{-3}$) and therefore we don't include them in our comparison here. For different choices of initial point scale and stepsize, we run the optimization algorithm for 50 iterations and report the best test accuracy. As we can see from Figure \ref{fig:robust}, MSBPG  is more robust than SGD to large initial points and large stepsizes. Training deep neural networks, which has millions or billions of parameters, is sensitive to the scale of the initial point and stepsize choice. It can be seen from Figure \ref{fig:robust} that SGD fails to converge with a slight increase of the initial point scale to 4.6 or increase the stepsize from 0.1 to 0.6. MSBPG , on the other hand, can converge with the initial point scale as large as 20 and stepsize as large as 5. This robustness of MSBPG  can ease the tuning of hyperparameters for training neural networks, and can also make the training process more robust to noises and errors. 
% \red{Add a concluding remark like this: for different application scenarios, SBPG can almost achieve the best performance simultaneously. While SGD is better in image classification task, Adam is better in language model.}

\begin{figure*}[th]
		%	\setcounter{figure}{3}
		%\setlength{\belowcaptionskip}{-2pt}
		\begin{center}
			\setlength{\tabcolsep}{0.0pt}  
			\scalebox{1}{\begin{tabular}{cc}
            \includegraphics[width=0.5\linewidth]{scale_result.pdf}&
            \includegraphics[width=0.5\linewidth]{stepsize_result.pdf}\\
            \footnotesize{(a) Initial point robustness}
            &  \footnotesize{(b) Stepsize robustness}
            \\
			\end{tabular}}
		\end{center}
	\vspace{-0.8em}
		\caption{Test accuracy (\%) of VGG16 on CIFAR10 dataset with different initial point scale and stepsize choice.} \label{fig:robust}
\end{figure*}

\section{Conclusion}\label{section-conclusion}
In this paper, we consider the problem of minimizing nonconvex composite objectives where the differentiable part does not satisfy Lipschitz smoothness, which is a fundamental assumption made by classical stochastic gradient methods. To overcome this limitation, we investigate a family of stochastic Bregman proximal gradient (SBPG) methods that only require smooth adaptivity of the differentiable part. From a modeling perspective, SBPG replaces the upper quadratic approximation used in SGD with the Bregman proximity measure, which captures the non-Lipschitz geometry and results in a better approximation model. We first formulate the vanilla SBPG and establish its convergence properties under the nonconvex setting without a finite sum structure. We then propose a momentum-based version of SBPG (MSBPG) that further improves its convergence properties, making it well-suited for large-scale applications. To demonstrate the effectiveness of MSBPG, we apply it to train deep neural networks with a polynomial kernel function that ensures the smooth adaptivity of the loss function. We also demonstrate the ability of MSBPG to alleviate gradient explosion during the training of deep neural networks. We conduct numerical experiments on quadratic inverse problems and training deep neural networks to validate the effectiveness of SBPG. The experimental results on sparse quadratic inverse problems show that SBPG is more robust than classical stochastic (proximal) gradient methods in terms of stepsize selection and initial point selection. Additionally, the experimental results on deep neural networks show that MSBPG outperforms state-of-the-art optimizers in terms of efficiency, robustness of stepsize selection, and yet achieves better generalization performance. In conclusion, our work demonstrates that MSBPG can constitute a valuable addition to the existing family of optimization methods for solving stochastic nonconvex optimization problems. The enhanced robustness, improved convergence results, the ability to alleviate gradient explosion, and negligible extra computational cost thus make MSBPG a promising approach for a broad range of machine learning applications and beyond. 

%\acks{}

\appendix
%\input{appendix proofs}
\section{Proofs in Preliminaries}\label{appendix-preliminary}
\noindent{\bf Proof of Lemma \ref{well-def-Breg-prox}}
First, we prove the uniqueness of the solution. Problem \eqref{Breg-prox-def} is equivalent to the following problem:
\[
\arg\min_{\vu\in\overline{C}}\;\Psi(\vu):=\alpha R(\vu)+\phi(\vu)-\inprod{\nabla\phi(\vx)}{\vu}.
\]
We have that
\[
\Psi(\vu)\geq \alpha R(\vu)
+\phi(\vu)-\norm{\nabla\phi(\vx)}\norm{\vu}
\geq\|\vu\|\Big(\frac{\alpha R(\vu)+\phi(\vu)}{\|\vu\|}-\|\nabla\phi(\vx)\|\Big).
\]
As $\|\vu\|\rightarrow\infty$, we have $\Psi(\vu)\geq\|\vu\|\Big(\frac{\alpha R(\vu)+\phi(\vu)}{\|\vu\|}-\|\nabla\phi(\vx)\|\Big)=\infty$, where we use the fact that $\phi$ is supercoercive and $R$ is convex. Since $\Psi$ is a proper lower-semicontinuous convex function, by the modern form of Weierstrass theorem \cite[Chapter 1]{rockafellar1997convex}, we know that the solution set of \eqref{Breg-prox-def} is a nonempty compact set. Also note that $\Psi$ is a strictly convex function, which implies the uniqueness of the solution. For any Legendre function $\phi$, from \cite[Chapter 26]{rockafellar1997convex}, we have $dom\,\partial\phi=int\,dom\phi$ with $\partial\phi(\vx)=\{\nabla\phi(\vx)\}$ for all $\vx\in int\,dom\,\phi$. The optimality condition implies that $\partial\phi({\rm Prox}^\phi_{\alpha R}(\vx))$ is nonempty, which automatically forces ${\rm Prox}^\phi_{\alpha P}(\vx)\in\text{int\;dom}\phi$. This completes the proof. $\Box$%Next, we prove the solution in $\text{int\;dom}\phi$. For any Legendre function $\phi$, from \cite[Chapter 26]{rockafellar1997convex}, we have $\text{dom}\partial\phi=\text{int\;dom}\phi$ with $\partial\phi(\vx)=\{\nabla\phi(\vx)\}$ for all $x\in\text{int\;dom}\phi$. The optimal condition of \eqref{Breg-prox-def} implies $\partial\phi({\rm Prox}^\phi_{\alpha R}(\vx))$ is nonempty, which automatically requires that ${\rm Prox}^\phi_{\alpha R}(\vx)\in\text{int\;dom}\phi$. This completes the proof. 


 \medskip
 \noindent{\bf Proof of Proposition \ref{prop-Bregman-mapping-subdifferential}}
  Note that $\|\nabla\phi(\vx^+)-\nabla\phi(\vx)\|\leq L_\phi\|\vx^+-\vx\|$ and $\|\nabla F(\vx^+)-\nabla F(\vx)\|\leq L_FL_\phi\|\vx^+-\vx\|$.
By the definition of $\vx^+$, we have 
\[
\nabla F(\vx^+)-\nabla F(\vx)+\frac{\nabla\phi(\vx)-\nabla\phi(\vx^+)}{\alpha}\in\nabla F(\vx^+)+\partial R(\vx^+).
\]
Thus, we obtain 
\[
{\rm dist}\left(0,\partial\Phi(\vx^+)\right)\leq\bigg\|\nabla F(\vx^+)-\nabla F(\vx)+\frac{\nabla\phi(\vx)-\nabla\phi(\vx^+)}{\alpha}\bigg\|\leq\left(L_FL_\phi+\frac{L_\phi}{\alpha}\right)\|\vx^+-\vx\|.
\]
Note that $\|\vx^+-\vx\|=\alpha\|\cG_\alpha(\vx)\|$, which completes the proof. $\Box$

\medskip
\noindent{\bf Proof of Proposition \ref{nonexpansive}}
By the definition of ${\rm Prox}_R^\phi(\cdot)$,  $x_i\in\partial R(\vx_i^+)+\nabla\phi(\vx^+_i)$, $i=1,2$. Since $\partial R(\cdot)$ is monotone, then 
 $\inprod{\vx_1-\vx_2-(\nabla\phi(\vx^+_1)-\nabla\phi(\vx^+_2))}{\vx^+_1-\vx^+_2}\geq0$. From the $\mu$-strong convexity of $\phi$, it follows that $\inprod{\vx_1-\vx_2}{\vx^+_1-\vx^+_2}\geq\inprod{\nabla\phi(\vx^+_1)-\nabla\phi(\vx^+_2)}{\vx^+_1-\vx^+_2}\geq\mu\|\vx^+_1-\vx^+_2\|^2$. Therefore, $\|\vx^+_1-\vx^+_2\|\leq\frac{1}{\mu}\|\vx_1-\vx_2\|$. $\Box$






\begin{comment}
\medskip
\noindent{\bf Proof of Lemma \ref{three-conditions-KL}}\\
Since $\E[\|\varepsilon_k\||\cF_k]\leq\sqrt{\E[\|\varepsilon_k\|^2|\cF_k]}\leq\frac{\sigma}{\sqrt{m_k}}$, then we have $\sum_{k\geq0}\E[\|\varepsilon_k\||\cF_k]<\infty$. On the other hand, $\|\varepsilon_k\|\vartriangleleft\cF_{k+1}$, by Levy's sharpening of the Borel-Cantelli lemma, see for example, \cite[Chapter 1, Theorem 21]{meyer2006martingales}, we have $\sum_{k\geq0}\|\varepsilon_k\|<\infty$ almost surely.

\noindent Claim 1: By Theorem \ref{vanilla-as-convergence}, we know that $\{x^k\}$ is bounded. Note that $\phi$ is locally Lipshitz, and domain of $\phi$ is the whole space, therefore, there exist a compact set $\cU\supset\{x^k\}$, and $\phi$ is ${L_\phi}$-Lipschitz smooth over $\cU$. By Lemma \ref{basic-descent-lemma-2}, we have 
\[
\Phi(x^{k+1})\leq\Phi(x^k)-\frac{\mu}{2}(\frac{2}{\alpha}-L_F)\|x^{k+1}-x^k\|^2+\inprod{\varepsilon_k}{x^{k+1}-x^k}.
\]
Since $\{x^k\}_{k\geq0}$ is bounded, we have that $\inprod{\varepsilon_k}{x^{k+1}-x^k}\leq C\|\varepsilon_k\|$ for some constant $C$. We prove the Claim 1 by setting $\alpha=\frac{\mu}{2}(\frac{2}{\alpha}-L_F)$, $\epsilon_k$ will be determined later.

\noindent Claim 2: By the optimality condition for \eqref{vanilla-SBPG}, we have 
\[
\left(\nabla F(x^{k+1})-\nabla F(\vx^k)\right)+\frac{1}{\alpha}\left(\nabla\phi(x^k)-\nabla\phi(x^{k+1})\right)+{\varepsilon_k}\in\partial\Phi(x^{k+1}).
\]
Therefore, we deduce that
\[
\|w^{k+1}\|\leq(L_F{L_\phi}+\frac{1}{\alpha}{L_\phi})\|x^k-x^{k+1}\|+\|\varepsilon_k\|,\;w^{k+1}\in\partial\Phi(x^{k+1}),
\]
which proves the Claim 2 by setting $\beta=(L_F{L_\phi}+\frac{{L_\phi}}{\alpha})$, $\epsilon_k=(C+1)\|\varepsilon_k\|$. $\Box$
\end{comment}


%\noindent{\bf Proof of Lemma \ref{SMAE-recursion-lemma}}\\



\begin{comment}
\section{Proofs of Lemmas in Variance Reduction SBPG}
\noindent{\bf Proof of Lemma \ref{descent-lemma-general}}\\
It follows from the relative smoothness of $F$ that
\[
\begin{aligned}
F(y)&\leq F(x)+\inprod{\nabla F(x)}{y-x}+L_F{D_\phi}(y,x)\\
F(x)&\leq F(z)+\inprod{\nabla F(x)}{x-z}+L_F{D_\phi}(z,x),
\end{aligned}
\]
Add the two inequality above, we have
\[
F(y)\leq F(z)+\inprod{\nabla F(x)}{y-z}+L_F{D_\phi}(y,x)+L_F{D_\phi}(z,x).
\]
Similar to argument of Lemma \ref{basic-descent-lemma-1}, we have
\[
R(y)\leq R(z)-\inprod{g}{y-z}+\frac{1}{\alpha}[{D_\phi}(z,x)-{D_\phi}(y,x)-{D_\phi}(z,y)].
\]
Combine inequalities above, we can obtain Lemma \ref{descent-lemma-general}. $\Box$

\medskip
\noindent{\bf Proof of Lemma \ref{svrg-Lypunov-lemma}}\\
Define $\bar x^s_{t+1}:={\rm Prox}^\phi_{\alpha R}(\nabla\phi^*(\nabla\phi(x^s_t)-\alpha\nabla F(x^s_t)))$ for the use of analysis. 
Apply Lemma \ref{descent-lemma-general} by setting $y=\bar x^s_{t+1},x=z=x^s_t, g=\nabla F(x^s_t)$, we have
\[
\Phi(\bar x^s_{t+1})\leq\Phi(x^s_t)+(L-\frac{1}{\alpha}){D_\phi}(\bar x^s_{t+1},x^s_t)-\frac{1}{\alpha}{D_\phi}(x^s_t,\bar x^s_{t+1})
\]
Apply Lemma \ref{descent-lemma-general} by setting $y=x^s_{t+1},x=x^s_t,z=\bar x^s_{t+1},g=v^{s}_t$, we have
\[
\begin{aligned}
\Phi(x^s_{t+1})\leq \Phi(\bar x^s_{t+1})&+\inprod{x^s_{t+1}-\bar x^s_{t+1}}{\nabla F(x^s_t)-v^{s}_t}\\
&+(L-\frac{1}{\alpha}){D_\phi}(x^s_{t+1},x^s_t)+(L+\frac{1}{\alpha}){D_\phi}(\bar x^s_{t+1},x^s_t)-\frac{1}{\alpha}{D_\phi}(\bar x^s_{t+1},x^s_{t+1})
\end{aligned}
\]
Combine two inequalities above, we have 
\be
\begin{aligned}
\Phi(x^s_{t+1})\leq \Phi( x^s_{t})&+\inprod{x^s_{t+1}- \bar x^s_{t+1}}{\nabla F(x^s_t)-v^{s}_t}-\frac{1}{\alpha}{D_\phi}(x^s_t,\bar x^s_{t+1})\\
&+(L-\frac{1}{\alpha}){D_\phi}(x^s_{t+1},x^s_t)+2L{D_\phi}(\bar x^s_{t+1},x^s_t)-\frac{1}{\alpha}{D_\phi}(\bar x^s_{t+1},x^s_{t+1})
\end{aligned}
\label{svrg-expt-proof-1}
\ee
Note that $x^s_{t+1}={\rm Prox}_{\alpha R}^\phi(\nabla\phi^*(\nabla\phi(x^s_t)-\alpha v^{s}_t))$, by the proof of Lemma \ref{inprod-bound}, we have
\[
\begin{aligned}
&|\inprod{x^s_{t+1}-\bar x^s_{t+1}}{\nabla F(x^s_{t})-v^{s}_t}|
\leq\frac{\alpha}{\mu}\|v^{s}_t-\nabla F(x^s_t)\|^2\\
=&\frac{\alpha}{\mu}\|\nabla f_{i^s_t}(x^s_t)-\nabla f_{i^s_t}(\tilde x^s)+\nabla F(\tilde x^s)-\nabla F(x^s_t)\|^2\\
\leq&\frac{2\alpha}{\mu}\|\nabla f_{i^s_t}(x^s_t)-\nabla f_{i^s_t}(\tilde x^s)\|^2+\frac{2\alpha}{\mu}\|\nabla F(\tilde x^s)-\nabla F(x^s_t)\|^2\\
\leq&\frac{4\kappa\alpha}{\mu}{D_\phi}(\tilde x^s,x^s_t).
\end{aligned}
\]
 Substitute the above two inequalities into \eqref{svrg-expt-proof-1}, we get 
 {
 \be
 \begin{aligned}
 \Phi(x^s_{t+1})-\Phi( x^s_{t})\leq &\frac{4\kappa\alpha}{\mu}{D_\phi}(\tilde x^s,x^s_t)-\frac{1}{\alpha}{D_\phi}(x^s_t,\bar x^s_{t+1})\\
&+(L-\frac{1}{\alpha}){D_\phi}(x^s_{t+1},x^s_t)+2L{D_\phi}(\bar x^s_{t+1},x^s_t)-\frac{1}{\alpha}{D_\phi}(\bar x^s_{t+1},x^s_{t+1}).
\end{aligned}
\label{svrg-proof-expt-2}
 \ee
 }
 Note that $\frac{2L}{c_\phi}\leq\frac{1}{2\alpha}$, thus, \eqref{svrg-proof-expt-2} becomes 
 \[
  \begin{aligned}
 \Phi(x^s_{t+1})-\Phi( x^s_{t})\leq &\frac{4\kappa\alpha}{\mu}{D_\phi}(\tilde x^s,x^s_t)-(\frac{1}{\alpha}-L){D_\phi}(x^s_{t+1},x^s_t)-\frac{1}{4\alpha}{D_\phi}(x^s_{t},\bar x^s_{t+1}).
\end{aligned}
 \]
Introduce Lyapunov function $\cE^{s}_t=\Phi(x^s_t)+r_t{D_\phi}(\tilde x^s,x^s_t)$, where $r_t$ will be determined later. Substitute ${D_\phi}(\tilde x^s,x^s_{t+1})\leq{{\hat c}_\phi}({D_\phi}(\tilde x^s,x^s_t)+{D_\phi}(x^s_{t+1},x^s_t))$ into above inequality, we get the following inequality,
{\small
\be
\cE^{s}_{t+1}\leq \cE^{s}_t+(r_{t+1}{{\hat c}_\phi}+\frac{4\kappa\alpha}{\mu}-r_t){D_\phi}(\tilde x^s,x^s_t)-(\frac{1}{\alpha}-L-r_{t+1}{{\hat c}_\phi}){D_\phi}(x^s_{t+1},x^s_t)-\frac{1}{4\alpha}{D_\phi}(x^s_{t},\bar x^s_{t+1}).
\label{svrg-proof-expt-3}
\ee}
Let $\{r_t\}_{t=0,...m}$ satisfy the condition:
\[
(i)\;r_{t+1}{{\hat c}_\phi}+\frac{4\kappa\alpha}{\mu}=r_t,\;(ii)\;\frac{1}{\alpha}-L-r_{t+1}{{\hat c}_\phi}\geq\frac{1}{2\alpha},\;(iii)\;r_m=0.
\]
By $(i),(iii)$, we have 
\[
r_t=\left\{\begin{array}{cc}
  \frac{{{\hat c}_\phi}^{m-t}-1}{{{\hat c}_\phi}-1}\frac{4\kappa\alpha}{\mu}   & if\;{{\hat c}_\phi}\neq1 \\
  \frac{4(m-t)\kappa\alpha}{\mu}   & if\;{{\hat c}_\phi}=1.
\end{array}\right.
\]
So, we only require $\frac{1}{2\alpha}-L\geq{{\hat c}_\phi}\frac{{{\hat c}_\phi}^{m-t}-1}{{{\hat c}_\phi}-1}\frac{4\kappa\alpha}{\mu}$ for ${{\hat c}_\phi}\neq1$ and $\frac{1}{2\alpha}-L\geq{{\hat c}_\phi}\frac{4(m-t)\kappa\alpha}{\mu}$ for ${{\hat c}_\phi}=1$, which holds for $\alpha\leq\min\left\{\frac{1}{4L},\sqrt{\frac{\mu(1-{{\hat c}_\phi})}{16\kappa(1-{\hat c}_\phi^m)}}\right\}$ if ${{\hat c}_\phi}\neq1$, and $\alpha\leq\min\{\frac{1}{4L},\sqrt\frac{\mu}{16m\kappa}\}$ if ${c_\phi}=1$. By the construction of $r_t$ and \eqref{svrg-proof-expt-3}, we obtain our main descent inequality:
\[
\cE^{s}_{t+1}\leq\cE^{s}_t-\frac{1}{2\alpha}{D_\phi}(x^s_{t+1},x^s_t)-\frac{1}{4\alpha}{D_\phi}(x^s_{t},\bar x^s_{t+1}).
\]
$\Box$

\medskip
\noindent{\bf Proof of Lemma \ref{svrg-three-conditions-KL}}\\
Claim 1 is directly from Lemma \ref{svrg-Lypunov-lemma} by setting $\rho_1=\frac{1}{2\alpha}$. From optimality condition of Step 9 in Algorithm \ref{SVRG-algorithm}, we 
have
\[
0\in\partial R(x^{s}_{t+1})+v^{s}_t+\frac{1}{\alpha}(\nabla\phi(x^{s}_{t+1})-\nabla\phi(x^{s}_t)),
\]
 then $w^{s}_{t+1}:=\nabla F(x^{s}_{t+1})-v^{s}_t-\frac{1}{\alpha}(\nabla\phi(x^{s}_{t+1})-\nabla\phi(x^{s}_t))$, and $w^{s}_{t+1}\in\partial\Phi(x^{s}_{t+1})$. Since $\{x^s_t\}$ is bounded, then, there exist compact set $\cU\supset\{x^s_t\}$. Since $\phi$ is locally Lipschitz smooth, $\phi$ is Lipschitz smooth over $\cU$ with Lipschitz constant ${L_\phi}$. Moreover,
 \[
 \begin{aligned}
 \|\nabla F(x^s_{t+1})-v^s_t\|=&\|\nabla F(x^s_{t+1})-\nabla F(x^s_{t})+\nabla f_{i^s_t}(\tilde x^s)-\nabla f_{i^s_t}(x^s_t)\|\\
 \leq&\|\nabla F(x^s_{t+1})-\nabla F(x^s_{t})\|+\|\nabla f_{i^s_t}(\tilde x^s)-\nabla f_{i^s_t}(x^s_t)\|\\
 \leq&L_F{L_\phi}\|x^s_{t+1}-x^s_t\|+L_F{L_\phi}\|\tilde x^s-x^s_t\|.
 \end{aligned}
 \]
 Therefore,
 \[
 \|w^s_{t+1}\|\leq(L_F{L_\phi}+\frac{{L_\phi}}{\alpha})\|x^s_{t+1}-x^s_t\|+L_F{L_\phi}\|\tilde x^s-x^s_t\|, 
 \]
which proves Claim 2.

\medskip
\noindent{\bf Proof of Theorem \ref{svrg-KL-rate}}\\
By \eqref{KL-recursion}, then, for any $\beta\in(0,\frac{1}{m})$, we have
\[
\begin{aligned}
\|x^{s}_{t+1}-x^{s}_t\|\leq&\beta\|x^{s}_t-x^{s}_{t-1}\|+\frac{\rho_2}{2\beta\rho_1\mu}\Delta^{s}_{t,t+1}+\beta\|\tilde x^{s}-x^{s}_{t-1}\|+\frac{\rho_3}{2\beta\rho_1\mu}\Delta^{s}_{t,t+1}\\
\leq&\beta(\|x^{s}_{t}-x^{s}_{t-1}\|+\|x^{s}_{t-1}-x^{s}_{t-2}\|...+\|x^{s}_1-x^{s}_0\|)+\frac{\rho_2+\rho_3}{2\beta\rho_1\mu}\Delta^{s}_{t,t+1}.
\end{aligned}
\]
Using the same notation in the proof of Theorem \ref{svrg-KL-convergence-thm}, we have 
\[
\begin{aligned}
\sum_{k=N}^{\infty}\|u^{k+1}-u^k\|\leq\beta&\sum_{k=N}^\infty\left(\|u^k-u^{k-1}\|+...+\|u^{ms(k)+1}-u^{ms(k)}\|\right)+\\
&\frac{ad(\rho_2+\rho_3)}{2\beta\rho_1\mu}\left(\Phi(u^N)-\Phi^{s(N)}_{t(N),*}+c_{t(N)}\|u^N-u^{ms(N)}\|^2\right)^{1-\theta}\\
\leq&\beta\sum_{k=N}^\infty\left(\|u^k-u^{k-1}\|+...+\|u^{ms(k)+1}-u^{ms(k)}\|\right)+\\
&\frac{ad(\rho_2+\rho_3)}{2\beta\rho_1\mu}\left(\left(\Phi(u^N)-\Phi^{s(N)}_{t(N),*}\right)^{1-\theta}+\left(c_{t(N)}\|u^N-u^{ms(N)}\|^2\right)^{1-\theta}\right).\\
\end{aligned}
\]
By the KL property with exponent $\theta$, $\left(\Phi(u^N)-\Phi^{s(N)}_{t(N),*}\right)^\theta\leq a(1-\theta){\rm dist}(0,\partial\Phi(u^N))$, then
\[
\begin{aligned}
\sum_{k=N}^{\infty}\|u^{k+1}-u^k\|
\leq&\beta\sum_{k=N}^\infty\left(\|u^k-u^{k-1}\|+...+\|u^{ms(k)+1}-u^{ms(k)}\|\right)+\\
&\frac{a(\rho_2+\rho_3)}{2\beta\rho_1\mu}\left(a(1-\theta)^{\frac{1-\theta}{\theta}}\left({\rm dist}(0,\partial\Phi(u_N))\right)^{\frac{1-\theta}{\theta}}+\left(c_{t(N)}\|u^N-u^{ms(N)}\|^2\right)^{1-\theta}\right)
.
\end{aligned}
\]
Note that ${\rm dist}(0,\partial\Phi(u^N))\leq\rho_2\|u^N-u^{N-1}\|+\rho_3\|u^{N-1}-u^{ms(N)}\|$, and $\|u^{N-1}-u^{ms(N)}\|\leq\|u^{N-1}-u^{N-2}\|+...+\|u^{ms(N)+1}-u^{ms(N)}\|$, then there exist $\hat\rho_1,\hat\rho_2$, such that
\[
\begin{aligned}
\sum_{k=N}^\infty\|u^{k+1}-u^k\|\leq&\beta\sum_{k=N}^\infty\left(\|u^k-u^{k-1}\|+...+\|u^{ms(k)+1}-u^{ms(k)}\|\right)\\
&+\hat\rho_1\left(\|u^N-u^{N-1}\|+...+\|u^{ms(N)+1}-u^{ms(N)}\|\right)^{\frac{1-\theta}{\theta}}\\
&+\hat\rho_2\left(\|u^N-u^{N-1}\|+...+\|u^{ms(N)+1}-u^{ms(N)}\|\right)^{2-2\theta}.
\end{aligned}
\]
Let $\Delta_N:=\sum_{k=N}^\infty\|u^{k+1}-u^k\|$, then $\|u^N-\bar u\|\leq\Delta_N$. By using this notation,
\[
(1-m\beta)\Delta_N\leq\beta\left((\Delta_{N-1}-\Delta_N)+...+(\Delta_{ms(N)}-\Delta_N)\right)+\hat\rho_1\left(\Delta_{ms(N)}-\Delta_N\right)^{\frac{1-\theta}{\theta}}+\hat\rho_2\left(\Delta_{ms(N)}-\Delta_N\right)^{2-2\theta}.
\]

\noindent$\blue\bullet$ If $\theta=0$, let $I:=\{k:u^{k+1}\neq u^k\}$, from K2, $\|u^{k+1}-u^k\|\geq{\rm dist}(0,\partial\Phi(u^{k+1}))>C>0$. By K1, $0\leq \cE^{s}_{t+1}\leq\cE^{s}_t-\rho_1{D_\phi}(u^k,u^{k+1})\leq\cE^{s}_t-\frac{\rho_1\mu}{2}C^2$. Therefore, $I$ must be a finite set, otherwise $\cE^s_t\rightarrow-\infty$.

\noindent$\blue\bullet$ If $\theta\in(0,\frac{1}{2}]$, then $\frac{1-\theta}{\theta}\geq1$ and $2-2\theta\geq1$. Since $\Delta_N$ is nonincreasing and $\Delta_N\rightarrow0$ as $N\rightarrow0$, then for sufficiently large $N$,
\[
(1-m\beta)\Delta_N\leq(\beta m+\hat\rho_1+\hat\rho_2)\left(\Delta_{ms(N)}-\Delta_N\right),
\]
so $\Delta_N\leq\frac{\beta m+\hat\rho_1+\hat\rho_2}{1+\hat\rho_1+\hat\rho_2}\Delta_{ms(N)}$. Let $N=m(s+1)$, then $\Delta_{m(s+1)}\leq\frac{\beta m+\hat\rho_1+\hat\rho_2}{1+\hat\rho_1+\hat\rho_2}\Delta_{ms}$. The Claim 2 is proven.

\bigskip
\noindent$\blue\bullet$ If If $\theta\in(1/2,1)$, then $0<\frac{1-\theta}{\theta}<2-2\theta<1$. Since
\[
(1-m\beta)\Delta_N\leq m\beta\left(\Delta_{ms(N)}-\Delta_N\right)+\hat\rho_1\left(\Delta_{ms(N)}-\Delta_N\right)^{\frac{1-\theta}{\theta}}+\hat\rho_2\left(\Delta_{ms(N)}-\Delta_N\right)^{2-2\theta}.
\]
and $\Delta_N\rightarrow0$ as $N\rightarrow0$, then for sufficiently large $N$, there exists a positive constant $c_2$, such that 
\[
\Delta_N^{\frac{\theta}{1-\theta}}\leq c_2\left(\Delta_{ms(N)}-\Delta_N\right).
\]
Let $N=m(s+1)$, then  
\(\Delta_{m(s+1)}^{\frac{\theta}{1-\theta}}\leq c_2(\Delta_{ms}-\Delta_{m(s+1)})
\). Let $h(r)=r^{-\frac{\theta}{1-\theta}}$, $c_3=\frac{1}{c_2}$, then $h(\Delta_{m(s+1)})(\Delta_{ms}-\Delta_{m(s+1)})\geq c_3$. Note that
\[
\int^{\Delta_{ms}}_{\Delta_{m(s+1)}}h(r)dr\geq h(\Delta_{m(s+1)})(\Delta_{ms}-\Delta_{m(s+1)})\geq c_3.
\]
Therefore,
\(
\frac{1-\theta}{2\theta-1}\left(\Delta_{m(s+1)}^{\frac{1-2\theta}{1-\theta}}-\Delta_{ms}^{\frac{1-2\theta}{1-\theta}}\right)\geq c_3,
\)
so $\Delta_{ms}^{\frac{1-2\theta}{1-\theta}}\geq\Delta_{0}^{\frac{1-2\theta}{1-\theta}}+\frac{2\theta-1}{1-\theta}c_3s$. There exists positive constant $c_4$, such that $\Delta_{ms}^{\frac{1-2\theta}{1-\theta}}\geq c_4^{\frac{1-2\theta}{1-\theta}}s$. Finally, 
\(
\Delta_{ms}\leq c_4s^{-\frac{1-\theta}{2\theta-1}}.
\)
Consequently, $\|\tilde x^s-\bar x\|\leq\Delta_{ms}\leq c_4s^{-\frac{1-\theta}{2\theta-1}}$.


\end{comment}
\begin{comment}
\section{Proofs in Section \ref{DNN implementation}}\label{appendix-DNN}
\noindent{\bf Proof of Proposition \ref{DNN-kernel-degree-prop}}\\
In order to compute the Hessian of $F$, two formulas are required, which can be verified directly.\vspace{-0.1cm}
\begin{lemma}
Let $\vu\in\R^{km}$, $\vg\in\R^m$, $A\in\R^{n\times k}$ and $\vb\in\R^m$. Consider two linear maps: $\vu\mapsto\Mat(\vu)g$ and $\vu\mapsto A(\Mat(\vu))^T\vb$, then, the Jacobian of the two maps are given by
\[
J_{\vu}\left[\Mat(\vu)\vg\right]=\vg^T\otimes\mathbb{I}_k,
\]
\[
J_{\vu}[A(\Mat(\vu))^Tb]=A\otimes \vb^T.
\]
\label{Jac-Hess-aux-lemma}
\end{lemma}\vspace{-0.5cm}
Let $\Mat(\vu)$ be denoted by $M$, and let $\text{Diag}(\cdot)$ be the operator that reshapes a vector into a diagonal matrix with diagonal elements equal to the corresponding entries of the vector. We use $\circ$ to denote the Hadamard product, which should not be confused with composition notation in this proof. Let $Jg$ denote the Jacobian of $g$, and let $J^Tg$ denote its transpose. Using Lemma \ref{Jac-Hess-aux-lemma}, we can compute the Jacobian and Hessian of $F$ as follows:

{\noindent\bf Jacobian of $F$ :}\\
\be 
\begin{aligned}
\frac{\partial F}{\partial\vu}&=\left(g(\vv)\otimes\mathbb I_k\right)\left[\sigma^\prime(Mg(\vv))\circ\sigma(Mg(\vv))\right],\\
\frac{\partial F}{\partial\vv}&=J^Tg(\vv)M^T\left[\sigma^\prime(Mg(\vv))\circ\sigma(Mg(\vv))\right].
\end{aligned}
\label{Jacobian_F}
\ee 

{\noindent\bf Hessian of $F$ :}\\
\be 
\begin{aligned}
&\frac{\partial^2F}{\partial\vu^2}=(1)+(2),\\
\text{where }(1)&=\left(g(\vv)\otimes\mathbb I_k\right)\text{Diag}\big(\sigma(Mg(\vv))\circ\sigma^{\prime\prime}(Mg(\vv))\big)(\vg^T(v)\otimes\mathbb I_k)\\
(2)&=\left(g(\vv)\otimes\mathbb I_k\right)\text{Diag}\big(\sigma^\prime(Mg(\vv))\circ\sigma^{\prime}(Mg(\vv))\big)(\vg^T(v)\otimes\mathbb I_k).
\end{aligned}
\label{Hessian_Fuu}
\ee 

\be 
\begin{aligned}
&\frac{\partial^2F}{\partial\vu\partial\vv}=(1)+(2)+(3),\\
\text{where }(1)&=\left(J^Tg(\vv)\right)\otimes\left[\sigma^\prime(Mg(\vv))\circ\sigma^{\prime}(Mg(\vv))\right]^T\\
(2)&=J^Tg(\vv)M^T\text{Diag}[\sigma(Mg(\vv))\circ\sigma^{\prime\prime}(Mg(\vv))]\left(g^T(\vv)\otimes\mathbb I_k\right)\\
(3)&=J^Tg(\vv)M^T\text{Diag}[\sigma^\prime(Mg(\vv))\circ\sigma^{\prime}(Mg(\vv))]\left(g^T(\vv)\otimes\mathbb I_k\right).
\end{aligned}
\label{Hessian_F_uv}
\ee 
\be 
\begin{aligned}
&\frac{\partial^2F}{\partial\vv^2}=(1)+(2)+(3),\\
\text{where }(1)&=D^2g(\vv)\left[M^T[\sigma^\prime(Mg(\vv))\circ\sigma(Mg(\vv))]\right]=\sum d_i\nabla^2g_i(\vv)\\
(2)&=J^Tg(\vv)M^T\text{Diag}[\sigma(Mg(\vv))\circ\sigma^{\prime\prime}(Mg(\vv))]MJg(\vv)\\
(3)&=J^Tg(\vv)M^T\text{Diag}[\sigma^\prime(Mg(\vv))\circ\sigma^{\prime}(Mg(\vv))]MJg(\vv),
\end{aligned}
\label{Hessian_F_vv}
\ee 
where $\vd=M^T[\sigma^\prime(Mg(\vv))\circ\sigma(Mg(\vv))]$. Now, we are ready to prove this proposition. For any $\vw\in\R^{km+n}$ and $\vh=[\vh^u;\vh^v]\in\R^{km+n}$, it suffices to prove that $\inprod{\nabla^2F(\vw)\vh}{\vh}=\mathcal O(\inprod{\nabla^2\phi(\vw)\vh}{\vh})$. From \eqref{Hessian_Fuu}\eqref{Hessian_F_uv}\eqref{Hessian_F_vv} and Assumption \ref{bound_assumption_bn}, we can easily get $\inprod{\nabla^2F(\vw)\vh}{\vh}=\mathcal O((1+\|\vw\|^2)\|\vh\|^2)$ if $\sigma$ is bounded.  On the other hand, $\nabla^2\phi(\vw)=I(1+\|\vw\|^{r-2})+(r-2)\|\vw\|^{r-4}\vw\vw^T$, $\inprod{\nabla^2\phi(\vw)\vh}{\vh}\geq(1+\|\vw\|^{r-2})\|\vh\|^2$. So, we only require $r-2\geq2$. This completes the proof. $\Box$
    
\end{comment}


%\input{appendix extra figures}

\section{Additional Figures}\label{app-fig}
\begin{figure}[H]
\centering
\subfigure[Illustration of SBPG and SGD updates]{
\includegraphics[width=5.5cm,height=4.5cm]{pull_back_graph.pdf}
%\caption{fig1}
}
%\quad
\subfigure[Different initial points for SBPG]{
\includegraphics[width=5.5cm]{initial_test-eps-converted-to.pdf}
}
\caption{Figure (a) depicts SGD and SBPG updates. SBPG includes a "pull back" mechanism that prevents the point from moving excessively in any given direction. $"P"$ and $"D"$ refer to the primal and dual spaces, respectively, and these terms are commonly used in the mirror descent method literature (see, e.g., \cite{bubeck2015convex,nemirovskij1983problem}). Figure (b) illustrates the effect of {choosing the initial point from a ball of radius $r$
on the SBPG when $r$ changes} for 
the QIP example with $d=100$ and $n=5000$. All initial step sizes are set to $1\times10^{-3}$. As shown in Figure (b), even for an initial point that is far from the optimal point, SBPG can pull back the iterates to the optimal point. }
\label{additional-figure}
\end{figure}



\vskip 0.2in
\bibliography{main}


\end{document}
